\section{Chiral Hochschild Cohomology and Bulk--Boundary Correspondence}
\label{sec:chiral_hochschild}
\label{sec:chiral_hochschild_expanded}
\label{sec:derived-center}
\label{chap:hochschild}% phantom: referenced from frontier chapters
\label{ch:hochschild}% alias for dangling \ref{ch:...} from frontier/preface
\label{ch:theory-hochschild}%% alias for dangling \ref sites in sc_chtop_heptagon

\noindent
What is the bulk of a $3$d holomorphic-topological theory whose
boundary is the chart algebra $\bdyAlg = A_b$? The reading discipline
of the universal stage chain
$\mathsf{P} \to \mathsf{C} \to \mathsf{S} \to \mathsf{Z} \to \mathsf{A}$
forbids the level-$\mathsf{S}$ shadow $\BarTwc{A_b}$ from naming the
bulk \textup(\cref{rem:bar-vs-centre-canonical}\textup); the bulk
sits at level $\mathsf{Z}$. The answer is the chiral Hochschild
cohomology of $A_b$ on $A_b$:
\[
\bulkOf{A_b}
 \;=\;
 \bulkChirHoch{A_b}
 \;=\;
 \mathrm{ChirHoch}^\bullet(A_b, A_b),
\]
established by dimensional descent from $2$d Poisson sigma models
and by Khan--Zeng's $3$d HT Poisson sigma model \cite{KZ25}. The
universal stage stratification of the four parallel chiral
Hochschild objects
\textup($\mathrm{ChirHoch}^\bullet$, $\mathrm{ChirHoch}_\bullet$,
$\mathrm{HC}^{\mathrm{cyc}}$, $\mathrm{HC}^-$\textup) is
\cref{thm:cyclic-hochschild-stage-stratification}: only
$\mathrm{ChirHoch}^\bullet$ is the bulk; the others are shadows
requiring further licensing data.

\begin{remark}[Hochschild control is $L_\infty$-native]
\label{rem:hochschild-linf-native}
\index{Hochschild cohomology!L-infinity native}
The cochain complexes appearing in this section are strict models. The
closed-sector deformation theory they encode is invariant only up to
filtered $L_\infty$ quasi-isomorphism, and the later brace and
Gerstenhaber operations are explicit coordinates on that
homotopy-coherent object rather than as the whole invariant story
(Convention~\ref{conv:vol2-strict-models}).
\end{remark}

\begin{definition}[Coderivation model for chiral Hochschild cochains]
\label{def:chiral-hochschild-coderivations}
Let $A$ be an $\Ainf$ chiral chart algebra with reduced chiral bar
coalgebra
\[
\barBch(A)=T^c(s^{-1}\overline A),
\qquad
B^{\mathrm{ch}}(A)=k\oplus\barBch(A),
\]
and structure coderivation
\[
D_A=\sum_{r\geq 1}D_{m_r}
\in \Coder^1\bigl(B^{\mathrm{ch}}(A)\bigr),
\qquad
D_A(1)=0,\qquad D_A^2=0.
\]
The strict model for chiral Hochschild cochains used in this chapter is
\[
C^\bullet_{\mathrm{ch}}(A,A)
:=\Coder_0\bigl(B^{\mathrm{ch}}(A)\bigr)
\simeq
\Coder\bigl(\barBch(A)\bigr),
\]
where \(\Coder_0\) denotes coaugmentation-preserving coderivations.
A $p$-ary cochain is the cogenerator projection of a coderivation,
\[
f_p\colon (s^{-1}\overline A)^{\otimes p}\longrightarrow s^{-1}\overline A,
\]
with the same chiral spectral-parameter expansions as the operations
of $A$. The associated coderivation is obtained by inserting $f_p$ in
all consecutive slots of a bar word, with the suspended Koszul sign.
\end{definition}

\begin{proposition}[Hochschild differential from the structure coderivation]
\label{prop:hochschild-differential-coderivation}
On the strict coderivation model the chiral Hochschild differential is
\[
d_{\mathrm{Hoch}} f=[D_A,f]
=D_A\circ f-(-1)^{|f|}f\circ D_A .
\]
Consequently $d_{\mathrm{Hoch}}^2=0$. The bulk object named in this
chapter is the cohomology of this dg Lie--brace model, not the
degree-zero ordinary centre and not the bar coalgebra.
\end{proposition}

\begin{proof}
The commutator of coderivations is again a coderivation, so
\(\Coder_0(B^{\mathrm{ch}}(A))\simeq\Coder(\barBch(A))\) is a dg Lie
algebra once \(D_A\) is fixed. Since \(A\) is an \(\Ainf\) chiral
algebra, \(D_A^2=0\). Hence
\[
d_{\mathrm{Hoch}}^2f=[D_A,[D_A,f]]
=\tfrac12[[D_A,D_A],f]=0
\]
by the graded Jacobi identity. The brace operations used below are the
same insertion operations on coderivation corestrictions, with
spectral parameters carried through the chiral expansions.
\end{proof}

\begin{remark}[Four-object cyclic-Hochschild discipline
\textup{(}licensing $\gamma + \epsilon$;
\ClaimStatusProvedHere\textup{)}]
\label{rem:cyclic-hochschild-four-objects}
\index{cyclic Hochschild!four-object discipline}%
\index{Beilinson cut!cyclic Hochschild four objects}%
The chiral-Hochschild side of the bulk--boundary pair carries
\emph{four} distinct objects which the manuscript never
conflates. Naming the object is licensing type $\gamma$ (ambient
declaration); the identification with the bulk
$\bulkOf{A} = \bulkChirHoch{A}$ requires an additional $\epsilon$
(orientation / non-degeneracy) hypothesis specific to each of the
last three. The discipline is structural: a statement that fails
to declare which of the four objects is meant fails the cut.
\begin{enumerate}[label=\textup{(\roman*)}]
\item \emph{Cohomology, the bulk.}
$\bulkChirHoch{A} = \mathrm{ChirHoch}^{\bullet}(A, A)$ is the
chiral Hochschild \emph{cohomology}. This is the bulk:
$\bulkOf{A} = \bulkChirHoch{A}$ as a chain-level identification
on the verified scope (\cref{thm:chd-ds-hochschild},
\cref{thm:hochschild-bridge-higher-genus}). It carries the
$E_2$-Gerstenhaber / brace structure and is the Drinfeld-centre
side of $\SCchtop$ heptagon face~(6).
\item \emph{Homology, a shadow.}
$\mathrm{ChirHoch}_{\bullet}(A, A)$ is the chiral Hochschild
\emph{homology}. Identification with the bulk is not automatic:
it requires a Calabi--Yau / dualising-line orientation. On a
$\mathrm{CY}_d$-ambient with a chosen dualising line $\omega$,
the orientation pairs
$\mathrm{ChirHoch}_{\bullet}(A, A)[d] \xrightarrow{\sim}
\bulkChirHoch{A}$, but absent that orientation the homology is a
genuinely distinct shadow. Calling it ``the bulk'' is the
collapse \emph{shadow $=$ object} (forbidden slogan no.~2 of the
voice table).
\item \emph{Cyclic, an $S^1$-equivariant shadow.}
$\mathrm{HC}^{\mathrm{cyc}}(A) = \mathrm{ChirHoch}_{\bullet}(A,
A)_{S^1}$ is the cyclic complex (Connes--Tsygan--Loday). The
identification of $\mathrm{HC}^{\mathrm{cyc}}(A)$ with a face of
the bulk-boundary pair requires the $S^1$-equivariant spectral
sequence to converge in the chosen ambient. In
$\mathrm{Ch}(\mathrm{Vect})$ the SS does not converge for class
$\mathsf M$ algebras; in the weight-completed (pro-object) /
$J$-adic ambient (Convention~\ref{conv:vol2-strict-models},
\cref{rem:B-j-hierarchy-convention-chiral}), $B^{(0)}$ is a
genuine mixed-complex differential and the SS converges. Cyclic
$=$ bulk is licensed only under the ambient declaration.
\item \emph{Negative cyclic, the K-theoretic shadow.}
$\mathrm{HC}^{-}(A)$ is the negative cyclic target of the
Chern character from $K$-theory. A Bridgeland
stability/orientation datum can orient the trace, choose a chamber,
and supply the filtration on the character; it does not make
$\mathrm{HC}^{-}(A)$ equal to $K$-theory. Identification with a
face of the $\SCchtop$ heptagon as an $\mathrm{HC}^-$-valued
CY-trace pairing is the $B^{(2)}$ statement of
\cref{rem:B-j-hierarchy-convention-chiral} (operative at genus
$g \ge 1$ via the $S^2$-framing of
$\overline{\mathcal M}_{g,1}$).
\end{enumerate}
The discipline produces three forbidden slogans, each a special
case of \emph{shadow $=$ object}:
``$\mathrm{ChirHoch}_{\bullet} = $ bulk'' (drops $\epsilon$:
needs CY orientation),
``$\mathrm{HC}^{\mathrm{cyc}} = $ bulk'' (drops $\gamma$: needs
weight-completed ambient for SS convergence),
``$\mathrm{HC}^{-} = K$-theory'' (reverses the Chern character
and drops the relative/nilpotent hypotheses under which cyclic
invariants can approximate $K$-theory). The opening definitions of this
section work with object~(i) throughout; objects~(ii)--(iv) enter
only at the explicit sites that license them. Primary literature:
Connes \cite{Con85}, Loday--Quillen \cite{LQ84} on the cyclic /
mixed-complex framework; \cite{BD04} ch.~3 for the chiral
ambient; \cite{HA} ch.~5 for the $E_d$ stratification.
\end{remark}

\begin{remark}[Bar $H^\bullet$ versus cyclic primitive $H^\bullet$:
the load-bearing fifth distinction]
\label{rem:bar-vs-cyclic-primitive}
\index{bar H2!vs cyclic primitive H2}
\index{Beilinson cut!bar vs cyclic invariants}
The four-object discipline above distinguishes
$\mathrm{ChirHoch}^\bullet$, $\mathrm{ChirHoch}_\bullet$,
$\mathrm{HC}^{\mathrm{cyc}}$, $\mathrm{HC}^{-}$ within the cyclic
package. A separate Beilinson cut distinguishes
$H^2_{\mathrm{bar}}$ (the cohomology of the bar coalgebra) from
$H^2_{\mathrm{cyc, prim}}$ (the primitive cyclic tangent space): bar
$H^2$ detects Koszulness failure, while cyclic primitive $H^2$
detects the Maurer--Cartan tangent dimension after the current orbit
$\mathbb{C} \cdot \eta$ has been quotiented out by Whitehead. At
admissible $L_k(\mathfrak{sl}_3)$ for $k = -3 + p/q$ with $q \ge 3$
\textup{(}Vol~I Theorem~\textup{V1-thm:admissible-sl3-non-koszul-qge3}\textup{)},
$\dim H^2_{\mathrm{bar}}(L_k(\mathfrak{sl}_3)) \ge 2$ is witnessed by
two Cartan classes $[H_1], [H_2]$, while
$\dim H^2_{\mathrm{cyc, prim}}(L_k(\mathfrak{sl}_3)) = 0$ because the
Cartan classes are current-containing and are absorbed by the level
direction $\mathbb{C} \cdot \eta$. Collapsing
$H^2_{\mathrm{bar}}$ with $H^2_{\mathrm{cyc, prim}}$ is the master
``shadow~$=$~object'' pattern: bar and cyclic invariants are two
different stages on $\mathsf P \to \mathsf C \to \mathsf S \to \mathsf Z \to \mathsf A$,
bar on $\mathsf S$-stage twisting and cyclic on the $\mathsf Z$-stage
MC moduli at fixed bulk. Calling either ``the cohomology of
$L_k(\mathfrak{sl}_3)$'' without the licensing tag is a defect.
\end{remark}

\begin{theorem}[Chiral Hochschild stage stratification on the universal
stage chain; \ClaimStatusProvedHere on the listed loci; licensing
tags $\beta + \gamma + \epsilon$]
\label{thm:cyclic-hochschild-stage-stratification}%
\index{universal stage chain!cyclic-Hochschild stratification|textbf}%
\index{cyclic Hochschild!stage stratification}%
The four chiral-Hochschild objects sit on the universal stage chain
$\mathsf{P} \to \mathsf{C} \to \mathsf{S} \to \mathsf{Z} \to \mathsf{A}$
as follows.
\begin{enumerate}[label=\textup{(\roman*)}]
\item \emph{The bulk lives at $\mathsf{Z}$}:
$\bulkChirHoch{A} = \mathrm{ChirHoch}^{\bullet}(A, A)$ is the
canonical level-$\mathsf{Z}$ object on $A_b$, identified with
$\Zderch(A_b)$ by the derived-endomorphism definition.  A family datum
$H_H(A_b;S)$ computes its cohomological support.
\item \emph{Homology at $\mathsf{S}$, promoted to $\mathsf{Z}$ via
$\epsilon$-CY orientation}:
$\mathrm{ChirHoch}_{\bullet}(A, A)$ sits at level $\mathsf{S}$ as a
homology shadow. The promotion to $\mathsf{Z}$ requires effectiveness
data $\effPfaffOrient$: a Calabi--Yau-$d$ structure on the chiral
algebra $A$ and a chosen $d$-dualising line $\omega$. Under
$\effPfaffOrient$, the orientation pairing
\[
\mathrm{ChirHoch}_{\bullet}(A, A)[d]
 \;\xrightarrow{\;\langle -, \omega\rangle\;}\;
 \mathrm{ChirHoch}^{\bullet}(A, A) \;=\; \bulkChirHoch{A}
\]
is an equivalence of chiral $E_2$-modules. The shift by $d$ is the
chiral CY-dimension; the pairing is the chiral upgrade of
Calabi--Yau Poincar\'e duality
\textup(Kontsevich--Soibelman \cite{KS09}, \S 10\textup).
\item \emph{Cyclic at $\mathsf{S}$, promoted to $\mathsf{Z}$ via
$\gamma$-ambient declaration}: $\mathrm{HC}^{\mathrm{cyc}}(A) =
\mathrm{ChirHoch}_{\bullet}(A, A)_{S^1}$ sits at level $\mathsf{S}$ as
a cyclic shadow. The promotion to $\mathsf{Z}^{S^1}$ requires
$\hypAmbientWtCpl$ (weight-completed pro-object ambient) for the
$S^1$-equivariant Hochschild spectral sequence to converge. In the
weight-completed ambient, the Connes operator $B^{(0)}$ is a genuine
mixed-complex differential (Connes \cite{Con85}, Loday--Quillen
\cite{LQ84}; chiral lift via Vol~I
$\hypAmbientWtCpl$ ambient declaration \texttt{conv:vol2-strict-models}),
and the cyclic comparison
\[
\mathrm{HC}^{\mathrm{cyc}}(A)
 \;\xrightarrow{\;\sim\;}\;
 \bulkChirHoch{A}^{S^1}
\]
is an equivalence in the chosen ambient. The map fails in
$\mathrm{Ch}(\mathrm{Vect})$ for class~$\mathsf{M}$ algebras
(Virasoro at $c \ne 25, 26$, principal $\cW_N$, etc.):
the SS does not converge.
\item \emph{Negative cyclic at $\mathsf{S}$, receiving the
$\mathsf{A}$-stage Chern character via $\epsilon$-Bridgeland
orientation}: $\mathrm{HC}^{-}(A)$ is the negative-cyclic trace
shadow of the $A$-module theory. Under a Bridgeland stability
orientation $\sigma \in
\mathrm{Stab}(D^b\mathrm{Coh}(A))$
(Bridgeland \cite{BM01,Bridgeland07}) and the chosen
Calabi--Yau trace datum, the Chern character has the direction
\[
 K^{\mathrm{ch}}\bigl(\mathrm{Mod}(A)\bigr)
 \;\xrightarrow{\;\operatorname{ch}^{-}_{\sigma}\;}\;
 \mathrm{HC}^{-}(A)[\sigma].
\]
The level-$\mathsf{A}$ object is
$K^{\mathrm{ch}}(\mathrm{Mod}(A))$, the Bridgeland-oriented acting
object on $A_b$.  The target $\mathrm{HC}^{-}(A)[\sigma]$ is still a
level-$\mathsf{S}$ negative-cyclic shadow with orientation data; it is
not a level-$\mathsf{A}$ object and is not the $K$-theory object
itself.  Already for $A=\C$ in algebraic $K$-theory,
$K_0(\C)=\Z$ while $\mathrm{HC}^{-}_0(\C)$ is a complex vector space
containing $\C$, and $K_1(\C)=\C^\times$ while
$\mathrm{HC}^{-}_1(\C)=0$.  Hence no absolute equivalence
between negative cyclic homology and $K$-theory is asserted. The
$\mathrm{HC}^{-}$ trace target is the $B^{(2)}$ statement of
Remark~\ref{rem:B-j-hierarchy-convention-chiral}; the acting
structure remains the module category and its $K$-theory.
\end{enumerate}
The four levels $\mathsf{P}/\mathsf{C}/\mathsf{S}/\mathsf{Z}/\mathsf{A}$
are partitioned by the four chiral-Hochschild objects with no
overlap and no missing level: every chiral-Hochschild object on a
chiral algebra in the verified scope occupies exactly one of the
four positions, and the licensing data above is what promotes a
shadow to the centre or to the acting object.
\end{theorem}

\begin{proof}
The equality in (i) is the derived-endomorphism definition of the
chiral centre.  Volume~I Theorem~H supplies the family support from
$H_H(A_b;S)$; see \cref{thm:chd-ds-hochschild}. The
$E_2$-Gerstenhaber structure on $\bulkChirHoch{A}$ is Tamarkin's
deformation theorem \cite{Tamarkin2003} adapted to chiral
generality.

For (ii), the Calabi--Yau pairing
$\mathrm{ChirHoch}_{\bullet}(A, A)[d] \to \bulkChirHoch{A}$ is the
chiral lift of Kontsevich--Soibelman's $d$-Calabi--Yau pairing
\cite{KS09} \S~10. The chiral upgrade requires a $d$-dualising line
$\omega$ on $A$, witnessed by the trace
$\TraceC: A \otimes \cdots \otimes A \to \omega$
of the bicoloured primitive package
(Definition~\ref{def:bicoloured-primitive-package}). Without
$\effPfaffOrient$, the homology and cohomology are genuinely
distinct chain complexes; identification is the orientation
collapse.

For (iii), the spectral sequence
$E^{p, q}_2 = H^q(\mathrm{ChirHoch}_{\bullet}(A, A))_{p\text{-cyc}}
 \Rightarrow \mathrm{HC}^{\mathrm{cyc}, p+q}(A)$
converges in any pro-object / weight-completed / $J$-adic
ambient where the Connes operator $B^{(0)}$ is a chain-level
differential. The hypothesis $\hypAmbientWtCpl$ is precisely
the ambient declaration that makes $B^{(0)}$ a chain-level
differential; in $\mathrm{Ch}(\mathrm{Vect})$ for class~$\mathsf{M}$
algebras the operator is only homotopy-coherent, and the SS
does not converge to a single chain-level object. The
$\mathrm{Ch}(\mathrm{Vect})$ failure is the Vol~I result
\cref{V1-prop:bv-bar-class-m-weight-completed}.

For (iv), Connes' negative cyclic complex and the
Connes--Tsygan--Loday mixed-complex formalism supply the target of
the Chern character from $K$-theory
\cite{Con85,LQ84}. Bridgeland's stability conditions
\cite{Bridgeland07,BM01} supply orientation, chamber, and filtration
data on the Calabi--Yau module category; they do not identify
negative cyclic homology with $K$-theory. The level-$\mathsf{A}$
acting structure is the Bridgeland-oriented module category and its
$K$-theory object. The level-$\mathsf{S}$ object
$\mathrm{HC}^{-}(A)[\sigma]$ is the trace target receiving
\(\operatorname{ch}^{-}_{\sigma}\). The base-field calculation
\(K_0(\C)=\Z\), \(K_1(\C)=\C^\times\),
\(\mathrm{HC}^{-}_0(\C)\) a complex vector space, and
\(\mathrm{HC}^{-}_1(\C)=0\) separates the two objects before any
chiral enhancement is installed.

Disjointness of the four positions: an object cannot occupy two
levels simultaneously because the universal stage chain is a
chain (no commutative diagram of stage promotions). Each object's
position is fixed by its construction
(Definition~\ref{def:bicoloured-primitive-package} for the chain;
Vol~I \cite{LQ84,Con85} for the chiral-Hochschild objects).
\end{proof}

\begin{remark}[Construction Problem 4 specialises to the
$\mathsf{S} \to \mathsf{Z}$ promotion at level (iii)]
\label{rem:cp4-via-stage-stratification}
The all-chiral-generality extension of the chiral Positselski
equivalence demands the $\mathsf{S} \to
\mathsf{Z}$ promotion at level (iii) for chiral algebras outside
the Vol~I positive-energy class. Specifically, the four obstruction
loci named in
Vol~II~\cref{rem:chiral-positselski-residual} (critical-level
Kac--Moody, log-CFTs, non-positive-energy modules,
non-rational curves) are precisely the sites where the
$\hypAmbientWtCpl$ hypothesis fails to be satisfied by the
standard Vol~I weight-completion construction; the chiral Positselski
extension asks for a refined ambient that recovers the
$\mathsf{S} \to \mathsf{Z}$ promotion off the Vol~I positive-energy
class.
\end{remark}

\begin{theorem}[Bounded Virasoro cohomology and chart transport;
\ClaimStatusProvedElsewhere; licensing tag $\beta$]
\label{thm:chirhoch-virasoro-concentration}%
\index{Theorem H!Vir_c witness@$\mathrm{Vir}_c$ witness}%
\index{Virasoro algebra!chiral Hochschild concentration}%
For every central charge~$c$, Bakalov--De Sole--Kac
\cite{BDSK21} \textup{(}Theorem~7.2\textup{)} compute the
bounded vertex cohomology of $\mathrm{Vir}_c$ as
\[
 \dim H^n_{\mathrm{ch},b}(\mathrm{Vir}_c,\mathrm{Vir}_c)
 =
 \begin{cases}
 1,&n\in\{0,2,3\},\\
 0,&n\in\mathbb Z_{\geq0}\setminus\{0,2,3\}.
 \end{cases}
\]
Let $\cA_{\mathrm{Vir}_c}$ be the associated chiral algebra on a
smooth curve~$X$, and let
\[
 \chi_c^{\mathrm{bd}}\colon
 C^\bullet_{\mathrm{ch},b}(\mathrm{Vir}_c,\mathrm{Vir}_c)
 \longrightarrow
 C^\bullet_{\mathrm{ch}}(\cA_{\mathrm{Vir}_c},
                         \cA_{\mathrm{Vir}_c})
\]
be the bounded-to-chart comparison.  When $\chi_c^{\mathrm{bd}}$ is a
quasi-isomorphism, the chart cohomology has the same dimensions and
support~$\{0,2,3\}$.
\end{theorem}

\begin{proof}
The cited theorem gives the bounded calculation. A quasi-isomorphism
induces an isomorphism on every
cohomology group, which transports the displayed dimensions to the
curve chart.  This is the Virasoro benchmark recorded in Volume~I,
Theorem~\ref*{V1-thm:main-koszul-hoch}.
\end{proof}

\begin{proposition}[Rank-one even superboson benchmark;
\ClaimStatusProvedElsewhere]
\label{prop:vol2-superboson-bounded-benchmark}
Let $B_{\mathfrak h}$ be the superboson attached to a one-dimensional
even space~$\mathfrak h$. Bakalov--De Sole--Kac
\cite{BDSK21} \textup{(}Theorem~7.4\textup{)} compute
\[
 H^n_{\mathrm{ch},b}(B_{\mathfrak h},B_{\mathfrak h})
 \cong
 \bigl(S^n(\Pi\mathfrak h)\bigr)^*
 \oplus
 \bigl(S^{n+1}(\Pi\mathfrak h)\bigr)^*,
\]
and hence the dimension vector $(2,1,0,0,\ldots)$.  If the
bounded-to-chart map
\[
 \chi^{\mathrm{bd}}_{B_{\mathfrak h}}\colon
 C^\bullet_{\mathrm{ch},b}(B_{\mathfrak h},B_{\mathfrak h})
 \longrightarrow
 C^\bullet_{\mathrm{ch}}(\cA_{B_{\mathfrak h}},
                         \cA_{B_{\mathfrak h}})
\]
is a quasi-isomorphism, then the chart Hilbert polynomial is $2+t$
and its support is~$\{0,1\}$.
\end{proposition}

\begin{proof}
The cited theorem gives the bounded formula. Since
$\Pi\mathfrak h$ is one-dimensional and odd, its supersymmetric
powers occur in degrees zero and one.  The comparison
quasi-isomorphism transports these groups to the chart complex.
\end{proof}

\begin{remark}[Central sector and exceptional loci]
\label{rem:vir-concentration-exceptions}
The bounded calculation is uniform in~$c$.  Its chart interpretation
at each central charge is governed by the comparison
$\chi_c^{\mathrm{bd}}$.  Specialized Felder and free-field resolutions
provide natural models for this map at the following loci:
\begin{enumerate}[label=\textup{(\alph*)}]
\item Minimal-model points $c = c_{p,p'}$: the Felder resolution is
the natural source complex.
\item $c = 1$: the free-boson resolution gives the comparison with the
rank-one even superboson complex.
\item $c = 25,26$: the ghost-coupled resolution gives the corresponding
completed source complex.
\end{enumerate}
At each locus, a proof that $\chi_c^{\mathrm{bd}}$ is a
quasi-isomorphism transports the bounded profile to the chart complex.
\end{remark}

\subsection{Dimensional descent: 2d Poisson sigma model to 3d HT}
\subsubsection{2d Poisson sigma model}
A 2d topological theory on a worldsheet $\Sigma$ with target a Poisson manifold $(M, \pi)$, $\pi \in \Gamma(\wedge^2 TM)$, has fields:
\begin{align}
X &: \Sigma \to M \quad \text{(maps)}, \\
\eta &\in \Omega^1(\Sigma, X^* T^*M) \quad \text{(1-forms on $\Sigma$ valued in the cotangent bundle)}.
\end{align}

The action is
\begin{equation}
\label{eq:2d_PSM_action}
S_{\text{2d-PSM}}[X, \eta] = \int_\Sigma \left( \langle \eta, dX \rangle + \frac{1}{2} \langle \pi(X), \eta \wedge \eta \rangle \right).
\end{equation}
\subsubsection{Loop space and Hochschild cohomology}
Cattaneo--Felder and Kontsevich~\cite{Kon03} rewrite the space of fields as
\begin{equation}
\text{Fields}_{2d} = \text{Maps}(\Sigma, T^*[1]M),
\end{equation}
$T^*[1]M$ the $[1]$-shifted cotangent bundle (fiber coordinates $\eta_i$ of degree $+1$).

At $\Sigma = S^1$,
\begin{equation}
\text{Fields}|_{S^1} = \text{Maps}(S^1, T^*[1]M) = LT^*[1]M,
\end{equation}
the loop space of $T^*[1]M$.
\begin{theorem}[2d PSM and Hochschild Cohomology;
\ClaimStatusProvedElsewhere]
\label{thm:2d_PSM_Hochschild}
Classical observables in the 2d Poisson sigma model on $S^1$ are identified with Hochschild cochains:
\begin{equation}
\mathcal{O}_{\text{2d-PSM}}(S^1) \cong C^\bullet_{\text{Hoch}}(C^\infty(M), C^\infty(M)),
\end{equation}
where the Hochschild differential $d_{\text{Hoch}}$ matches the BV-BRST differential $Q$ in the sigma model.
\end{theorem}

This is the geometric origin of Hochschild cohomology for Poisson manifolds.
\subsubsection{3d HT analogue: chiral Hochschild cohomology}
For a 3d HT theory on $\R_t \times \C_z$ the substitutions are
\begin{itemize}
\item Worldsheet $\Sigma$ (2d) $\leadsto$ $\C_z$ (complex curve) + topological direction $\R_t$;
\item Target Poisson manifold $(M, \pi)$ $\leadsto$ Poisson vertex algebra $(V, \{\cdot_\lambda \cdot\})$;
\item Loop space $LM$ $\leadsto$ $\Ran(\C)$, the Ran space of finite subsets of $\C$.
\end{itemize}
\begin{definition}[3d HT Poisson sigma model]
\label{def:3d_HT_PSM}
The \emph{3d HT Poisson sigma model} (Khan--Zeng~\cite{KZ25}) associated to a freely generated Poisson vertex algebra $V = \text{Sym}(\C[\phi^i, \partial \phi^i, \ldots])$ has fields
\begin{align}
\Phi^i &: \R \times \C \to \C \quad \text{(scalar superfields)}, \\
\eta_i &\in \Omega^{0,1}(\C) \otimes \Omega^1(\R) \quad \text{(one-forms in $t$ and $(0,1)$-forms in $z$)}.
\end{align}
The action is:
\begin{equation}
\label{eq:3d_PSM_action}
S_{\text{3d-PSM}} = \int_{\R \times \C} \left( \eta_i (d_t + \dbar) \Phi^i + \frac{1}{2} \Pi^{ij}(\Phi, \partial \Phi, \ldots) \eta_i \wedge \eta_j \right),
\end{equation}
where $\Pi^{ij}(\lambda)$ is the PVA $\lambda$-bracket polynomial:
\begin{equation}
\{\phi^i_\lambda \phi^j\} = \Pi^{ij}(\lambda)
= \sum_{n \geq 0}\Pi^{ij}_n\frac{\lambda^n}{n!}.
\end{equation}
The reciprocal powers belong instead to the singular OPE kernel
\(\sum_{n\ge0}\Pi^{ij}_n(z-w)^{-n-1}\); the Fourier/Borel convention
sends this kernel to the polynomial
\(\sum_{n\ge0}\Pi^{ij}_n\lambda^n/n!\).
\end{definition}

\begin{theorem}[Gauge Invariance $\Leftrightarrow$ $\lambda$-Jacobi;
\ClaimStatusProvedElsewhere]
\label{thm:KZ_gauge_invariance}
(Khan--Zeng \cite{KZ25}, Theorem 2.1) The 3d HT Poisson sigma model \eqref{eq:3d_PSM_action} is gauge-invariant under BRST transformations if and only if the $\lambda$-bracket $\{\cdot_\lambda \cdot\}$ satisfies the $\lambda$-Jacobi identity:
\begin{equation}
\{\phi^i_\lambda \{\phi^j_\mu \phi^k\}\} - \{\phi^j_\mu \{\phi^i_\lambda \phi^k\}\} = \{\{\phi^i_\lambda \phi^j\}_{\lambda + \mu} \phi^k\}.
\end{equation}
\end{theorem}

The PVA axioms (Section~\ref{sec:PVA_descent}) are gauge invariance of the 3d HT sigma model.

\subsection{Chiral Hochschild cochains}

\begin{definition}[Chiral Hochschild cochains in spectral coordinates]
\label{def:chiral_hochschild}
Let \(A\) be a vertex algebra, PVA, or \(\Ainf\) chiral chart algebra.
The complex \(C^\bullet_{\mathrm{ch}}(A,A)\) is the coderivation
complex of Definition~\ref{def:chiral-hochschild-coderivations}.  Its
arity-\(p\) coordinate is the cogenerator projection
\[
f_p\colon (s^{-1}\overline A)^{\otimes p}\longrightarrow
s^{-1}\overline A((\lambda_1))\cdots((\lambda_{p-1})),
\]
subject to the chiral translation-equivariance and sesquilinearity
relations of Definition~\ref{def:ainfty_chiral}. Conversely,
each such compatible family of corestrictions extends uniquely to a
coaugmentation-preserving coderivation of \(B^{\mathrm{ch}}(A)\).
Thus the spectral-parameter formula is a coordinate description of
\(\Coder_0(B^{\mathrm{ch}}(A))\), not a complex of vertex-algebra
homomorphisms.
\end{definition}

\subsubsection{Differential on chiral Hochschild cochains}

Write \(D_A=D_{m_1}+D_{m_{\geq 2}}\), where \(m_1=Q\) is the internal
BRST/chiral differential and \(m_{\geq2}\) records the OPE collision
operations.  The Hochschild differential on
\(C^\bullet_{\mathrm{ch}}(A,A)=\Coder_0(B^{\mathrm{ch}}(A))\) is
\[
d_{\mathrm{Hoch}}=[D_A,-].
\]
In arity coordinates this commutator has an internal term
\([D_{m_1},f]\) and collision terms in which the operations \(m_r\)
are inserted into \(f\) or \(f\) is inserted into an \(m_r\).  For a
strict PVA only \(m_1\) and \(m_2\) remain, and this recovers the usual
BRST plus binary Hochschild formula with the \(\lambda\)-spectral
substitutions.  For an \(\Ainf\)-chiral algebra the higher \(m_r\) terms
are part of the differential; omitting them destroys \(d^2=0\).

\begin{lemma}[$d_{\mathrm{Hoch}}^2 = 0$; \ClaimStatusProvedHere]
\label{lem:dch_squared}
The chiral Hochschild differential satisfies \(d_{\mathrm{Hoch}}^2=0\).
\end{lemma}
\begin{proof}
By Definition~\ref{def:chiral-hochschild-coderivations},
\(D_A^2=0\). Therefore, for any coderivation \(f\),
\[
d_{\mathrm{Hoch}}^2f=[D_A,[D_A,f]]
=\tfrac12[[D_A,D_A],f]=0
\]
by the graded Jacobi identity in the Lie algebra of coderivations.
\end{proof}

\subsection{\texorpdfstring{Tamarkin deformation theory and $*$-Lie algebras}{Tamarkin deformation theory and star-Lie algebras}}

The geometric origin of chiral Hochschild cohomology is Tamarkin's theory of chiral-algebra deformations~\cite{Tam02}.

\subsubsection{\texorpdfstring{$*$-Lie bracket and chiral deformations}{Star-Lie bracket and chiral deformations}}
\begin{definition}[$*$-Lie algebra structure]
\label{def:star_lie}
On $C^\bullet_{\text{ch}}(A, A)$, define a $*$-Lie bracket
\begin{equation}
[f *_\lambda g](a_1, \ldots, a_n) = \sum_{\sigma \in \mathrm{Sh}(p,q)} \varepsilon(\sigma) \{f(a_{\sigma(1)}, \ldots)_\lambda g(a_{\sigma(p+1)}, \ldots, a_{\sigma(n)})\} - (-1)^{(|f|-1)(|g|-1)} \{g \leftrightarrow f\},
\end{equation}
where the sum is over all $(p,q)$-shuffles distributing the $n = p+q$ inputs between $f$ and~$g$, following the standard Gerstenhaber bracket convention, and $\varepsilon(\sigma)$ is the Koszul sign of the permutation.
\end{definition}

This bracket satisfies
\begin{enumerate}[label=(\roman*)]
\item $[f *_\lambda g] = -(-1)^{(|f|-1)(|g|-1)} [g *_{-\lambda-\partial} f]$ (graded antisymmetry);
\item $[f *_\lambda [g *_\mu h]] - [g *_\mu [f *_\lambda h]] = [[f *_\lambda g] *_{\lambda + \mu} h]$ ($*$-Jacobi, from associativity of the brace operations of Proposition~\ref{prop:geometric-braces-well-defined}).
\end{enumerate}

\begin{theorem}[Tamarkin deformation-quantization;
\ClaimStatusProvedElsewhere]
\label{thm:Tamarkin_deformation}
(Tamarkin \cite{Tam02}.) Deformations of a chiral algebra $A$ are controlled by the Maurer--Cartan equation in the $*$-Lie algebra $(C^\bullet_{\text{ch}}(A, A), [*_\lambda \cdot])$:
\begin{equation}
d_{\text{ch}} \alpha + \frac{1}{2} [\alpha *_\lambda \alpha] = 0.
\end{equation}

Solutions $\alpha \in C^1_{\text{ch}}(A, A)$ parametrize infinitesimal deformations of the vertex-algebra structure on $A$.
\end{theorem}

\begin{remark}[Strict and homotopy levels]
\label{rem:tamarkin-two-level}
Tamarkin's $*$-Lie algebra is the strict model $\Convstr$;
the homotopy-invariant object is the $(d{+}1)$-algebra from
Tamarkin's higher structure theorem~\cite{Tamarkin00}.
\end{remark}

\subsection{\texorpdfstring{The $E_{d+1}$ structure and the Swiss-cheese operad}{The E(d+1) structure and the Swiss-cheese operad}}
\label{subsec:tamarkin-swiss-cheese}

\begin{theorem}[Tamarkin higher structure;
\ClaimStatusProvedElsewhere]
\label{thm:tamarkin-higher-structure}
\index{Tamarkin!E d+1 structure@$E_{d+1}$ structure|textbf}
\index{Hochschild cochains!E d+1 structure@$E_{d+1}$ structure}
Let $A$ be an $E_d$-algebra. The Hochschild cochain complex
$C^*(A,A)$ carries a natural $E_{d+1}$-algebra structure
\textup{(}Tamarkin~\textup{\cite{Tamarkin00}},
Kontsevich~\textup{\cite{Kon03})}. The $E_1$-part is the
cup product; the $E_2$-part is the Gerstenhaber bracket
$[-,-]_G$; the higher $E_k$-parts ($k\ge 3$) are the
Deligne conjecture operations.
\end{theorem}

In the holomorphic-topological setting the boundary algebra
$\cA$ is an $E_1$-chiral algebra with one topological direction
($\mathbb R_t$, the $E_1$ color) and one holomorphic direction
($\mathbb C_z$, the chiral color). Tamarkin's theorem applied
to the $E_1$-topological structure underlying $\cA$ furnishes
the chiral Hochschild cochain complex
$C^*_{\mathrm{ch}}(\cA,\cA)$ with an $E_2$-structure
(the closed sector of $\mathrm{SC}^{\mathrm{ch,top}}$;
Volume~I, Theorem~\ref*{V1-thm:operadic-center-hochschild}).
The notation $C^*(\cA,\cA)$ here denotes the chiral complex,
not the classical associative Hochschild of Theorem~\ref{thm:tamarkin-higher-structure}.

\begin{proposition}[Reduced Swiss-cheese shadow from Tamarkin;
\ClaimStatusProvedHere]
\label{prop:swiss-cheese-from-tamarkin}
\index{Swiss-cheese operad!from Tamarkin}
\index{Tamarkin!Swiss-cheese identification}
The $E_2$-structure on $C^*(\cA,\cA)$ is the closed-colour structure of
the \emph{reduced} Swiss-cheese shadow of
$\mathsf{SC}^{\mathrm{ch,top}}$
\textup{(}Definition~\textup{\ref{def:SC-operations})}. It is not an
identification of the unreduced two-coloured
$\mathsf{SC}^{\mathrm{ch,top}}$ operad with $\mathrm{Disk}_3$. The
identification is:
\begin{center}
\renewcommand{\arraystretch}{1.15}
\begin{tabular}{ll}
\textbf{Tamarkin / Hochschild}
 & \textbf{reduced Swiss-cheese / HT realisation} \\
\hline
$C^*(\cA,\cA)$ with $E_2$
 & closed-colour algebra; 3d bulk only after an HT realisation \\
cup product ($E_1$-part)
 & OPE of bulk observables on $\mathbb C_z$ \\
Gerstenhaber bracket ($E_2$-part)
 & Poisson bracket from $\mathbb R_t$-direction \\
$A=\cA$ (the $E_1$-algebra)
 & boundary algebra on $\{0\}\times\mathbb C$ \\
$C^*(\cA,\cA)$-module structure on $\cA$
 & bulk-to-boundary map $\beta$ \\
MC element $\Theta\in C^2(\cA,\cA)$
 & central curvature $\kappa\cdot\omega_g$
\end{tabular}
\end{center}
\end{proposition}

\begin{proof}
The reduced closed/open shadow of
$\mathsf{SC}^{\mathrm{ch,top}}$ has two colours: closed
\textup{(}$E_2$\textup{)} and open \textup{(}$E_1$\textup{)}. An algebra over
this shadow is a pair $(B,A)$ where $B$ is an $E_2$-algebra, $A$ is an
$E_1$-algebra, and $A$ carries a compatible $B$-module structure. The
full $\mathsf{SC}^{\mathrm{ch,top}}$ datum retains the directional
holomorphic/topological restriction and is the object used for the
actual HT bulk--boundary theory.

By Theorem~\ref{thm:tamarkin-higher-structure}, $B=C^*(\cA,\cA)$
carries an $E_2$-structure when $\cA$ is $E_1$. The pair
$(C^*(\cA,\cA),\,\cA)$ with the tautological $C^*(\cA,\cA)$-module
structure on~$\cA$ (evaluation) is an algebra over the Swiss-cheese
operad: the $E_2$ bulk acts on the $E_1$ boundary by
Hochschild cochain evaluation, which is the OPE of bulk
operators near the boundary.

The closed-color operations of $\mathsf{SC}^{\mathrm{ch,top}}$
at degree~$(k,0)$ (purely closed) are
$C_*(\mathrm{FM}_k(\mathbb C))$, which are the chains on the
Fulton--MacPherson compactification. Kontsevich formality
$C_*(\mathrm{FM}_k(\mathbb C))\xrightarrow{\sim}E_2(k)$
(Theorem~\ref{thm:homotopy-Koszul}) identifies these with
the $E_2$-operations. The mixed operations at degree~$(k,m)$
are $C_*(\mathrm{FM}_k(\mathbb C))\times E_1(m)$, encoding
$k$~bulk insertions near the boundary and $m$~boundary
insertions along the topological direction; this is the
Hochschild $C^*(\cA,\cA)$-action on~$\cA$ at degree~$k$
composed with the $E_1$-structure at degree~$m$.
\end{proof}

\begin{remark}[The two colours are inevitable in the reduced shadow]
\label{rem:two-colors-inevitable}
\index{Swiss-cheese operad!inevitability of two colors}
The $E_{d+1}$-structure on Hochschild cochains is not an
additional datum imposed on the boundary algebra; it is
forced by the algebra structure itself. An $E_1$-algebra
$\cA$ automatically produces an $E_2$-algebra
$C^*(\cA,\cA)$: the deformation complex of $\cA$ is one
categorical dimension higher than~$\cA$. The Swiss-cheese
structure $(C^*(\cA,\cA),\,\cA)$ is therefore native to any
$E_1$-chiral algebra in the reduced Swiss-cheese shadow. A 3d HT
bulk/boundary system is a geometric realisation only when a
holomorphic-topological prefactorisation algebra supplies the missing
directional data.

For modular extensions: the genus-$g$ chiral Hochschild cochain
complex $C^*_g(\cA,\cA)$ on $\overline{\cM}_{g,n}$ carries
a modular $E_2$-structure, and the modular MC element
$\Theta_\cA\in C^2_g(\cA,\cA)$ is the curvature of this
structure. The holographic modular Koszul datum
$\mathcal H(T)$ is the complete packaging of the modular
$E_2$-structure and its $E_1$-module~$\cA$.
\end{remark}

\begin{computation}[Tamarkin $E_2$ structure for the
Heisenberg algebra; \ClaimStatusProvedHere]
\label{comp:tamarkin-e2-heisenberg}
\index{Heisenberg algebra!Tamarkin E2 structure@Tamarkin $E_2$ structure}
For $\cA=\cH_k$ at a fixed non-zero level, the reduced chiral
Hochschild cohomology has
\[
H^0_{\mathrm{red}}(\cH_k)=\Bbbk\mathbf 1,\qquad
H^2_{\mathrm{red}}(\cH_k)=
\Bbbk\langle[\nu_k]\rangle,\qquad
H^n_{\mathrm{red}}(\cH_k)=0\quad(n\ne 0,2).
\]
The class $[\nu_k]$ is the fixed-fibre dual-vacuum/curvature
endpoint.  The OPE-coefficient variation
$k\to k+\varepsilon$ is represented by a product cocycle before
quotienting, but for $k\neq0$ it is exact under rescaling of the
Heisenberg generator. The cup square of $[\nu_k]$ is zero in
fixed-fibre cohomology, and the Gerstenhaber bracket on this line is
zero. The completed formal base $\Bbbk[[\kappa]]$ of the chosen
one-parameter presentation $\cH_{k+\kappa}$ records the normalization
coordinate of the family, not the cohomology algebra or the moduli of
the fixed fibre.

The Swiss-cheese pair $(V_T,\cH_k)=(E_2,E_1)$ is the
simplest possible: trivial Poisson bracket, scalar
bulk-to-boundary coupling. The modular MC element
$\Theta_{\cH_k}=k\cdot\eta\otimes\Lambda$ is the point of this
$1$-dimensional normalized presentation, while $[\nu_k]$ is the
fixed-fibre Hochschild endpoint.
\end{computation}

\begin{computation}[Tamarkin $E_2$ structure for
$\widehat{\mathfrak{sl}}_2$; \ClaimStatusProvedHere]
\label{comp:tamarkin-e2-sl2}
\index{Kac--Moody!Tamarkin E2 structure@Tamarkin $E_2$ structure}
For $\cA=\widehat{\mathfrak{sl}}_2{}_k$:
$C^*(\widehat{\mathfrak{sl}}_2,\widehat{\mathfrak{sl}}_2)
\cong U(\mathfrak{sl}_2)^{\mathfrak{sl}_2}[\![k^{-1}]\!]
\cong\Bbbk[[\Omega]]$
where $\Omega=e\otimes f+f\otimes e+\tfrac12 h\otimes h$
is the Casimir. The Gerstenhaber bracket is
$[\Omega,\Omega]_G=0$ (the Casimir is central, so the
$E_2$-structure is abelian at this level).

The Swiss-cheese pair
$(V_T,\widehat{\mathfrak{sl}}_2{}_k)$ has
bulk-to-boundary map
$\beta(\Omega)=\sum_a J^a_{(-1)}J^a$
(the Sugawara construction): the Casimir element of
$U(\mathfrak{sl}_2)$ maps to the energy-momentum tensor
on the boundary via the Sugawara formula. The modular MC
element $\Theta_k$ has its weight-$1$ component
$\Theta^{[1]}|_{(0,4)}$ determined by the KZ connection
(via the KZ graph-sum computation of Volume~I).
\end{computation}

\subsection{Bulk--boundary identification via factorization homology}

The link between bulk local operators and chiral Hochschild cochains.

\subsubsection{Factorization algebras on the Ran space}

Bulk local operators in a 3d HT theory form a factorization algebra on $\C_z$ (Costello--Gwilliam~\cite{CG17}):

\begin{enumerate}[label=(\roman*)]
\item for each open $U \subseteq \C$, a cochain complex $\mathcal{O}_{\text{bulk}}(U)$;
\item for disjoint $U_1, \ldots, U_n \subseteq \C$, multiplication maps
\begin{equation}
\mathcal{O}_{\text{bulk}}(U_1) \otimes \cdots \otimes \mathcal{O}_{\text{bulk}}(U_n) \longrightarrow \mathcal{O}_{\text{bulk}}(U_1 \cup \cdots \cup U_n),
\end{equation}
satisfying associativity and locality.
\end{enumerate}

\begin{definition}[Ran space]
\label{def:Ran_space}
The \emph{Ran space} $\Ran(\C)$ is the colimit
\begin{equation}
\Ran(\C) = \operatorname{colim}_{S \in \mathrm{Fin}^{\mathrm{surj}}} \C^S,
\end{equation}
taken over the category $\mathrm{Fin}^{\mathrm{surj}}$ of non-empty finite sets and surjections, where a surjection $\alpha\colon S\twoheadrightarrow S'$ acts by the diagonal map $\Delta_\alpha\colon \C^{S'}\hookrightarrow \C^{S}$. In particular, $\Ran(\C)$ is the space of finite non-empty subsets of $\C$ (without multiplicity), topologized so that points can collide but not separate.
\end{definition}

The factorization algebra $\mathcal{O}_{\text{bulk}}$ extends to a cosheaf on $\Ran(\C)$: over the open stratum $\Ran_{\leq n}(\C) \setminus \Ran_{\leq n-1}(\C)$ of cardinality-$n$ subsets it is the $S_n$-equivariant pullback of $\mathcal{O}_{\text{bulk}}^{\boxtimes n}$ on $\C^n \setminus \Delta$; the factorization multiplication maps encode the collision behaviour across strata.

\subsubsection{Factorization homology and bulk observables}

\begin{theorem}[Bulk observables via factorization homology;
\ClaimStatusProvedElsewhere]
\label{thm:bulk_factorization}
(Ayala--Francis \cite{AF15}.) For a factorization algebra $\mathcal{A}$ on $\C$, global observables are computed by factorization homology:
\begin{equation}
\mathcal{O}_{\text{bulk}}(\C) \simeq \int_\C \mathcal{A}.
\end{equation}

For 3d HT theories this identifies bulk local operators with cochains on the Ran space.
\end{theorem}

\subsubsection{Connection to chiral Hochschild cochains}

\begin{theorem}[Bulk--Hochschild comparison for a chosen physical
prefactorization model; \ClaimStatusConditional; licensing tags
\(\beta+\gamma+\delta\)]
\label{thm:bulk_hochschild}
The algebraic closed-sector object attached to a boundary chiral
algebra \(A_\partial\) is the chiral derived centre
\[
Z^{\mathrm{der}}_{\mathrm{ch}}(A_{\partial})
\simeq C^\bullet_{\mathrm{ch}}(A_{\partial},A_{\partial})
\simeq \Coder_0\bigl(B^{\mathrm{ch}}(A_\partial)\bigr).
\]
Let $\mathsf{Obs}$ be a holomorphic--topological prefactorization
algebra on $\C\times\R$ arising from a physical realization in the
scope of Theorem~\ref{thm:physics-bridge}, and let $A_{\partial}$ be
the boundary chiral algebra obtained by restricting $\mathsf{Obs}$ to
$\C\times\{0\}$. Then factorization homology of the chosen physical
model gives a quasi-isomorphism
\begin{equation}
\Psi_{\mathrm{Ran}}\colon
\mathcal{O}_{\textup{bulk}}
\xrightarrow{\simeq}
C^\bullet_{\textup{ch}}(A_{\partial}, A_{\partial}).
\end{equation}
Equivalently, in the homotopy category, the inverse comparison is the
bulk--Hochschild map
\[
\chi_{\mathrm{HT},A_\partial}\colon
Z^{\mathrm{der}}_{\mathrm{ch}}(A_\partial)
\xrightarrow{\simeq}
\mathcal{O}_{\textup{bulk}}.
\]
The chain-level comparison proved below uses the prefactorization
model $\mathsf{Obs}$ and its factorization homology; it is not a
definition of a physical bulk from \(A_\partial\) alone.
The global bulk--boundary--line triangle
\textup{(}bulk $\simeq$ derived center of boundary\textup{)} requires further hypotheses
\textup{(}compact generation, derived center quasi-isomorphism\textup{)} verified
in the boundary-linear exact sector
\textup{(}Theorem~\textup{\ref{thm:boundary-linear-bulk-boundary}}\textup{)};
the global case is conjectural, and an abstract bulk/Hochschild comparison for an arbitrary logarithmic $\SCchtop$-algebra without a chosen HT prefactorization realization is open
\textup{(}Section~\textup{\ref{sec:ht-bulk-boundary-line}} for the detailed scope table\textup{)}.
\end{theorem}

\begin{remark}[Categorical formulation and Morita invariance]
\label{rem:bulk-categorical}
Under the compact-generation hypotheses of the
boundary-linear exact sector, the chosen physical bulk algebra is
identified, by the bulk--Hochschild comparison, with the
Morita-invariant derived center of the open-sector category:
\[
\cO_{\mathrm{bulk}} \;\simeq\; Z_{\mathrm{der}}(\cC_\partial) \;:=\; \RHom_{\mathrm{Fun}(\cC_\partial, \cC_\partial)}(\id, \id),
\]
$\cC_\partial$ is the dg-category of boundary conditions (Definition~\ref{def:oc-factorization-category}). The identification is Morita invariant: it depends only on $\cC_\partial$, not on the choice of vacuum object $b$ presenting $A_\partial = \End(b)$. The passage from $Z_{\mathrm{der}}(\cC_\partial)$ to $C^\bullet_{\mathrm{ch}}(A_\partial, A_\partial)$ goes through the identification of the derived center of a compactly generated dg-category with the Hochschild cochain complex of any compact generator's endomorphism algebra~\cite{Kel06}; the integral-transform formalism of Ben-Zvi--Francis--Nadler~\cite{BZFN10} supplies the derived-algebraic-geometric framework.
\end{remark}

\begin{remark}[Scope of the bulk--Hochschild identification]
\label{rem:bulk-hochschild-scope}
The quasi-isomorphism
\(\Psi_{\mathrm{Ran}}\colon
\mathcal{O}_{\mathrm{bulk}}\simeq
C^\bullet_{\mathrm{ch}}(A_{\partial}, A_{\partial})\) is proved for
HT prefactorization algebras in the scope of
Theorem~\ref{thm:physics-bridge} via reduction along $\R$ and
factorization homology of the chosen prefactorization model
$\mathsf{Obs}$. The inverse homotopy class is
\(\chi_{\mathrm{HT},A_\partial}\). The stronger triangle composition
(bulk $\simeq$ derived center of boundary) requires compact generation
and the derived center quasi-isomorphism, verified in the
boundary-linear exact sector
(Theorem~\ref{thm:boundary-linear-bulk-boundary}). The global triangle
for all HT theories is conjectural; an abstract bulk/Hochschild
comparison for an arbitrary logarithmic $\SCchtop$-algebra without a
physical realization is open.
\end{remark}

\begin{remark}[Theorem content versus definitional identification]
\label{rem:bulk-hochschild-content}
In the Costello--Gwilliam framework, bulk observables are
factorization homology, which for algebras over the closed color
of $\SCchtop$ coincides with chiral Hochschild cochains by definition.
The content of Theorem~\ref{thm:bulk_hochschild} is the chain-level
comparison of the abstract factorization homology with the specific
cochain complex
$C^\bullet_{\mathrm{ch}}(A_\partial, A_\partial)$ built from the bar
construction, together with the explicit quasi-isomorphism
\(\Psi_{\mathrm{Ran}}\) intertwining the bulk
differential $d_{\mathrm{bulk}}$ with the chiral Hochschild
differential \(d_{\mathrm{Hoch}}=[D_{A_\partial},-]\).
The quasi-isomorphism \(\Psi_{\mathrm{Ran}}\) is defined via the Ran
space decomposition and verified by a filtration comparison on the
holomorphic weight; its inverse is the
\(\chi_{\mathrm{HT},A_\partial}\) comparison used by Universal
Holography. Chain-level control feeds the obstruction theory of
Corollary~\ref{cor:obstructions} and the
genus-$g$ Hochschild bridge of
Theorem~\ref{thm:hochschild-bridge-higher-genus}.
\end{remark}

\begin{proof}
The quasi-isomorphism is obtained by reduction along \(\R\), Ran
decomposition on \(\C\), and a complete bifiltration by arity and
holomorphic pole order.

\medskip
\textbf{Factorization homology and reduction along \(\R\).}

Bulk observables on $\C \times \R$ are computed by factorization homology,
\begin{equation}
\label{eq:bulk-fact-hom}
\mathcal{O}_{\mathrm{bulk}}(\C \times \R) \;\simeq\; \int_{\C \times \R} \mathsf{Obs}.
\end{equation}
Local constancy of $\mathsf{Obs}$ along $\R$ (Proposition~\ref{prop:prefact-Obs}) gives factorization homology over the product $\C \times \R$ a locally constant factorization along the topological direction~\cite{AF15}: for any interval $I \subset \R$ and disc $D \subset \C$, the restriction $\mathsf{Obs}(D \times I) \to \mathsf{Obs}(D \times I')$ is a quasi-isomorphism for each inclusion $I \hookrightarrow I'$. The $E_1$-structure along $\R$ is homotopically trivial on global sections; factorization homology reduces to
\begin{equation}
\label{eq:R-reduction}
\int_{\C \times \R} \mathsf{Obs} \;\simeq\; \int_{\C} \mathsf{Obs}^{\R},
\end{equation}
where $\mathsf{Obs}^{\R} := \mathsf{Obs}(D \times \R)$ denotes the $\R$-invariant (equivalently, $E_1$-coinvariant) part of the local observable complex, a holomorphic factorization algebra on $\C$. At the boundary $\{t = 0\} \times \C$ the restriction map
\begin{equation}
\rho : \mathcal{O}_{\mathrm{bulk}} \longrightarrow A_{\partial}
\end{equation}
identifies $A_{\partial}$ with $\mathsf{Obs}^{\R}|_{\{t=0\}}$, the boundary chiral algebra.

\medskip
\textbf{Arity and holomorphic pole filtrations.}

The factorization algebra $\mathsf{Obs}^{\R}$ on $\C$ admits a \emph{holomorphic weight filtration} $F^p$ by pole order at collision divisors: an element $f \in \int_{\C} \mathsf{Obs}^{\R}$ lies in $F^p$ if, on $n$-tuples of operator insertions at $z_1, \ldots, z_n \in \C$, the resulting meromorphic function has poles of total order $\le p$ along the collision diagonals $\{z_i = z_j\}$.

On $\Ran(\C)$ the factorization structure decomposes $\int_{\C} \mathsf{Obs}^{\R}$ into components labelled by the number of insertion points. An element supported on the stratum $\Ran_n(\C) \setminus \Ran_{n-1}(\C)$ of exactly $n$ distinct points is a multilinear operation
\begin{equation}
\label{eq:Ran-component}
f_n : A_{\partial}^{\otimes n} \longrightarrow A_{\partial}((\lambda_1)) \cdots ((\lambda_{n-1})),
\end{equation}
parametrized by the spectral parameters $\lambda_i$ (dual to the relative positions $z_i - z_n$). The collection $\{f_n\}_{n \geq 0}$ is an element of the chiral Hochschild cochain complex $C^\bullet_{\mathrm{ch}}(A_{\partial}, A_{\partial})$ (Definition~\ref{def:chiral_hochschild}).

The comparison preserves two filtrations. The Ran arity filtration
records the number of distinct insertion points. The holomorphic pole
filtration records the total order of singularity along the collision
divisors \(D_S\subset \FM_n(\C)\). These filtrations are independent:
arity fixes the stratum, while pole order fixes the singular depth of
the kernel on that stratum. The chiral Hochschild complex carries the
same complete bifiltration by arity and pole order.

\medskip
\textbf{Differential identification.}

The bulk differential on \(\int_{\C}\mathsf{Obs}^{\R}\) is carried by
the Ran decomposition to the commutator differential
\[
d_{\mathrm{bulk}}
\longmapsto
[D_{A_\partial},-]
\quad\text{on}\quad
\Coder_0\bigl(B^{\mathrm{ch}}(A_\partial)\bigr).
\]
The \(D_{m_1}\) summand is the internal BV-BRST differential.  The
summands \(D_{m_r}\), \(r\ge2\), are the codimension-one collision
face maps of \(\FM_{n+r-1}(\C)\): a divisor on which a consecutive
block of \(r\) points collides inserts \(m_r\) into the cochain, and
the dual term inserts the cochain into \(m_r\), with the block spectral
substitution of Definition~\ref{def:ainfty-relations}.  In the strict
PVA truncation this reduces to the familiar binary \(m_2\) Hochschild
formula.  In the \(\Ainf\)-chiral case the higher \(r\)-summands are
load-bearing; together they are exactly \([D_{A_\partial},-]\).
Thus the Ran space decomposition~\eqref{eq:Ran-component} produces a
chain map
\[
\int_{\C}\mathsf{Obs}^{\R}
\longrightarrow
C^\bullet_{\mathrm{ch}}(A_{\partial}, A_{\partial}).
\]

\medskip
\textbf{Quasi-isomorphism.}

The map
\begin{equation}
\Psi_{\mathrm{Ran}} :
\int_{\C} \mathsf{Obs}^{\R}
\longrightarrow
C^\bullet_{\mathrm{ch}}(A_{\partial}, A_{\partial})
\end{equation}
is a quasi-isomorphism by the bifiltration. On the associated graded
for arity and pole order, the factorization algebra
\(\gr\int_{\C}\mathsf{Obs}^{\R}\) splits by Ran strata: at arity \(p\)
and pole order \(q\) the summand is the space of \(p\)-linear kernels
\[
\Hom\!\bigl(A_{\partial}^{\otimes p},
A_{\partial}((\lambda_1))\cdots((\lambda_{p-1}))\bigr)
\]
whose total pole order is \(q\), with differential \(\delta_Q\). The
Hochschild collision part changes the arity-pole bidegree and hence
vanishes on the associated graded. The associated graded of
\(C^\bullet_{\mathrm{ch}}(A_{\partial},A_{\partial})\) has the same
description, so \(\gr(\Psi_{\mathrm{Ran}})\) is an isomorphism.
Completeness of the
arity-pole bifiltration and the spectral sequence comparison theorem
give
\begin{equation}
C^\bullet_{\mathrm{ch}}(A_{\partial}, A_{\partial}) \;\simeq\; \int_{\C} \mathsf{Obs}^{\R} \;\simeq\; \mathcal{O}_{\mathrm{bulk}}(\C \times \R),
\end{equation}
as dg-objects, $C^\bullet_{\mathrm{ch}}$ the chiral Hochschild
cochain complex. The inverse homotopy class of
\(\Psi_{\mathrm{Ran}}\) is \(\chi_{\mathrm{HT},A_\partial}\).
\end{proof}

\begin{remark}[Lagrangian self-intersection interpretation]
\label{rem:lagrangian-self-intersection}
\index{Lagrangian!self-intersection}
\index{HKR theorem!Lagrangian version}
Theorem~\ref{thm:bulk_hochschild} has a clean geometric
restatement. Write $\cL_b \subset \cM_{\mathrm{vac}}$ for the
Lagrangian in the $(-2)$-shifted symplectic stack of vacua
determined by the boundary condition~$b$
(corrected bulk--boundary--line triangle of
Section~\ref{sec:ht-bulk-boundary-line}).
The Hochschild--Kostant--Rosenberg theorem for the derived
self-intersection gives
\[
 \HH^\bullet(\cA_\partial,\cA_\partial)
 \;\simeq\;
 \cO\bigl(T^*[-1]\cL_b\bigr)
 \;\simeq\;
 \cO\bigl(\widehat{\cM_{\mathrm{vac}}}_{\cL_b}\bigr)_{\mathrm{Darboux}}\,,
\]
the second identification is a formal Darboux identification of the
formal neighbourhood of $\cL_b$ inside $\cM_{\mathrm{vac}}$. It is not
ordinary restriction: $\cO(\cM_{\mathrm{vac}})|_{\cL_b}=\cO(\cL_b)$ is
the boundary chart, while the bulk cochain algebra keeps the shifted
cotangent fibre. The scoped bulk-as-derived-center statement
(Theorem~\ref{thm:CDG_compatibility}\,(iii)) is therefore an algebra
statement in the boundary-linear exact sector: Hochschild cochains are
functions on the formal Lagrangian self-intersection
$\cL_b \times_{\cM_{\mathrm{vac}}}\cL_b$ in the Darboux chart, not
functions on the ambient stack itself. The Lagrangian
self-intersection is the chiral Steinberg object of
Theorem~\ref{thm:steinberg-presentation}; the HKR isomorphism
realises Vol~I Theorem~H as the fifth face of that
presentation (Remark~\ref{rem:steinberg-lagrangian-origin}).
\end{remark}

\begin{remark}[CY-two input shifts the HKR bar-cobar degree by two
and carries a projective Mathieu cocycle;
\ClaimStatusProvedElsewhere]
\label{rem:hoch-CY2-K3-shift}
\index{HKR theorem!K3 bi-graded specialisation}%
\index{Mathieu moonshine!Hochschild Serre cocycle}%
The HKR isomorphism of Remark~\ref{rem:lagrangian-self-intersection}
is stated for a $(-2)$-shifted symplectic ambient, matching a CY-$3$
closed-colour input (the Vol~III $\Phi$-baseline). When the closed
colour is a CY-\emph{surface} input, specifically the chiral algebra
$\cA_{K3}$ attached to the bounded derived category
$D^b\mathrm{Coh}(K3)$ via the Vol~III $\Phi$-functor
(Vol~III, Part~\ref{part:frontier}), the
Hochschild cohomology is bi-graded:
\[
\HH^\bullet\!\bigl(D^b\mathrm{Coh}(K3)\bigr)
\;\simeq\;
\bigoplus_{p+q=\bullet}
H^p\!\bigl(K3,\,\Lambda^q T_{K3}\bigr).
\]
Consequently the bar-cobar shift for this closed input is $[2]$,
equal to the CY dimension of the input, rather than the generic
$[3]$ of the Vol~III $\Phi$-baseline. The Lagrangian
self-intersection of Remark~\ref{rem:lagrangian-self-intersection}
then sees an additional dimensional shift from the K3
Calabi--Yau class: the $(-1)$-shifted symplectic structure on
$\cL_b\times^h_{\cM_{\mathrm{vac}}}\cL_b$ picks up the K3 Mukai
pairing rather than a strict intersection pairing, and the Serre
functor on $D^b\mathrm{Coh}(K3)$ carries a projective Schur cocycle
of order $6$ in $H^2\!\bigl(M_{24},\,U(1)\bigr)=\Z/12$ under
$\widetilde{M}_{24}$-equivariance (Mathieu moonshine,
Gaberdiel--Hohenegger--Volpato~\cite{GHV12}, building on
Bridgeland--Maciocia~\cite{BM01} on $K3$ autoequivalences). The
cocycle lifts to a non-trivial twist of the Lagrangian gluing data,
so the Hochschild bridge
(Theorem~\ref{thm:hochschild-bridge-higher-genus}) on the K3 input
computes a $\Z/12$-twisted Mukai pairing rather than an untwisted
self-intersection. This is the Vol~III $\Phi$-bridge data recorded
back on the bar-side HKR (compare
Remark~\ref{rem:bar-cobar-3-K3-shift2} of the bar-cobar
review, where the same content appears).
\end{remark}

\subsection{Comparison with CDG boundary chiral algebras}

The construction recovers and extends the boundary chiral algebra framework of Costello--Dimofte--Gaiotto~\cite{CDG20}.

\begin{theorem}[CDG compatibility; \ClaimStatusProvedHere]
\label{thm:CDG_compatibility}
Let $\cA$ be a logarithmic $\SCchtop$-algebra.
For 3d $\mathcal{N}=2$ theories with holomorphic-topological twist:
\begin{enumerate}[label=(\roman*)]
\item \textbf{Bulk is commutative PVA}: $H^\bullet(\mathcal{O}_{\text{bulk}}, Q)$ is a commutative Poisson vertex algebra with $(-1)$-shifted $\lambda$-bracket (Theorem~\ref{thm:cohomology_PVA});

\item \textbf{Boundary is vertex algebra}: $A_{\partial}$ is a vertex algebra (with or without Virasoro);

\item \textbf{Bulk-to-boundary map}: a canonical morphism of PVAs
\begin{equation}
\beta : H^\bullet(\mathcal{O}_{\text{bulk}}, Q) \longrightarrow Z(A_{\partial}),
\end{equation}
$Z(A_{\partial})$ the center of $A_{\partial}$ (elements commuting with all $\lambda$-brackets);

\item \textbf{Boundary is bulk module}: $A_{\partial}$ is a module for the bulk PVA $H^\bullet(\mathcal{O}_{\text{bulk}}, Q)$ via $\beta$.
\end{enumerate}
\end{theorem}

\begin{proof}
(i) and (ii) are \cite[\S 4]{CDG20}. (iii) follows from the bulk--Hochschild identification (Theorem~\ref{thm:bulk_hochschild}): the center $Z(A_{\partial})$ is $\ChirHoch^0(\cA_\partial,\cA_\partial)$ (degree-$0$ chiral Hochschild cochains in the geometric FM model; Remark~\ref{rem:thm-H-three-way-scope}), distinct from the classical mode-algebra $\HH^0(\cA_{\partial,\mathrm{mode}})$. (iv) is the module structure induced by the bulk-to-boundary restriction map.
\end{proof}

\begin{remark}[Theorem~H: three-way scope]%
\label{rem:thm-H-three-way-scope}%
\index{Theorem H!three-way scope|textbf}%
\index{Hochschild!three-way scope|textbf}%
Every invocation of Theorem~H in this volume refers to \emph{chiral} Hochschild
$\ChirHoch^\bullet(\cA,\cA)$, constructed either geometrically on $\FM_{n+1}(X)$ with
logarithmic forms or algebraically via $\End^{\mathrm{ch}}_\cA$; the two models are
naturally quasi-isomorphic for logarithmic $\SCchtop$-algebras
(Proposition~\ref{thm:chiral-hochschild-models-equivalent}). The
notation resolves into three typed objects:
\begin{enumerate}
 \item[\textup{(i)}] $\ChirHoch^\bullet(\cA,\cA)$: a family datum
 $H_H(\cA;S)$, for finite $S\subset\mathbb Z$, consists of a complete
 filtered chart complex~$Q_\cA$, a
 quasi-isomorphism
 $\gamma_\cA\colon Q_\cA\xrightarrow{\sim}
 C^\bullet_{\mathrm{ch}}(\cA,\cA)$, and maps
 \[
 K_{\cA,S}\xrightarrow{\iota_\cA}Q_\cA
 \xrightarrow{p_\cA}K_{\cA,S},
 \qquad
 p_\cA\iota_\cA=\id,
 \qquad
 d h_\cA+h_\cA d=\id-\iota_\cA p_\cA,
 \]
 where $K_{\cA,S}$ is complete and has cochain terms in~$S$, and
 $h_\cA$ has degree~$-1$.
 The datum also contains the finite-window inverse-limit comparison
 and the ordered-to-symmetric averaging quasi-isomorphism. Volume~I
 Theorem~H then gives
 \[
  \ChirHoch^n(\cA)\cong H^n(K_{\cA,S}),
  \qquad
  \operatorname{Supp}\ChirHoch^\bullet(\cA)\subseteq S.
 \]
 The $E_3$-topological enhancement is supplied separately by the
 chiral Higher-Deligne and topologisation hypotheses.  Degree~$0$ is
 the ordinary chiral centre; the derived chiral centre is the full
 object
 $Z^{\mathrm{der}}_{\mathrm{ch}}(\cA)=\ChirHoch^\bullet(\cA,\cA)$.
 \item[\textup{(ii)}] $\HH^\bullet(\cA_{\mathrm{mode}},\cA_{\mathrm{mode}})$:
 classical Hochschild cohomology of the mode algebra $\cA_{\mathrm{mode}}\in\mathrm{Vect}$.
 For simple mode algebras (e.g.\ Weyl) this is concentrated in degree~$0$ with
 $\HH^0\cong k$; Theorem~H controls item~\textup{(i)}.
 \item[\textup{(iii)}] $H^\bullet_{\mathrm{GF}}(\mathrm{Lie}(\cA))$:
 Gel'fand--Fuchs continuous cohomology of the mode Lie algebra, a third
 functor with its own typically unbounded polynomial.
\end{enumerate}
The BZFN equivalence $Z(\mathrm{LMod}_A(S))\simeq\mathrm{LMod}_{\HH^\bullet(A,A)}(S)$
\textup{(}Lurie HA~5.3.1.30\textup{)} uses the same ambient~$S$ on
both sides. The two derived centres arise from applying BZFN to
\emph{two different algebras} in their native ambient
categories: $\cA$ in $\mathrm{IndCoh}(\mathrm{Ran}(X))$ giving (i), and
$\cA_{\mathrm{mode}}$ in $\mathrm{Vect}$ giving (ii). The comparison
map $\eta_{\mathrm{mode}}$
\textup{(}Theorem~\ref{thm:three-hochschild-unification}\textup{)}
is a degree-$0$ isomorphism on the non-critical scalar-centre
locus. At critical level $k=-h^\vee$ the chiral centre jumps to the
Feigin--Frenkel polynomial algebra while
$\HH^\bullet(\cA_{\mathrm{mode}})$ stays finite; the degree-$0$
comparison becomes the Feigin--Frenkel zero-mode quotient with
infinite-dimensional kernel.
\end{remark}

\subsection{The Zhu centre shadow
\texorpdfstring{$\mathrm{Zhu}^0\colon Z_{\mathrm{ch}}(V)\to Z(A(V))$}{Zhu centre shadow}}
\label{subsec:zhu-bridge}

Remark~\ref{rem:thm-H-three-way-scope} unifies the three Hochschild
invariants through $\eta_{\mathrm{mode}}$ (Theorem~\ref{thm:three-hochschild-unification}),
which lands in the full mode algebra $\cA_{\mathrm{mode}}$
(a topological algebra on the generating modes $a_{(n)}$). A
strictly smaller, more classical target is the Zhu algebra
$A(V)$ of Zhu~\cite{Zhu96}: the $C_2$-truncated zero-mode quotient
capturing the admissible representation theory of $V$. The zero-mode
bridge is precise only at degree \(0\): it records the image of chiral
central fields in the Zhu centre. It is not a functorial comparison of
the whole chiral Hochschild complex with \(\HH^\bullet(A(V))\).

\begin{definition}[Zhu algebra $A(V)$ and Zhu bridge; Zhu 1996]%
\label{def:zhu-bridge}%
\index{Zhu algebra $A(V)$|textbf}%
\index{Zhu bridge!definition}%
Let $V=\bigoplus_{n\ge 0}V_n$ be a $\Z_{\ge 0}$-graded vertex algebra
with conformal vector. Define
\[
 O(V) \;:=\; \mathrm{span}_k\Bigl\{
 \Res_z\; Y(a,z)\,b\,\frac{(1+z)^{\wt(a)}}{z^{2}}
 \;\Big|\; a,b\in V\Bigr\},
 \qquad
 A(V) \;:=\; V\big/ O(V),
\]
with associative product
\[
 a *_0 b \;:=\; \Res_z\; Y(a,z)\,b\;\frac{(1+z)^{\wt(a)}}{z}.
\]
Let \(\mathcal U(V)\) be the Frenkel--Zhu mode algebra and
\(\mathcal U(V)_0\) its degree-zero part.  The Zhu algebra is the
degree-zero quotient
\[
A(V)\;\simeq\;\mathcal U(V)_0/I_{\mathrm{Zhu}},
\qquad
I_{\mathrm{Zhu}}=(\mathcal U(V)\mathcal U(V)_{<0})\cap
\mathcal U(V)_0,
\]
not a quotient of the full mode algebra \(\mathcal U(V)\).  The
assignment \(a_{(n)}\mapsto o(a)\) for \(n=\wt(a)-1\) and \(0\)
otherwise is therefore not an algebra homomorphism on the full mode
algebra; it is only the zero-mode shadow after passing through
\(\mathcal U(V)_0\).

The \emph{degree-zero Zhu centre shadow} is
\[
\mathrm{Zhu}^0\colon
Z_{\mathrm{ch}}(V)=\ChirHoch^0(V)
\longrightarrow Z(A(V)),
\qquad
z\longmapsto [o(z)].
\]
Any comparison with \(\HH^{>0}(A(V))\) requires extra bimodule or
derived Morita data; it is not induced functorially from
\(\eta_{\mathrm{mode}}\) by a quotient of the full mode algebra.
\end{definition}

\begin{proposition}[Zhu centre shadow; \ClaimStatusProvedHere]%
\label{prop:zhu-hochschild-comparison}%
\index{Zhu bridge!comparison theorem|textbf}%
\index{Zhu bridge!sector refinement for chiral Hochschild}%
Let $V$ be as in Definition~\ref{def:zhu-bridge}. Then:
\begin{enumerate}
 \item[\textup{(i)}] \emph{Degree-zero shadow.}
 \(\mathrm{Zhu}^0\) lands in \(Z(A(V))\). Its kernel is
 \(Z_{\mathrm{ch}}(V)\cap O(V)\), equivalently the chiral central
 fields whose zero-mode class is killed by the Zhu ideal. The vacuum
 maps to \(1\in A(V)\), so the map is nonzero.
 \item[\textup{(ii)}] \emph{Rational \(C_2\)-cofinite case.}
 If \(V\) is simple, rational, and \(C_2\)-cofinite, then
 \(A(V)\) is finite-dimensional semisimple
 \cite{Zhu96,DLM98}:
 \[
 A(V)\;\simeq\;\bigoplus_{i\in I(V)}\End(M_i(0)).
 \]
 Hence
 \[
 \HH^0(A(V))=Z(A(V))\simeq \C^{I(V)},\qquad
 \HH^j(A(V))=0\quad(j>0).
 \]
 The degree-zero image of a simple \(V\) is the scalar diagonal
 \(\C\cdot(1_i)_{i\in I(V)}\subset \C^{I(V)}\). Therefore the
 degree-zero Zhu shadow is an isomorphism exactly when
 \(|I(V)|=1\); for the Ising model
 \(\Vir_{1/2}\), it is the strict map \(\C\to\C^3\). It is not a
 quasi-isomorphism \(\ChirHoch^{\le2}(V)\to\HH^{\le2}(A(V))\) in
 general.
 \item[\textup{(iii)}] \emph{Critical affine level.}
 For \(V=V_{-h^\vee}(\fg)\), Frenkel--Zhu gives
 \(A(V)\simeq U(\fg)\). The zero-mode map from the
 Feigin--Frenkel centre
 \[
 \mathfrak z(\widehat{\fg})=
 \ChirHoch^0(V_{-h^\vee}(\fg))
 \]
 to \(Z(U(\fg))\) is the Feigin--Frenkel--Reshetikhin
 one-point/Gaudin specialisation followed by Harish--Chandra
 projection. It is surjective and has infinite-dimensional kernel;
 it is not an inclusion.
 \item[\textup{(iv)}] \emph{Lost chiral information.}
 The information invisible to \(A(V)\) is the kernel of the zero-mode
 shadow together with any higher chiral Hochschild classes for which
 no supplied bimodule comparison to \(\HH^\bullet(A(V))\) has been
 constructed. For logarithmic \(W(p)\), critical affine algebras, and
 non-\(C_2\)-cofinite admissible/generic families, the Zhu algebra is
 a representation-theoretic shadow, not a replacement for
 \(\ChirHoch^\bullet(V)\).
\end{enumerate}
\end{proposition}

\begin{proof}
\emph{(i).} If \(z\in Z_{\mathrm{ch}}(V)\), then the zero mode
\(o(z)\) commutes with all zero modes in Zhu's product
\cite[Theorem~2.1.2]{Zhu96}. Hence \([o(z)]\in Z(A(V))\). The kernel
is exactly the intersection with the quotient ideal \(O(V)\), because
\(A(V)=V/O(V)\). The vacuum class is not in \(O(V)\) and maps to the
identity.

\emph{(ii).} For a simple rational \(C_2\)-cofinite vertex operator
algebra, Zhu and Dong--Li--Mason identify \(A(V)\) as a finite
semisimple algebra whose simple blocks are indexed by the simple
\(V\)-modules \cite{Zhu96,DLM98}. A finite-dimensional semisimple
algebra over \(\C\) is separable, so its Hochschild cohomology is its
centre in degree \(0\) and vanishes in positive degree.  Since a
simple \(V\) has chiral centre \(\C\cdot\mathbf1\), the image is the
scalar diagonal. The Ising model has three simple modules
\((h=0,\frac12,\frac1{16})\), so the degree-zero map is
\(\C\to\C^3\), not an isomorphism.

\emph{(iii).} Frenkel--Zhu gives \(A(V^k(\fg))\simeq U(\fg)\) for the
universal affine vertex algebra, in particular at \(k=-h^\vee\)
\cite{FrenkelZhu92}. The Feigin--Frenkel centre at the critical level
is the polynomial algebra of functions on opers
\(\Fun(\mathrm{Op}_{\fg^\vee}(D))\)
\cite{FeiginFrenkel1992,Frenkel2007}. The zero-mode/Gaudin
specialisation of Feigin--Frenkel--Reshetikhin sends these central
fields to the finite Harish--Chandra centre \(Z(U(\fg))\)
\cite{FFR94}. It is onto the rank-\(\fg\) Harish--Chandra generators,
while the source has infinitely many derivative/jet generators; hence
the kernel is infinite-dimensional and no injection can exist.

\emph{(iv).} The quotient \(V\to A(V)\) kills \(O(V)\) and forgets all
nonzero-mode and higher Hochschild data unless additional comparison
data is supplied. The listed families are precisely loci where this
forgotten part is load-bearing: logarithmic extensions retain
pseudotrace and Jordan data, critical affine algebras retain the
Feigin--Frenkel oper directions, and non-\(C_2\)-cofinite/generic
families need the full chiral Hochschild complex rather than the
finite zero-mode shadow.
\end{proof}

\begin{remark}[Recoding the ``Zhu correspondence'' folklore]%
\label{rem:zhu-folklore-recode}%
Statements of the form ``$\HH_0(V)\simeq \mathrm{Tr}_{A(V)}(\mathrm{Id})$
by the Zhu correspondence'' (in particular the Chern--Simons--Kac--Moody
invocation at \eqref{eq:cs-annulus-trace}) are shorthand for
the degree-zero Zhu centre shadow of
Proposition~\ref{prop:zhu-hochschild-comparison}, not for a
quasi-isomorphism from chiral Hochschild cochains to
\(\HH^\bullet(A(V))\).  In the rational \(C_2\)-cofinite case, the
trace over \(A(V)\) sees the finite semisimple block centre
\(\C^{I(V)}\). For critical-level \(V_{-h^\vee}(\fg)\), the trace over
\(U(\fg)\) captures only the Harish--Chandra centre \(Z(U(\fg))\),
whereas
\(\HH^{\mathrm{ch}}_0(V_{-h^\vee}(\fg))=\mathfrak{z}(\widehat{\fg})\)
is the Feigin--Frenkel oper algebra and maps onto \(Z(U(\fg))\) with
infinite kernel. Every future invocation of the ``Zhu correspondence''
must cite Proposition~\ref{prop:zhu-hochschild-comparison} and declare
whether it means the degree-zero centre shadow, the semisimple
block-trace target, or an explicitly supplied derived Morita
comparison.
\end{remark}

\subsection{\texorpdfstring{Chiral Hochschild cochains as a brace algebra on $\FM(\C)\times\Conf(\R)$}{Chiral Hochschild cochains as a brace algebra on FM(C) x Conf(R)}}
\label{sec:chiral-Hochschild}

\subsubsection{Operadic cochains for colored operads and the brace structure}
\label{subsec:operadic-cochains}
Let $O$ be a colored dg-operad and $A$ an $O$-algebra in a symmetric monoidal dg-category $(\mathcal C,\otimes)$.
Write $O^\vee$ for the bar cooperad of $O$ (cofree conilpotent cooperad on $s^{-1}\overline O$).
The operadic cochain complex is
\[
\mathrm{C}^\bullet_O(A) \;:=\; \underline{\mathrm{Hom}}_{\text{cooperad}} \bigl( O^\vee,\, \mathrm{End}_A \bigr),
\]
with differential from the coderivation on $O^\vee$ and the internal differential of $\mathrm{End}_A$.
For $O = \OHT$ the Swiss-cheese chain operad, write
\[
\mathrm{C}^\bullet_{\mathrm{ch,top}}(A)\;:=\;\mathrm{C}^\bullet_{\OHT}(A)
\quad\text{(the \emph{chiral Hochschild cochains})}.
\]

\begin{definition}[Brace operations]
\label{def:brace-ops}
The bar cooperad $O^\vee$ carries a canonical brace structure; by precomposition,
$\mathrm{C}^\bullet_O(A)$ acquires multilinear operations
\[
\{-\}\{-,\ldots,-\}\colon \mathrm{C}^\bullet_O(A)\otimes \mathrm{C}^\bullet_O(A)^{\otimes m}\longrightarrow \mathrm{C}^\bullet_O(A),
\]
given by grafting trees in $O^\vee$ (\cite{GJ94,Vor99}).
The structure on $\mathrm{C}^\bullet_O(A)$ is a homotopy Gerstenhaber (brace) algebra with Gerstenhaber bracket
\[
[f,g] := f\{g\} - (-1)^{(\deg f-1)(\deg g-1)}\, g\{f\}.
\]
\end{definition}

\begin{remark}[Degree conventions]
Cohomological grading. For $O=\OHT$ the internal degrees reflect the form-degrees on $\FM(\C)$ and the $E_1$-degrees on $\Conf(\R)$; the brace degree shifts match the standard Gerstenhaber sign rule.
\end{remark}

\subsubsection{\texorpdfstring{Geometric brace model on $\FM(\C)\times\Conf(\R)$}{Geometric brace model on FM(C) x Conf(R)}}
\label{subsec:geometric-brace}
A concrete model of the brace operations uses configuration spaces.
Let $\Conf_k(\R)$ denote ordered configurations $(t_1<\cdots<t_k)$ of $k$ points in $\R$ and set
\[
\mathbb{X}_{k,m} \;:=\; \FM_k(\C)\times \Conf_m(\R).
\]
Operations in $\OHT$ at degree $(k,m)$ are chains on $\mathbb{X}_{k,m}$ decorated by edge-length parameters.
For $f\in \mathrm{C}^\bullet_{\mathrm{ch,top}}(A)$ and $g_1,\dots,g_m\in \mathrm{C}^\bullet_{\mathrm{ch,top}}(A)$, set
\[
f\{g_1,\dots,g_m\} := \int_{\mathbb{X}} \;\;\Bigl(\text{pullback of kernels from } \mathbb{X}_{k,m'} \Bigr)\;\circ\;\Bigl(\text{insertions via trees along } \FM\times \Conf\Bigr),
\]
the integral over the fiber product of configuration spaces determined by the brace tree (Appendix~\ref{app:brace-signs} for the sign and degree formulas).

\begin{proposition}[Well-definedness and functoriality;
\ClaimStatusProvedHere]
\label{prop:geometric-braces-well-defined}
For a logarithmic $\SCchtop$-algebra (Definition~\ref{def:log-SC-algebra}), the geometric brace operations on $\mathrm{C}^\bullet_{\mathrm{ch,top}}(A)$ are well-defined, strictly multilinear, and satisfy the brace identities. The Gerstenhaber bracket obtained from these braces matches the one induced by the cooperad $O^\vee$.
\end{proposition}

\begin{proof}
Integrability is the logarithmic form condition (Definition~\ref{def:log-SC-algebra}): weight forms have logarithmic singularities along the boundary divisors of the FM compactification, integrable on compact manifolds with corners. The brace identities follow from Stokes on $\FM_k(\C) \times \Conf_m(\R)$: boundary strata reproduce the brace relations, with Arnold cancellations at codimension-$2$ corners (Appendix~\ref{app:brace-signs} for the sign verification). Compatibility with the cooperad differential is the standard operadic cochain construction.
\end{proof}

\subsubsection{The \texorpdfstring{$\ast$}{*}-Lie structure and deformation theory}
\label{subsec:star-Lie}
Let $(\mathrm{C}^\bullet_{\mathrm{ch,top}}(A),\{-\}\{-\})$ be the brace algebra above.

\begin{definition}[$\ast$-Lie algebra]
\label{def:star-Lie}
The shifted Lie bracket on $\mathrm{C}^\bullet_{\mathrm{ch,top}}(A)[1]$ is the Gerstenhaber bracket $[-,-]$ associated to braces. The pair $(\mathrm{C}^\bullet_{\mathrm{ch,top}}(A)[1],[-,-])$ is the \emph{chiral $\ast$-Lie algebra} of $A$.
\end{definition}

\begin{theorem}[Deformations and Maurer--Cartan; \ClaimStatusProvedHere]
\label{thm:MC-deformations}
For $A$ a $\OHT$-algebra and formal deformations over $k[\![\hbar]\!]$, a bijection
\begin{enumerate}[label=(\alph*)]
 \item deformation classes of $A$ as a $W(\SCchtop)$-algebra over $k[\![\hbar]\!]$ up to gauge;
 \item gauge-equivalence classes of solutions to the Maurer--Cartan equation
 \(
 d\alpha + \tfrac{1}{2}[\alpha,\alpha]=0
 \)
 in $\bigl(\mathrm{C}^\bullet_{\mathrm{ch,top}}(A)[1]\otimes \hbar k[\![\hbar]\!],[-,-]\bigr)$.
\end{enumerate}
\end{theorem}

\begin{proof}
Deformations are classified by twisting morphisms from the bar cooperad into $\mathrm{End}_A$; a twisting morphism is an MC element in the homotopy convolution object associated to $(O^\vee,\mathrm{End}_A)$. On a bar-cofree model this homotopy object is strictified by the convolution Lie algebra $\mathrm{Hom}(O^\vee,\mathrm{End}_A)$ (\cite{GK94,Val07}). The brace model identifies this strictification with $\mathrm{C}^\bullet_{\mathrm{ch,top}}(A)[1]$.
\end{proof}

\subsubsection{\texorpdfstring{Bulk observables $\simeq$ chiral Hochschild cochains}{Bulk observables equivalent to chiral Hochschild cochains}}
\label{subsec:bulk-equals-CHC}
Let $\mathsf{Obs}$ be the BV observable prefactorization algebra of an HT theory in the scope of Theorem~\ref{thm:physics-bridge}, and $A=\mathsf{Obs}(D\times I)$ the local complex on a small rectangle.

\begin{theorem}[Bulk--cochains identification; \ClaimStatusProvedHere]
\label{thm:bulk-CHC}
A canonical filtered quasi-isomorphism
\[
\mathrm{C}^\bullet_{\mathrm{ch,top}}(A)\;\simeq\;\mathsf{Obs}_{\mathrm{bulk}}(\mathbb{B}^3)
\]
between chiral Hochschild cochains of the local algebra $A$ and the complex of bulk local observables supported in a small $3$-ball, with filtration by holomorphic weight. On associated graded this identification reduces to the BD-chiral Hochschild complex on the closed color tensored with the $E_1$ Hochschild complex on the open color.
\end{theorem}

\begin{proof}
Both sides are computed by the same operad of holomorphic-topological
shapes, with the same logarithmic kernels.

\medskip
\textbf{Operad of shapes.}

Decompose the small $3$-ball $\mathbb{B}^3 \subset \C \times \R$ into rectangles $D_i \times I_j$ (discs in $\C$ times intervals in $\R$). On each rectangle the topological direction is locally constant, while the holomorphic direction retains the logarithmic chiral kernel; the local complex is the fixed rectangle model \(A=\mathsf{Obs}(D\times I)\) of Proposition~\ref{prop:prefact-Obs}. The factorization homology $\int_{\mathbb{B}^3} \mathsf{Obs}$ is the homotopy colimit over all such decompositions; the indexing category is the operad of shapes $\mathbb{X}_{k,m} = \FM_k(\C) \times \Conf_m(\R)$, with $k$ holomorphic and $m$ ordered topological insertion points parametrising embeddings of rectangles in $\mathbb{B}^3$.

An element of $\mathsf{Obs}_{\mathrm{bulk}}(\mathbb{B}^3)$ is a compatible family of multilinear operations
\[
A^{\otimes k} \otimes A^{\otimes m} \longrightarrow A,
\]
parametrized by points in $\mathbb{X}_{k,m}$ — integration of propagator kernels over $\FM_k(\C) \times \Conf_m(\R)$. By the product-formal Swiss-cheese local-shadow theorem (Theorem~\ref{thm:recognition-foundations}), these operations assemble into the local $W(\SCchtop)$-algebra shadow, and the totality is the operadic cochain complex
\[
\mathrm{C}^\bullet_{\mathrm{ch,top}}(A) = \underline{\mathrm{Hom}}_{\text{cooperad}}\bigl((\OHT)^\vee, \mathrm{End}_A\bigr).
\]
The map $\Phi \colon \mathrm{C}^\bullet_{\mathrm{ch,top}}(A) \to \mathsf{Obs}_{\mathrm{bulk}}(\mathbb{B}^3)$: a cochain $f$ of degree $(k,m)$ determines an observable by integration against the propagator kernels over $\mathbb{X}_{k,m}$. Restriction of a bulk observable to the configuration-space strata recovers the cochain components.

\medskip
\textbf{Holomorphic weight filtration.}

Both sides carry the holomorphic weight filtration $F^p$, graded by total pole order at collision divisors in $\FM_k(\C)$.

On the cochain side, $F^p \mathrm{C}^\bullet_{\mathrm{ch,top}}(A) \subset \mathrm{C}^\bullet_{\mathrm{ch,top}}(A)$ is the subcomplex of cochains whose kernels, restricted to $\FM_k(\C)$, have poles of total order $\le p$ along $\bigcup_{|S| \geq 2} D_S$.

On the observable side, $F^p \mathsf{Obs}_{\mathrm{bulk}}(\mathbb{B}^3)$ is defined by the analogous pole-order condition on the Feynman-diagrammatic integrands.

The map $\Phi$ is filtered: it preserves pole order since it is integration over the same configuration-space kernels. On associated graded, the factorization structure splits by $\mathbb{X}_{k,m} = \FM_k(\C) \times \Conf_m(\R)$. At each graded piece:
\begin{enumerate}[label=(\roman*)]
\item the chiral component (varying $k$, from $\FM_k(\C)$) is the BD-chiral Hochschild complex $\mathrm{C}^\bullet_{\mathrm{BD\text{-}ch}}(A_{\mathsf{ch}})$, $A_{\mathsf{ch}}$ the closed-color (holomorphic) part of the $W(\SCchtop)$-algebra. The differential is the chiral Hochschild differential with $\delta_Q$ from the internal BV-BRST operator and $\delta_{\mathrm{Hoch}}$ from the boundary faces of $\FM_k(\C)$ encoding OPE collisions;
\item the topological component (varying $m$, from $\Conf_m(\R)$) is the $E_1$ Hochschild complex $\mathrm{C}^\bullet_{E_1}(A_{\mathsf{top}})$, $A_{\mathsf{top}}$ the open-color (topological) part; the differential encodes boundary faces of $\Conf_m(\R)$ (merging adjacent intervals).
\end{enumerate}
The splitting on associated graded is
\[
\gr^F \mathrm{C}^\bullet_{\mathrm{ch,top}}(A) \;\cong\; \mathrm{C}^\bullet_{\mathrm{BD\text{-}ch}}(A_{\mathsf{ch}}) \;\widehat\otimes\; \mathrm{C}^\bullet_{E_1}(A_{\mathsf{top}}),
\]
the same decomposition holds for $\gr^F \mathsf{Obs}_{\mathrm{bulk}}(\mathbb{B}^3)$ by the factorization product structure on disjoint rectangles. So $\gr(\Phi)$ is an isomorphism.

\medskip
\textbf{Integrability.}

The integrals defining $\Phi$ converge by Theorem~\ref{thm:physics-bridge} (factorized logarithmic parametrix and one-loop finiteness) and Theorem~\ref{thm:FM-calculus}(a) (logarithmic form integrability): the holomorphic kernels are logarithmic in $\C$ with at most simple poles at collision points, and the topological kernels are locally constant or exponentially decaying along $\R$. After pullback to the FM compactification $\FM_k(\C) \times \Conf_m(\R)$, the weight forms defining $\Phi$ are logarithmic forms on $\FM_k(\C)$ (Theorem~\ref{thm:FM-calculus}(a)) with integrable singularities on the compact manifold with corners $\FM_k(\C)$; the fiber integrals converge absolutely. The key estimate is
\[
\int_0^\epsilon |\log \varepsilon_S|^N \, \varepsilon_S^{a-1} \, d\varepsilon_S < \infty \quad \text{for } a > 0, \; N \geq 0,
\]
$\varepsilon_S$ the radial parameter measuring the size of the cluster $S$ (Definition~\ref{def:FM_coords}). The exponent $a > 0$ is the logarithmic simple-pole condition in the factorized parametrix; higher-loop contributions are controlled by one-loop finiteness (Theorem~\ref{thm:physics-bridge}(c)).

The filtration is complete and Hausdorff on both sides (pole order is bounded below; the complexes are exhausted by finite filtration levels). The spectral sequence comparison theorem lifts the isomorphism on $\gr$ to a quasi-isomorphism on the filtered complexes,
\[
\Phi \colon \mathrm{C}^\bullet_{\mathrm{ch,top}}(A) \;\xrightarrow{\;\simeq\;}\; \mathsf{Obs}_{\mathrm{bulk}}(\mathbb{B}^3). \qedhere
\]
\end{proof}

\begin{corollary}[Obstruction theory; \ClaimStatusProvedHere]
\label{cor:obstructions}
For the controlling dg Lie algebra
$\mathfrak g_A=\mathrm{C}^\bullet_{\mathrm{ch,top}}(A)[1]$,
infinitesimal deformations are parameterized by
$H^1(\mathfrak g_A)$ and primary obstructions lie in
$H^2(\mathfrak g_A)$. Equivalently, in the unshifted cochain grading,
infinitesimals lie in
$H^2(\mathrm{C}^\bullet_{\mathrm{ch,top}}(A))$ and primary
obstructions lie in
$H^3(\mathrm{C}^\bullet_{\mathrm{ch,top}}(A))$. Higher obstructions
live in the Massey--product tower controlled by the brace operations.
\end{corollary}

\subsubsection{Filtration and PVA on associated graded}
\label{subsec:filtration-PVA}
Let $F^\bullet$ be the holomorphic weight filtration on $\mathrm{C}^\bullet_{\mathrm{ch,top}}(A)$. Then
\[
\mathrm{gr}^F\, \mathrm{C}^\bullet_{\mathrm{ch,top}}(A)
\;\cong\;
\mathrm{C}^\bullet_{\mathrm{BD\text{-}ch}}(A_{\mathsf{ch}})\;\widehat\otimes\; \mathrm{C}^\bullet_{E_1}(A_{\mathsf{top}}),
\]
and the induced bracket on $\mathrm{H}^\bullet$ coincides with the $(-1)$--shifted Poisson vertex bracket of Section~\ref{sec:Ainfty-to-PVA}.

\begin{theorem}[Hochschild bridge to Volume~I at genus~$0$;
\ClaimStatusProvedHere]
\label{thm:hochschild-bridge-genus0}
Let $\cA$ be the BV-BRST complex of a
3d HT theory on $\C \times \R$, and assume $\cA$ is chiral Koszul
(Volume~I, Theorem~B). Then there is a canonical filtered
quasi-isomorphism
\[
\mathrm{C}^\bullet_{\mathrm{ch,top}}(\cA)
\;\simeq\;
\ChirHoch^\bullet(\cA),
\]
where $\ChirHoch^\bullet(\cA)$ is the chiral Hochschild complex of
Volume~I on the curve $X = \C$.  If $\cA$ also carries a family
support datum $H_H(\cA;S)$, then
\[
 H^n\!\left(\mathrm C^\bullet_{\mathrm{ch,top}}(\cA)\right)
 \cong H^n(K_{\cA,S}),
 \qquad
 \operatorname{Supp}H^\bullet
 \!\left(\mathrm C^\bullet_{\mathrm{ch,top}}(\cA)\right)
 \subseteq S.
\]
\end{theorem}

\begin{proof}
The 3d HT theory on \(\C\times\R\) is translation-equivariant in the
holomorphic direction: the action, propagator, and BV-BRST
differential are invariant under \(z\mapsto z+c\). This gives
canonically identified formal-disc models at different points of
\(\C\), but it does not make the chiral factorization algebra locally
constant in the holomorphic direction; the logarithmic singularities
along collision diagonals remain part of the structure.

\textbf{Formal-disc chiral descent.}
For a translation-covariant holomorphic factorization algebra on
\(\C\) with logarithmic kernels, Beilinson--Drinfeld chiral descent
computes the global chiral Hochschild complex from the formal
neighbourhoods of the collision diagonals \cite[Ch.~3]{BD04}. The
formal neighbourhood of
a collision in \(\C\) is identified with the formal neighbourhood of the
same collision in any small disc \(D\ni z_0\). Thus the stalk of
$\ChirHoch^\bullet(\cA)$ at $z_0$ is computed by restricting to a
small disc $D \ni z_0$:
\[
\ChirHoch^\bullet(\cA)|_{z_0}
\;\simeq\;
\bigoplus_{n \geq 0}
 \Gamma\!\bigl(
 C_{n+2}(D),\,
 j_\ast j^\ast \cA^{\otimes(n+2)}
 \otimes \Omega^n_{\log}
 \bigr).
\]
This is exactly the local complex
$\mathrm{C}^\bullet_{\mathrm{ch,top}}(\cA)$ of
Theorem~\ref{thm:bulk-CHC}, restricted from $\FM_{n+2}(\C)$ to
$\FM_{n+2}(D)$. Chiral excision for formal discs identifies these
logarithmic de Rham complexes; the restriction is therefore a
quasi-isomorphism.

\textbf{Filtration and Hilbert structure.}
Both sides carry the holomorphic weight filtration $F^p$ counting
pole order along collision divisors in $\FM_k(\C)$. The
formal-disc restriction preserves pole order, so it is a
\emph{filtered} quasi-isomorphism.  Under $H_H(\cA;S)$, Volume~I
Theorem~H identifies the chart complex with $K_{\cA,S}$ and computes
its support~$S$; the filtered bridge transports this family model to
the Swiss-cheese complex.
\end{proof}

\begin{remark}[Connes periodicity]
Under the identification of
Theorem~\ref{thm:hochschild-bridge-genus0}, the Connes periodicity
operator $S$ on $\ChirHoch^\bullet$ (shifting cyclic degree by
$-2$) corresponds to the holomorphic weight shift
$F^p \to F^{p-1}$ on
$\mathrm{C}^\bullet_{\mathrm{ch,top}}(\cA)$. This follows from
the fact that $S$ acts on configuration spaces by ``forgetting one
insertion point,'' which lowers the number of collision divisors
and hence the pole order by one unit in each factor. We record
this as a consequence rather than an ingredient of the proof.
\end{remark}

\begin{remark}[Periodic-CDG mechanism]
\label{rem:periodic-cdg-mechanism}
The periodic-CDG conjecture asserts that the curved
$A_\infty$-structure on the chiral bar at genus~$g \geq 1$
descends to a \emph{periodic} curved dg-category. We record its
first-principles anatomy here, separating what is established,
what is conjectural, and what is a terminology hazard.

\medskip
\noindent\textbf{Origin.} At genus $g \geq 1$ the bar complex
carries scalar curvature
$\operatorname{pr}_{\mathrm{sc}}(d_B^2) = \kappaChHodge(\cA)\,\omega_g$
(Proposition~\ref{prop:curved-R-factorization}). Koszul-dualising
the curved $A_\infty$-algebra $(B(\cA), d_B, \omega_g)$ produces a
curved dg-category $\mathrm{CDG}(\cA, g)$ in the sense of
Positselski.

\medskip
\noindent\textbf{Periodicity candidate.} Let $u$ be the
degree-$(-2)$ generator of $H^\bullet(BS^1) = k[u]$. Negative
cyclic homology $\mathrm{HC}^-_\bullet(\cA)$ is a module over
$k[u]$, and
$\mathrm{HP}_\bullet(\cA) =
\mathrm{HC}^-_\bullet(\cA)[u^{-1}]$ is $2$-periodic by the
Connes--Tsygan $SBI$ sequence. The conjecture is that the
genus-$g$ bar, viewed through
Theorem~\ref{thm:hochschild-bridge-higher-genus}, extends the
$k[u]$-module structure on the genus-$0$ stalk to a GLOBAL
$k[u]$-action on the curved cdg-category, and that this action
becomes \emph{invertible} (yielding $u^{-1}$) after passing to
$\overline{\mathcal{M}}_g$-local sections.

\medskip
\noindent\textbf{Precise statement} (periodicity form). For
chirally Koszul $\cA$ at non-critical, non-admissible level:
\begin{enumerate}[label=\textup{(\roman*)}]
\item Genus~$1$: the bar $B(\cA)_1$ underlies
$\mathrm{HC}^-_\bullet(\cA)$ as a $k[u]$-module, and
$u$-inversion recovers $\mathrm{HP}_\bullet$. This is
Connes periodicity and is PROVED
(Theorem~\ref{thm:hochschild-bridge-genus0} plus the
mixed-complex $(b, B^{(0)})$ formalism).
\item Genus~$g \geq 2$: the curved twisting cochain
$\tau_g : B(\cA)_g \to \cA$ is $u$-compatible,
i.e.\ commutes with $S$ up to homotopy. This is CONJECTURAL;
the obstruction is an element of
$H^2(\overline{\mathcal{M}}_g,\, \mathrm{End}_{k[u]}(\cA))$
whose vanishing would follow once the curved-Dunn tower
$H^2 = 0$ bridge is tensored with the $k[u]$-action: a step
not yet carried out.
\item At admissible level $k = -h^\vee + p/q$: the conjecture
collapses to Koszul purity on \(O^{\mathrm{int}}\) as used in
derived Langlands. The semisimplified KL adjunction closes globally if
the \(k[u]\)-linear compatibility in (ii) holds.
\end{enumerate}

\medskip
\noindent\textbf{Terminology hazards}.
\begin{enumerate}[label=\textup{(\alph*)}]
\item \emph{Right}: the $u$-variable here IS the
negative-cyclic $u \in H^2(BS^1)$, and $u^{-1}$-inversion IS
the $\mathrm{HC}^- \to \mathrm{HP}$ periodicisation. The
conjecture is an assertion about the $k[u]$-module structure
on the curved bar at genus~$\geq 1$.
\item \emph{Wrong}: identifying this periodicity with Bott
periodicity of $KU$. Bott is a statement about complex
$K$-theory spectra; the $u$-periodicity here is an
$S^1$-action on cyclic complexes and has no direct
$KU$-content (the two agree only on the rational topological
point after
$\mathrm{HP}_\bullet \simeq K_\bullet \otimes \mathbb{Q}$ for
smooth proper $\cA$, which is ORTHOGONAL to the genus-$g$
chiral setting). Conflating them would import topological
$2$-periodicity as a free gift; it is not.
\item \emph{Correct}: ``$2$-periodic'' here means
$u^{-1}$-inverted $\mathrm{HC}^-$-module structure on
$B(\cA)_g$ along the $S^1$-rotation of
$\overline{\mathcal{M}}_{g,1}$-framings. The content is
$u$-compatibility of the curved twisting cochain $\tau_g$;
once (ii) is established, downstream uses
(modular-bootstrap $H^2 = 0$, admissible-level KL adjunction,
$\Theta_A$ curvature classes) inherit periodicity as
corollaries, not hypotheses.
\end{enumerate}

\medskip
\noindent\textbf{Genus-\(\geq2\) \(u\)-compatibility.} Genus~$\geq 2$ $u$-compatibility is
the remaining obstruction. It is equivalent to combining the
modular-bootstrap-to-curved-Dunn bridge map~$\Phi$
(Proposition~\ref{prop:modular-bootstrap-to-curved-dunn-bridge})
with the Getzler--Jones Gauss--Manin connection on
$\mathrm{HC}^-_\bullet$ over $\overline{\mathcal{M}}_g$; flatness
of the latter would upgrade $\Phi$ to a $k[u]$-linear morphism at
every genus. This reduces periodic-CDG to a \emph{flatness}
statement on Gauss--Manin, within reach of the irregular-singular
KZB framework (Jimbo--Miwa).

\medskip
\noindent Downstream uses that currently depend on this
conjecture: the derived-Langlands semisimplified KL equivalence at
admissible level, modular-bootstrap $H^2$ vanishing in the
$u^{-1}$-inverted form, and the admissible-level bar-cobar
adjunction purity clause
(Theorem~\ref{V1-thm:kl-bar-cobar-adjunction}).
Each downstream use must carry an explicit
``\textup{[conditional on Remark~\ref{rem:periodic-cdg-mechanism}]}''
tag until (ii) of the precise statement is proved.
\end{remark}

\begin{remark}[Periodic-CDG reduction: Bezrukavnikov, Arkhipov--Gaitsgory,
Positselski]
\label{rem:periodic-cdg-proof-path}
For
Remark~\textup{\ref{rem:periodic-cdg-mechanism}}\,(ii), at admissible
level $k = -h^\vee + p/q$ with $p, q$ coprime and $q \geq h^\vee$,
the reduction uses four mathematical inputs; the first three are
theorems in the cited literature.

\medskip
\noindent\textbf{Bezrukavnikov equivalence~\cite{Bezrukavnikov16}:
admissible Hecke $\leftrightarrow$ Steinberg).}
The affine Hecke category at level $k$ admissible is equivalent, as a
monoidal dg category, to the derived category of
$G^\vee \!\times\! \mathbb{G}_m$-equivariant coherent sheaves on the
Steinberg variety
$\mathrm{St}^\vee = \widetilde{\mathcal{N}}^\vee
\times_{\mathfrak{g}^\vee}\!\widetilde{\mathcal{N}}^\vee$
of the Langlands dual group. This equivalence sends the regular block
of $O^{\mathrm{int}}_k(\widehat{\mathfrak{g}})$ to the heart of the
standard $t$-structure on
$D^b\mathrm{Coh}^{G^\vee\!\times\!\mathbb{G}_m}(\mathrm{St}^\vee)$, and
$\mathbb{G}_m$-equivariance supplies a natural degree-$(-2)$ loop
generator $\ell$ acting by the second Chern class of the tautological
bundle on $\widetilde{\mathcal{N}}^\vee$.

\medskip
\noindent\textbf{Arkhipov--Gaitsgory twisted Satake~\cite{ArkhipovGaitsgory09}.}
At the admissible twist $\kappa = p/q$, twisted geometric Satake
promotes the Steinberg equivalence to an equivalence of braided monoidal categories
$\mathrm{KL}_k \simeq \mathrm{Rep}_q(G^\vee)$ with
$q = e^{\pi i / q}$ a primitive $q$-th root of unity, restricted to
the integrable block. The class $\ell$ on the Steinberg side is transported to the
$\mathbb{G}_m$-equivariant parameter of the $G^\vee$-quantum-group
side, which on the bar complex
$B^{\mathrm{ch}}(U_q(\mathfrak{g}^\vee))$ realises the Connes
generator $u \in H^2(BS^1)$ through the circle action on framed
$\overline{\mathcal{M}}_{0,3}$ governing the braiding.

\medskip
\noindent\textbf{Positselski filtered curved Koszul duality~\cite{Positselski11,Positselski18}.}
The filtered CDG$\leftrightarrow$filtered dg duality of Positselski,
applied to the adic filtration on
$U_q(\mathfrak{g}^\vee)_{\mathrm{adm}}$ induced by the quantum
parameter, yields a curved dg-category
$\mathrm{CDG}_q(\mathfrak{g}^\vee)$ whose curvature element is the
Casimir shift $c_q - c_\infty$. Imposing the $2$-periodicity
$u \in \mathrm{End}^{-2}(\mathrm{id})$ from the Satake side; that is,
inverting $u$ on $\mathrm{Hom}$-complexes converts this CDG pair
into a genuine $2$-periodic curved dg-category
$\mathrm{CDG}_q^{\mathrm{per}}(\mathfrak{g}^\vee)$ in Positselski's
sense. The resulting $\mathrm{CDG}_q^{\mathrm{per}}$ is Morita-equivalent
to the $u^{-1}$-inverted negative-cyclic completion of the Steinberg
dg-category.

\medskip
\noindent\textbf{Affine modular law \(q\leftrightarrow 1/q\).}
At admissible level $k = -h^\vee + p/q$ the Kac--Wakimoto modular
$S$-matrix satisfies
$S(q) \,S(1/q) = \mathrm{id}$ on the integrable block, which at the
Steinberg/Satake level is the Serre duality involution
$\mathbb{D} : \mathrm{St}^\vee \to \mathrm{St}^\vee$
shifted by the canonical class. Iterating this involution provides the
invertible degree-$(-2)$ endomorphism that realises
$u^{-1}$. This gives a \emph{conceptual} source of periodicity; it
does not, however, automatically supply chain-level invertibility of
$u$ on $\mathrm{Hom}$-complexes, because the Kac--Wakimoto $S$ acts on
characters and is defined at the level of the Grothendieck group.

\medskip
\noindent\textbf{Chain-level obstruction.}
The first three inputs compose into a categorical equivalence
\[
O^{\mathrm{int}}_k(\widehat{\mathfrak{g}})
\;\simeq\;
\mathrm{CDG}_q^{\mathrm{per}}(\mathfrak{g}^\vee)
\]
\emph{conditional on} $u$ acting invertibly on the
$\mathrm{Hom}$-complexes of the Steinberg dg-category at admissible
twist. The affine modular law supplies this invertibility in the Grothendieck group
but not at chain level; the obstruction is a class in
$H^2\bigl(\mathrm{St}^\vee,\,
\mathrm{End}^{\mathbb{G}_m}\!(\mathrm{id})\bigr)^{G^\vee}$, which
vanishes iff the Serre-duality involution $\mathbb{D}$ lifts to a
$k[u]$-linear autoequivalence of the Steinberg dg-category
compatible with the integrable $t$-structure. That lift is the
remaining obstruction: it is the $u$-invertibility of the
Connes operator on Hom-complexes at admissible level, matching
Remark~\textup{\ref{rem:periodic-cdg-mechanism}}\,(ii).

\medskip
\noindent\textbf{Serre-duality lift.} Once the
Serre-duality lift is established, the equivalence above upgrades the
admissible bar-cobar adjunction
(Theorem~\ref{V1-thm:kl-bar-cobar-adjunction}, purity clause) from
\emph{conditional} to \emph{unconditional}; semisimplified
Kazhdan--Lusztig at admissible level closes globally; and the
modular-bootstrap $H^2 = 0$ statement inherits $u^{-1}$-inversion
uniformly in genus via the Gauss--Manin-flatness route recorded in
Remark~\textup{\ref{rem:periodic-cdg-mechanism}}.
\end{remark}

\begin{remark}[Serre-duality \(u\)-invertibility obstruction:
Serre-duality $u$-invertibility]
\label{rem:serre-duality-u-invertibility-decomposition}
\index{periodic-CDG!Serre-duality decomposition}%
\index{Steinberg variety!Serre duality lift}%
The remaining obstruction of
Remark~\textup{\ref{rem:periodic-cdg-proof-path}}, namely $u$-invertibility
of the Serre-duality involution $\mathbb{D}$ compatibly with the
integrable $t$-structure on the Steinberg dg-category at admissible
twist, factorises into three conditions.

\medskip
\noindent\textbf{\(t\)-exactness of \(\mathbb{S}\).}
The Serre-duality functor
$\mathbb{S} \colon D^b(\mathrm{St}^\vee) \to D^b(\mathrm{St}^\vee)$
is a $t$-exact autoequivalence at admissible level with respect to
the integrable $t$-structure transported from
$O^{\mathrm{int}}_k(\widehat{\mathfrak{g}})$ by the Bezrukavnikov
equivalence. This is
\emph{known} for integrable modules by
Bezrukavnikov--Mirkovic~\cite{Bezrukavnikov16}, where the
exotic $t$-structure on coherent sheaves on
$\widetilde{\mathcal{N}}^\vee$ is shown to be Serre-self-dual on the
integrable block; what remains is a \emph{global} compatibility
argument extending the check from the integrable block to the full
derived category $D^b(\mathrm{St}^\vee)$ at admissible twist. The
argument reduces to $t$-exactness checking on a finite set of
standard modules (the tilting generators of
Bezrukavnikov~\cite{Bezrukavnikov16}); see also~\cite{BD04} for the
chiral-algebra ambient in which the check is natural.

\medskip
\noindent\textbf{Commutation with \(u\).}
The functor $\mathbb{S}$ commutes, up to shift, with the action of
$u$ generated by the affine Grassmannian's $\mathbb{G}_m$-action on
$\mathrm{St}^\vee$ (the loop rotation underlying the Connes
generator on the Steinberg side). The obstruction to strict commutation is a
cohomology class
\[
\alpha_{\mathbb{S},u} \;\in\;
H^2\!\bigl(\mathrm{Aut}(\mathrm{St}^\vee),\,k\bigr),
\]
which, at admissible level, reduces to a finite computation in the
cohomology of the affine Weyl group $\widetilde{W}$ acting on the
dual Cartan: a standard combinatorial object, directly accessible
via the geometric Satake framework of
Arkhipov--Gaitsgory~\cite{ArkhipovGaitsgory09}. The claim is
$\alpha_{\mathbb{S},u} = 0$ at admissible twist.

\medskip
\noindent\textbf{Periodic localisation.}
Assuming (a) and (b), the composite $\mathbb{S}\circ u = u\circ\mathbb{S}$
holds on $D^b(\mathrm{St}^\vee)$ up to a coherent shift; passing to
the periodic derived category
$D^{\mathrm{per}} = D^b(\mathrm{St}^\vee)[u,u^{-1}]$ in the sense of
Positselski~\cite{Positselski11,Positselski18} then produces the
desired $k[u,u^{-1}]$-linear autoequivalence. Condition (c) is
formal once (a) and (b) are in place: the periodic localisation
intertwines the $t$-exact $\mathbb{S}$ with the invertible $u$-action
by a standard Bousfield-localisation argument in the curved dg-module
category.

\medskip
\noindent\textbf{Vanishing criterion.} Conditions (a), (b), (c)
\emph{collectively} closes the $u$-invertibility obstruction of
Remark~\textup{\ref{rem:periodic-cdg-proof-path}}, and consequently
closes the periodic-CDG conjecture (Remark~\ref{rem:periodic-cdg-mechanism}). Each condition is concrete:
(a) is a finite $t$-exactness check on standard modules;
(b) is an affine-Weyl-group cohomology class whose vanishing is
amenable to direct computation; (c) is a formal consequence of
(a) and (b) via Positselski's periodic curved Koszul
duality~\cite{Positselski11,Positselski18}.
\end{remark}

\begin{proposition}[Explicit computation of $\alpha_{\mathbb{S},u}$
for $\widehat{\mathfrak{sl}}_2$ at admissible level]
\label{prop:alpha-su-sl2-admissible}
\index{periodic-CDG!alpha class sl2 admissible}%
\index{affine Weyl group!central extension!Serre-u commutator}%
\index{Steinberg variety!sl2 admissible}%
Fix $\mathfrak{g} = \mathfrak{sl}_2$ and admissible level
$k = -2 + 1/q$ with $q \in \mathbb{Z}_{\geq 2}$ coprime to the
denominator lattice ($\gcd(q,2) = 1$ or $q$ even; both cases will be
distinguished below). In the decomposition of
Remark~\textup{\ref{rem:serre-duality-u-invertibility-decomposition}},
condition~\textup{(b)} reduces at this level to an affine-Weyl-group
cohomology computation.

\medskip
\noindent\textbf{Identifications.}
\begin{enumerate}[label=\textup{(\roman*)}]
\item Geometric Satake
(Arkhipov--Gaitsgory~\cite{ArkhipovGaitsgory09}) identifies the
Langlands-dual Steinberg variety
$\mathrm{St}^{\vee}_{\widehat{\mathfrak{sl}}_2}$ at level $q$ with the
affine flag variety $\mathrm{Fl}_{\mathrm{PGL}_2}$ at level~$q$; its
automorphism group at the
Kac--Moody / loop-rotation level contains the affine Weyl group
\[
\widetilde{W} \;=\; \mathbb{Z}\rtimes\mathbb{Z}/2
\]
generated by the simple reflection $s_0$ and the translation $t$
(with $s_0\, t\, s_0 = t^{-1}$).
\item The cohomology of $\widetilde{W}$ with trivial coefficients in
a field $k$ of characteristic $\neq 2$ is
\[
H^\bullet(\widetilde{W},\,k) \;=\;
H^\bullet(\mathbb{Z},k)\,\otimes\,H^\bullet(\mathbb{Z}/2,k)^{\mathbb{Z}}
\;=\; k[\eta]/\eta^2\ \text{ in degrees }\le 2,
\]
so $H^2(\widetilde{W},k) = k$, one-dimensional, detected by the
central-extension class of the metaplectic double cover
$\widetilde{W}_{\mathrm{met}} \to \widetilde{W}$.
\item The Serre-duality functor $\mathbb{S}$ acts on
$H^2(\widetilde{W},k)$ via its induced action on central extensions;
by Bezrukavnikov--Mirkovic~\cite{Bezrukavnikov16} (exotic
$t$-structure Serre-self-duality on the integrable block)
$\mathbb{S}$ descends to the identity on
$H^2(\mathrm{Aut}(\mathrm{St}^{\vee}),k)$ up to the shift by the
dualising twist, and the obstruction class $\alpha_{\mathbb{S},u}$
is the image of $1\in k$ under this shift.
\end{enumerate}

\medskip
\noindent\textbf{Computation.}
The shift in (iii) is determined by the level-$q$ central extension
of $\widetilde{W}$. For $\mathfrak{sl}_2$ at level~$q$ the metaplectic
cocycle is the Kubota cocycle~\cite{Kubota69}
\[
c_q \colon\ \widetilde{W}\times\widetilde{W}\ \longrightarrow\ k,\qquad
c_q(t^a s_0^\epsilon,\,t^b s_0^{\epsilon'}) \;=\;
(-1)^{\epsilon\epsilon' q\cdot ab}.
\]
Under $\mathbb{S}$ (acting as Cartan involution on $\widetilde{W}$
with loop-rotation twist of parity~$q$) the class $[c_q]$ maps to
$[c_q]^{(-1)^{q+1}}$ in $H^2(\widetilde{W},k) = k$. Hence
\[
\alpha_{\mathbb{S},u} \;=\; \tfrac{1}{2}\bigl(1 + (-1)^q\bigr)\ \in\ k,
\qquad\text{equivalently}\qquad
\alpha_{\mathbb{S},u} =
\begin{cases}
0 & q\ \text{odd},\\
1 & q\ \text{even},
\end{cases}
\]
where the even-$q$ value represents the non-trivial metaplectic
class of $H^2(\widetilde{W},k) = k$ (conventionally labelled
``$1/2$'' in metaplectic coset representatives) and vanishing is
understood in $k$ of characteristic~$\neq 2$.

\medskip
\noindent\textbf{Consequence.}
\begin{enumerate}[label=\textup{(\alph*)}]
\item \emph{Odd $q$}: $\alpha_{\mathbb{S},u} = 0$, condition~(b) of
Remark~\textup{\ref{rem:serre-duality-u-invertibility-decomposition}}
vanishes; together with conditions~(a) and~(c) this closes
$u$-invertibility of $\mathbb{S}$ on
$D^b(\mathrm{St}^{\vee}_{\widehat{\mathfrak{sl}}_2})$. At admissible
levels $k = -2 + 1/q$ with odd $q$, the periodic-CDG conjecture (Remark~\ref{rem:periodic-cdg-mechanism})
holds for $\widehat{\mathfrak{sl}}_2$.
\item \emph{Even $q$}: $\alpha_{\mathbb{S},u} = 1/2\neq 0$ in
characteristic~$\neq 2$; condition~(b) does \emph{not} vanish
directly, and the $u$-invertibility of $\mathbb{S}$ fails strictly
on the integer lattice. In characteristic~$2$ the class collapses
to~$0$ and odd-$q$-style closure applies; in characteristic~$\neq 2$
the residual obstruction is detected by the metaplectic
$\mathbb{Z}/2$-cover and requires passage to a level-$2q$ refinement.
\end{enumerate}
Thus $\widehat{\mathfrak{sl}}_2$ admissible level
$k = -2 + 1/q$ closes the periodic-CDG conjecture (Remark~\ref{rem:periodic-cdg-mechanism}) for odd~$q$
unconditionally (via Bezrukavnikov--Mirkovic~\cite{Bezrukavnikov16}
and Arkhipov--Gaitsgory~\cite{ArkhipovGaitsgory09}) and leaves the
even-$q$ residue as the locus where the original
Remark~\textup{\ref{rem:periodic-cdg-mechanism}}
$u$-compatibility obstruction remains live; the residue is a
Kubota-cocycle $\mathbb{Z}/2$-class, strictly weaker than the
original obstruction and amenable to a separate metaplectic
refinement.
\end{proposition}

\begin{proposition}[Extension of $\alpha_{\mathbb{S},u}$ to
$\widehat{\mathfrak{sl}}_N$ at admissible level; Kubota tuple and
$N\geq 3$ non-vanishing]
\label{prop:alpha-su-slN-admissible}
\index{periodic-CDG!alpha class slN admissible}%
\index{affine Weyl group!extended!slN}%
\index{Kubota cocycle!tuple!slN}%
\index{metaplectic obstruction!non-vanishing N geq 3}%
Fix $\mathfrak{g} = \mathfrak{sl}_N$ with $N\geq 2$ and admissible
level $k = -N + p/q$ with $\gcd(p,q)=1$ and $q\geq N$ coprime to the
Coxeter number $h^{\vee} = N$ in the sense of
Kac--Wakimoto~\cite{KacWakimoto88}. The extended affine Weyl group
is
\[
\widetilde{W}_{\mathrm{aff}}(\mathfrak{sl}_N) \;=\;
Q^{\vee}\rtimes W,\qquad
Q^{\vee} \cong \mathbb{Z}^{N-1},\quad W = S_N,
\]
with $Q^{\vee}$ the coroot lattice and $W$ the finite Weyl group
acting by permutation of the coordinate basis on $Q^{\vee}$.

\medskip
\noindent\textbf{Cohomology and the Kubota tuple.}
\begin{enumerate}[label=\textup{(\roman*)}]
\item By Bezrukavnikov--Finkelberg--Mirkovic-type
computation~\cite{Bezrukavnikov16,FinkelbergTsymbaliuk19} applied to
the metaplectic cover at level $q$,
\[
H^2\bigl(\widetilde{W}_{\mathrm{aff}}(\mathfrak{sl}_N),\,k\bigr)
\;\cong\; k^{N(N-1)/2},
\]
the dimension equal to the number of positive roots
$\alpha_{ij} = \varepsilon_i - \varepsilon_j$, $1\leq i<j\leq N$.
\item The obstruction class $\alpha_{\mathbb{S},u}$ is the
Kubota-type tuple
\[
\alpha_{\mathbb{S},u} \;=\; \bigl(\alpha_{ij}\bigr)_{1\leq i<j\leq N}
\;\in\; k^{N(N-1)/2},
\]
with per-root coefficients (generalising the scalar
$\mathfrak{sl}_2$ computation of
Proposition~\ref{prop:alpha-su-sl2-admissible})
\[
\alpha_{ij} \;=\; \tfrac{1}{2}
\bigl(1 + (-1)^{q\,\langle \alpha_{ij},\rho\rangle}\bigr)
\;=\; \tfrac{1}{2}\bigl(1 + (-1)^{q(j-i)}\bigr),
\]
where $\rho = \tfrac12\sum_{\alpha>0}\alpha$ is the Weyl vector of
$\mathfrak{sl}_N$ and $\langle\alpha_{ij},\rho\rangle = j-i$.
The Kubota cocycle~\cite{Kubota69} on each $\mathfrak{sl}_2$-triple
subsystem $(e_{ij},h_{ij},f_{ij})\subset\mathfrak{sl}_N$ produces
the per-root factor; the tuple assembles via restriction to the
$N(N-1)/2$ commuting $\mathfrak{sl}_2$-subsystems.
\end{enumerate}

\medskip
\noindent\textbf{Vanishing analysis.}
The full class vanishes,
$\alpha_{\mathbb{S},u} = 0 \in k^{N(N-1)/2}$, iff
$q(j-i)$ is odd for every $1\leq i<j\leq N$.
\begin{enumerate}[label=\textup{(\alph*)}]
\item For $N=2$: the pair $(i,j)=(1,2)$ gives the single
condition $q$ odd, recovering
Proposition~\ref{prop:alpha-su-sl2-admissible}.
\item For $N\geq 3$: take $(i,j)=(1,2)$, forcing $q$ odd; then
take $(i,j)=(1,3)$, forcing $2q$ odd; \emph{impossible} for
any $q\in\mathbb{Z}$. Hence $\alpha_{\mathbb{S},u}\neq 0$ at
\emph{every} admissible level of $\widehat{\mathfrak{sl}}_N$,
$N\geq 3$, in characteristic~$\neq 2$.
\end{enumerate}

\medskip
\noindent\textbf{Metaplectic consequence.}
The Kubota-cocycle route to the periodic-CDG conjecture (Remark~\ref{rem:periodic-cdg-mechanism}) closes
\emph{only} at $N=2$ with odd-denominator admissible levels
$k=-2+1/q$. For $\widehat{\mathfrak{sl}}_N$ with $N\geq 3$ the
metaplectic obstruction tuple is non-vanishing at every admissible
level, so condition~(b) of
Remark~\ref{rem:serre-duality-u-invertibility-decomposition} cannot
be closed by metaplectic cohomology alone. The Langlands-dual
Steinberg identification via Arkhipov--Gaitsgory~\cite{ArkhipovGaitsgory09}
survives unchanged; Serre-duality $u$-compatibility does \emph{not} reduce to a
single scalar and carries $\binom{N}{2}-1$ residual classes (all
but the $(1,2)$-entry can fail independently); periodic closure
requires a replacement mechanism: either a higher-metaplectic
cover (level-$2q$ refinement extended across all simple roots,
obstructed by the $(1,3),\ldots,(1,N)$ entries) or an entirely
different route via Finkelberg--Tsymbaliuk quantum geometric
Langlands at Yangian level~\cite{FinkelbergTsymbaliuk19}.
\end{proposition}

\begin{remark}[$B^{(j)}$ hierarchy convention on the chiral side;
first-principles]
\label{rem:B-j-hierarchy-convention-chiral}
\index{Connes hierarchy!chiral convention}%
\index{negative cyclic!framing level}%
We fix once and for all the hierarchy level at which
each Connes-type operator is invoked on the chiral side.
\begin{enumerate}[label=\textup{(\alph*)}]
\item \emph{$B^{(0)}$ \textup{(}classical Connes\textup{)}:} the
mixed-complex operator of
Connes--Tsygan~\textup{\cite{Con85,LQ84}} acting on the cyclic bar
$(\ChirHoch^\bullet(\cA), b)$. It satisfies
$\{b, B^{(0)}\} = 0$ \emph{per~$k$} unconditionally (mixed-complex
axiom); this is the operator underlying
$\mathrm{HC}^-_\bullet$, $\mathrm{HP}_\bullet$, and the
$SBI$ sequence at genus~$0$. This is the operator used in
Remark~\ref{rem:periodic-cdg-mechanism}: the periodicity candidate
$u \in H^2(BS^1)$ and its $u^{-1}$-inversion act through the
$(b, B^{(0)})$ formalism.
\item \emph{$B^{(j)}$ for $j\geq 1$ \textup{(}Connes hierarchy /
$S^d$-framing lift\textup{)}:} the higher Connes operators arising
from $S^d$-framing cycle data on the cyclic bar\textup{;}
$j=2$ is the $S^2$-framing operator native to chiral $\mathrm{CY}_2$
geometry \textup{(}Vol~II operative dimension, since the closed
colour of $\SCchtop$ is $E_2$-chiral\textup{)}. Unlike~$B^{(0)}$,
individual $\{b_k, B^{(j)}\}\neq 0$ for non-formal
$\cA$\textup{;} only the total
$\{b, B^{(j)}\} = 0$ holds, and only as a Costello
TCFT-supertrace identity \textup{(}cross-arity Stasheff
cancellation\textup{)}; \emph{never} per-$k$. Extending a
per-$k$ identity from $j=0$ to $j\geq 1$ is a scope error
\textup{:} the higher framed operators see cross-arity cancellations.
\item \emph{Vol~II operative level.} In this volume the chiral
derived-center pair $(Z^{\mathrm{der}}_{\mathrm{ch}}(\cA),\cA)$
carries the $\SCchtop$-structure. The Connes
periodicity mechanism of Remark~\ref{rem:periodic-cdg-mechanism}
and the $\mathrm{HC}^-_\bullet$ invocation in
Proposition~\ref{prop:derived-center-E2} use $B^{(0)}$
\emph{exclusively}: both are genus-$0$ / mixed-complex statements.
At genus $g\geq 1$, where scalar curvature
$\operatorname{pr}_{\mathrm{sc}}(d_B^2) = \kappaChHodge(\cA)\,\omega_g$ enters, the operative framing is
$B^{(2)}$ via the $S^2$-framing of
$\overline{\mathcal{M}}_{g,1}$-strata; any global Lagrangian
complementarity statement over
$\overline{\mathcal{M}}_{g,n}$ \textup{(}Volume~I
Theorem~C; $\mathrm{HC}^-$-valued CY-trace pairing\textup{)} is
a $B^{(2)}$-statement, and the vanishing required for
Lagrangianness is the \emph{total} TCFT identity, not a sum of
per-$k$ terms.
\end{enumerate}
\emph{Consequence.} The Vol~II frontier~\#3 periodic-CDG mechanism is a
$B^{(0)}$-conjecture \textup{(}Connes classical\textup{)}; the Vol~II
global Lagrangian complementarity over
$\overline{\mathcal{M}}_{g,n}$ is a $B^{(2)}$-statement. Universal
Holography \textup{(}Vol~II climax\textup{)} couples both via the
DS--Hochschild compatibility bridge of
Theorem~\ref{thm:chd-ds-hochschild}: chain-level
$\mathrm{HC}^-$-framing on the bulk uses $B^{(0)}$
level-wise on each weight, while the genus tower uses $B^{(2)}$
only in aggregate via Costello TCFT supertrace. Any future
occurrence of ``Connes~$B$'' or ``$\mathrm{HC}^-$'' in Vol~II
chapters must cite this remark and declare the framing level.
\end{remark}

\begin{definition}[Filtered curved comparison]
\label{def:filtered-curved-comparison}
Fix \(g\geq1\) and write
\[
R_g=(\Omega^\bullet(\overline{\mathcal M}_g),d_{\mathrm{dR}},\omega_g)
\]
for the curved base with central degree-two curvature class
\(\omega_g\).  A filtered \(R_g\)-curved complex of scalar curvature
\(\kappa\omega_g\) is a filtered \(R_g\)-module \((C,F^\bullet)\) with
a degree-one operator \(d_C\) such that
\[
d_C^2=\kappa\,\omega_g\cdot \mathrm{id}_C .
\]
A \emph{filtered curved comparison}
\[
f\colon (C,F,d_C,\kappa\omega_g)\longrightarrow
(D,F,d_D,\kappa\omega_g)
\]
is a filtration-preserving \(R_g\)-linear map satisfying
\[
f d_C=d_D f,\qquad
f(\kappa\omega_g x)=\kappa\omega_g f(x),
\]
and whose special fibre at \(\omega_g=0\), equivalently the associated
graded obtained before the curvature twist, is a filtered
quasi-isomorphism of ordinary dg complexes.  When \(\kappa=0\), this is
just a filtered quasi-isomorphism of dg complexes.  When
\(\kappa\ne0\), the phrase ``quasi-isomorphism of curved complexes'' is
not used: the datum is the filtered curved comparison above.
\end{definition}

\begin{theorem}[Filtered curved Hochschild bridge at genus~$g \geq 1$;
\ClaimStatusProvedHere; licensing tags $\beta+\gamma$ via the
genus-$g$ curved base \(R_g\), factorisation descent, and the
curvature-compatible comparison]
\label{thm:hochschild-bridge-higher-genus}
Let\/ $\cA$ be a logarithmic\/ $\SCchtop$-algebra. Let $\cA$ be a chiral Koszul algebra
on a smooth projective curve $\Sigma_g$ of genus~$g \geq 1$.
Then there is a canonical filtered curved comparison over
\(\overline{\mathcal{M}}_g\):
\begin{equation}\label{eq:hochschild-bridge-genus-g}
\Psi_g\colon
\mathrm{C}^\bullet_{\mathrm{ch,top}}(\cA;\, \Sigma_g)
\;\longrightarrow\;
\ChirHoch^\bullet_g(\cA;\, \Sigma_g),
\end{equation}
where the left side is the genus-$g$ chiral Hochschild cochains
(the curved complex with $d^2 = \kappaChHodge(\cA) \cdot \omega_g$, as
in Proposition~\textup{\ref{prop:curved-R-factorization}}), and
the right side is Volume~I's chiral Hochschild descent object on
\(\Sigma_g\) equipped with the same scalar curvature class
\(\kappaChHodge(\cA)\omega_g\). The special fibre of \(\Psi_g\) at
\(\omega_g=0\) is a filtered quasi-isomorphism of ordinary dg
complexes. If \(\kappaChHodge(\cA)=0\), \(\Psi_g\) itself is an
ordinary filtered quasi-isomorphism.

The quasi-isomorphism has the following structure:
\begin{enumerate}[label=\textup{(\roman*)}]
\item \emph{Stalk identification.}
 On each fiber over a point $[\Sigma_g] \in \mathcal{M}_g$,
 the stalk of $\ChirHoch^\bullet(\cA;\, \Sigma_g)$ at any
 $z_0 \in \Sigma_g$ is quasi-isomorphic to
 $\mathrm{C}^\bullet_{\mathrm{ch,top}}(\cA)$ (the local
 complex of Theorem~\textup{\ref{thm:hochschild-bridge-genus0}}).

\item \emph{Global chiral descent.}
 The global complex on $\Sigma_g$ is the factorization descent
 totalization on the Ran space of $\Sigma_g$. After a symplectic
 marking of \(H_1(\Sigma_g,\Z)\), the transition functions are the
 Gauss--Manin local system together with the Hodge-line and prime-form
 factors of Volume~I's genus tower. A tensor factor
 \(\Lambda^\bullet H^1(\Sigma_g)\) appears only after trivialising this
 local system; it is not canonical on \(\mathcal M_g\).

\item \emph{Curvature.}
The curved differential on the left satisfies
$d^2 = \kappaChHodge(\cA) \cdot \omega_g$
(Proposition~\textup{\ref{prop:curved-R-factorization}}).
Under the identification, this curvature maps to the
 genus-$g$ cyclic obstruction class
 \(\Theta_g\in H^2_{\mathrm{cyc}}(\cA,\cA)\) of Volume~I
 (Theorem~D\textsubscript{scal}): the physical curvature IS the
 algebraic obstruction to extending the Hochschild complex
 from genus~$0$ to genus~$g$ as an uncurved dg object.

\item \emph{Complementarity.}
 For a Koszul pair $(\cA, \cA^!)$, the curvatures satisfy
 $\kappaChHodge(\cA) + \kappaChHodge(\cA^!) = K(\cA)$
 (Volume~I, Theorem~C), where $K(\cA)$ is a level-independent
 complementarity constant ($K = 0$ for the affine lineage;
 $K \neq 0$ for DS reductions). When $K = 0$, the tensor product
 $\ChirHoch^\bullet(\cA) \otimes \ChirHoch^\bullet(\cA^!)$
 is uncurved at every genus, and the combined Hochschild
 complex extends to a flat family over
 $\overline{\mathcal{M}}_g$. When $K \neq 0$, the combined
 complex retains residual curvature $K(\cA) \cdot \omega_g$.
\end{enumerate}
\end{theorem}

\begin{proof}
\textbf{Stalk identification.}
Fix $[\Sigma_g] \in \mathcal{M}_g$ and $z_0 \in \Sigma_g$.
Choose a holomorphic coordinate on a small disc \(D\ni z_0\).  This
identifies \(D\) with a small analytic disc in \(\C\), and its formal
completion with \(\operatorname{Spf}\C[[z]]\); no global
uniformisation of \(\Sigma_g\) is used.  On this formal disc the local
computation on \(D\times I\) is identical to the genus-$0$ case. By
Theorem~\ref{thm:hochschild-bridge-genus0} applied to the
restriction $\cA|_D$, the stalk of $\ChirHoch^\bullet(\cA)$ at
$z_0$ is quasi-isomorphic to the local complex
$\mathrm{C}^\bullet_{\mathrm{ch,top}}(\cA)$.

\textbf{Global assembly via chiral descent.}
Cover $\Sigma_g$ by formal discs and punctured formal
neighbourhoods of the pairwise diagonals. On each formal disc the
preceding stalk identification applies; on overlaps the transition
functions are the coordinate-change action on logarithmic kernels,
the Gauss--Manin connection on \(H^1(\Sigma_g)\), and the Hodge-line
factor in the prime-form normalisation. The global sections are
computed by the \v{C}ech--Thom--Sullivan totalization of the
factorization cosheaf
\cite[Chapter~6]{CG17}:
\begin{equation}\label{eq:cech-ts}
\Gamma(\Sigma_g,\, \ChirHoch^\bullet(\cA))
\;\simeq\;
\mathrm{Tot}\!\bigl(
 \check{C}^\bullet(\{U_\alpha\},\,
 \ChirHoch^\bullet(\cA))
\bigr).
\end{equation}
After a marking and a flat trivialisation of the Gauss--Manin local
system, the associated graded of this totalization contains the usual
\(\Lambda^\bullet H^1(\Sigma_g)\) factor. Without that trivialisation,
the descent object is the corresponding local system on
\(\mathcal M_g\), not a canonical exterior tensor product.
The local maps glue because the genus-$0$ bridge is natural under
coordinate change and because the same Gauss--Manin, Hodge-line, and
prime-form transition functions act on both sides.  This glued map is
\(\Psi_g\).  It preserves the arity/weight filtration, and on the
curvature-zero special fibre its associated graded is the stalkwise
genus-$0$ quasi-isomorphism glued by factorisation descent.

\textbf{Curvature identification.}
The curved bar differential on $\barB^{(g)}(\cA)$ satisfies
$d_{\mathrm{fib}}^2 = \kappaChHodge(\cA) \cdot \omega_g$
(Proposition~\ref{prop:curved-R-factorization}), where
$\kappaChHodge(\cA)$ is the modular characteristic (Volume~I,
Theorem~D) and $\omega_g$ is the Arakelov $(1,1)$-form.
Volume~I identifies $\kappaChHodge(\cA) \cdot \omega_g$ with the
genus-$g$ obstruction class
\(\Theta_g\in H^2_{\mathrm{cyc}}(\cA,\cA)\) via the Chern--Weil map
from the Hodge bundle $\mathbb{E} \to \overline{\mathcal{M}}_g$
(Theorem~D\textsubscript{scal}).
Under the stalk identification above, the physical curvature
(failure of $d^2 = 0$ on the fiberwise bar complex) maps to the
algebraic obstruction (failure of the Hochschild complex to
extend as a dg complex from genus~$0$ to genus~$g$).
Thus \(\Psi_g\) is a map in the filtered curved sense of
Definition~\ref{def:filtered-curved-comparison}; it is not an
ordinary quasi-isomorphism unless the curvature vanishes.

\textbf{Complementarity.}
For a Koszul pair $(\cA, \cA^!)$, Volume~I's Theorem~C gives
$\kappaChHodge(\cA) + \kappaChHodge(\cA^!) = K(\cA)$, where $K(\cA)$ is
the complementarity constant. The tensor product
$\ChirHoch^\bullet(\cA) \otimes \ChirHoch^\bullet(\cA^!)$
therefore has curvature
$(\kappaChHodge(\cA) + \kappaChHodge(\cA^!)) \cdot \omega_g
= K(\cA) \cdot \omega_g$.
When $K = 0$ (affine lineage: Kac--Moody, free fields, lattice
VOAs), the combined complex is a genuine dg complex (not curved)
at every genus, and the Hochschild bridge for the combined system
is a quasi-isomorphism of honest (uncurved) dg algebras
over $\overline{\mathcal{M}}_g$. When $\kappaChHodge(\cA) + \kappaChHodge(\cA^{!}) \neq 0$ (DS reductions:
$\kappa + \kappa^{!} = 13$ for $\mathrm{Vir}_c$,
$\kappa + \kappa^{!} = 250/3$ for $\mathcal{W}_3$, both nonzero
 projections of the complementarity conductor $K = c + c^{!}$ of Vol~I
\cref{V1-chap:universal-conductor} under the anomaly-ratio factor
$\varrho$), the combined complex retains a level-independent
residual curvature.
\end{proof}

\begin{proposition}[Moduli action on chiral Hochschild;
\ClaimStatusProvedHere]
\label{prop:hochschild-mbar-action}
Let $\cA$ be a chirally Koszul logarithmic $\SCchtop$-algebra on a
smooth projective curve. Then the chiral Hochschild complex
$\ChirHoch^\bullet(\cA,\cA)$ carries a natural action of the
rational homology of the Deligne--Mumford compactification:
\begin{equation}\label{eq:hochschild-mbar-action}
\rho_{\cA}\colon\; H_\bullet\bigl(\overline{\mathcal{M}}_{g,n};\,
\mathbb{Q}\bigr) \;\longrightarrow\;
\mathrm{End}_{\mathrm{Vect}_\mathbb{Q}^{\mathbb{Z}}}\!\bigl(
\ChirHoch^\bullet(\cA,\cA)^{\otimes n}\bigr),
\end{equation}
of total cohomological degree $-d$ on a class of homological
degree~$d$, via the chiral analogue of Costello's Hochschild--Segal
assignment. For each class $\gamma \in
H_d(\overline{\mathcal{M}}_{g,n};\mathbb{Q})$, the operation
$\rho_\cA(\gamma)$ integrates the chiral factorization density over
a cycle representative of $\gamma$:
\[
\rho_\cA(\gamma)(\alpha_1,\dots,\alpha_n)
\;=\; \int_{\gamma}\; \mu^{\mathrm{ch}}_{g,n}(\alpha_1,\dots,\alpha_n),
\]
where $\mu^{\mathrm{ch}}_{g,n}$ is the genus-$g$ $n$-point chiral
factorization density associated to $\cA$ over
$\overline{\mathcal{M}}_{g,n}$. The action has the following
properties.
\begin{enumerate}[label=\textup{(\roman*)}]
\item \emph{Operadic compatibility with higher Deligne.}
 The restriction of $\rho_\cA$ to the nodal boundary strata of
 $\overline{\mathcal{M}}_{g,n}$ recovers the brace operations of
 Theorem~\textup{\ref{thm:chd-deligne-tamarkin}}: on a codimension-one
 boundary stratum parameterised by a stable dual graph $\Gamma$ with
 one node, $\rho_\cA(\partial_\Gamma)$ is the brace insertion
 associated to that node.
\item \emph{Cup product from a point.}
 The fundamental class $[\mathrm{pt}] \in
 H_0(\overline{\mathcal{M}}_{0,3};\mathbb{Q})$ maps under $\rho_\cA$
 to the cup product on $\ChirHoch^\bullet(\cA,\cA)$, recovering the
 $E_2$-algebra multiplication of the chiral derived centre.
\item \emph{Curvature and the Hodge bundle.}
 The Chern class of the Hodge bundle
 $c_1(\mathbb{E}) \in H^2(\overline{\mathcal{M}}_{g,n};\mathbb{Q})$
 acts, via Poincar\'e duality, as multiplication by the curvature
 cocycle $\kappaChHodge(\cA)\cdot\omega_g$ of
 Proposition~\textup{\ref{prop:curved-R-factorization}}, intertwining
 the geometric $\overline{\mathcal{M}}_{g,n}$-action with the
 algebraic genus tower.
\end{enumerate}
\end{proposition}

\begin{proof}
\textbf{Cohomological construction.}
On the Koszul locus, Kontsevich formality (in the guise used for
Theorem~\textup{\ref{thm:chd-deligne-tamarkin}}) identifies
$\ChirHoch^\bullet(\cA,\cA)$ with an algebra over a cofibrant model
of the chiral $E_2$-operad. Costello's Hochschild--Segal construction
produces, for any algebra over such a model, a homological action of
$H_\bullet(\overline{\mathcal{M}}_{g,n};\mathbb{Q})$ via integration
of the associated modular descent datum on
$\overline{\mathcal{M}}_{g,n}$. Applied to the chiral factorization
density $\mu^{\mathrm{ch}}_{g,n}$
(Proposition~\textup{\ref{prop:curved-R-factorization}}), this gives
the map~\eqref{eq:hochschild-mbar-action} at the level of cohomology.

\textbf{Chain level on the chirally Koszul locus.}
Kontsevich formality further lifts the cohomological action to a
chain-level action on the chirally Koszul locus, because the formality
quasi-isomorphism is operadic: it intertwines the modular operadic
composition of Getzler--Kapranov with the brace operations on
$\ChirHoch^\bullet(\cA,\cA)$. The transfer of the homological
$\overline{\mathcal{M}}_{g,n}$-action through the formality zigzag
preserves operadic coherence by the standard homotopy transfer lemma.

\textbf{Boundary stratum compatibility.}
A codimension-one boundary stratum of $\overline{\mathcal{M}}_{g,n}$
is a clutching locus for a stable dual graph with one node.
Costello's construction identifies the integral over such a stratum
with the operadic composition along the node, which, under the
operadic formality quasi-isomorphism above, is precisely the brace operation
indexed by the node's combinatorial type. This yields~(i).

\textbf{Genus-zero three-point case.}
For $g=0$, $n=3$, the moduli space is a point. The fundamental class
maps to the operadic composition $\mu^{\mathrm{ch}}_{0,3}$, which on
$\ChirHoch^\bullet(\cA,\cA)$ is the Gerstenhaber cup product, proving~(ii).

\textbf{Hodge bundle and curvature.}
The Chern class $c_1(\mathbb{E})$ is represented by the Arakelov
$(1,1)$-form $\omega_g$ via Mumford's isomorphism
$c_1(\mathbb{E}) = [\omega_g]$ in $H^2(\overline{\mathcal{M}}_g;\mathbb{Q})$.
Poincar\'e duality converts cap product with $c_1(\mathbb{E})$ into
the cohomological operator associated to the class
$\omega_g$, which by
Proposition~\textup{\ref{prop:curved-R-factorization}} acts as
multiplication by $\kappaChHodge(\cA)\cdot\omega_g$.
\end{proof}

\begin{remark}[Scope and Costello's universal property]
\label{rem:hochschild-mbar-scope}
Part~(i)--(ii) of
Proposition~\textup{\ref{prop:hochschild-mbar-action}} hold at chain
level on the chirally Koszul locus, where Kontsevich formality lifts
the cohomological construction. Outside that locus, the
$\overline{\mathcal{M}}_{g,n}$-action is constructed cohomologically
only; a chain-level lift requires a Drinfeld associator and is
equivalent to the associator-dependent version of the chiral
higher Deligne conjecture
(Theorem~\textup{\ref{thm:chd-deligne-tamarkin}}). The universal
property of Costello's construction, namely that the
$\overline{\mathcal{M}}_{g,n}$-action is the unique natural
extension of the $E_2$-algebra structure compatible with the modular
operadic composition on the Deligne--Mumford boundary, is inherited
by the chiral case because the formality zigzag is natural in the
algebra argument.
\end{remark}

\begin{remark}[The full Hochschild bridge]
Theorems~\ref{thm:hochschild-bridge-genus0}
and~\ref{thm:hochschild-bridge-higher-genus} together constitute
the Hochschild bridge in its honest form.  At genus \(0\), the bridge
is an ordinary dg quasi-isomorphism between the local
holomorphic-topological cochain complex and the chiral Hochschild
complex.  At genus \(g\geq1\), the bridge is the filtered curved
comparison
\(\Psi_g\) of
Theorem~\ref{thm:hochschild-bridge-higher-genus}: the two sides carry
the same scalar curvature
\(\kappaChHodge(\cA)\omega_g\), the special fibre at
\(\omega_g=0\) is a filtered quasi-isomorphism of ordinary dg
complexes, and the comparison becomes an ordinary quasi-isomorphism
only when the curvature vanishes.  Thus the physical bulk observables
of a supplied 3d HT realization are identified with chiral Hochschild
cochains in genus \(0\), and in genus \(g\geq1\) only in this filtered
curved \(R_g\)-sense unless the \(K=0\) or
\(\kappaChHodge(\cA)=0\) hypothesis removes the curvature.

In this form the bridge supplies the Vol~II Hochschild counterpart to
the Vol~I algebraic theorem without erasing the genus-\(g\) obstruction
class:
\begin{center}
\renewcommand{\arraystretch}{1.2}
\begin{tabular}{lll}
\textbf{Theorem} & \textbf{Vol~I (algebra)} &
 \textbf{Vol~II (physics)} \\
\hline
A (bar-cobar) & Factorization coalgebra &
 Swiss-cheese algebra on $\C \times \R$ \\
B (inversion) & Quillen equivalence &
 On the chirally Koszul locus, line operators $\simeq$
 $\cA^!_{\mathrm{line}}$-modules \\
C (complementarity) & $\kappa + \kappa^! = \varrho_{\cA}\,K$ &
 Anomaly cancellation ($K = 0$ for affine lineage,
 $\varrho_{\mathrm{Vir}} = 1/2$ at $K = 26$,
 $\varrho_{\cW_N} = H_N - 1$) \\
D (leading coeff.) & $\kappaChHodge(\cA) \cdot \omega_g$ &
 Curved HT at genus $g$ \\
H (Hochschild) & chiral derived centre &
 algebraic closed sector $\simeq$ chiral Hochschild at genus \(0\);
 genus \(g\geq1\) through \(\Psi_g\), not an ordinary
 quasi-isomorphism unless the curvature vanishes
\end{tabular}
\end{center}
\end{remark}

\begin{corollary}[Bridge to Volume~I Theorem~H; \ClaimStatusProvedHere]
\label{cor:hochschild-bridge}
When specialized to a chiral algebra $\cA$ on a curve $X$ with
trivial topological direction ($\R = \mathrm{pt}$), the chiral
Hochschild complex of this volume recovers the chiral Hochschild
complex $\ChirHoch^*(\cA)$ of Volume~I\@.  In particular, let
$\mathcal W^k(\fg)$ carry a family support datum
$H_H(\mathcal W^k(\fg);S_{\mathcal W})$ with finite-dimensional
cohomology.  Then
\[
 \operatorname{Supp}\ChirHoch^\bullet(\mathcal W^k(\fg))
 \subseteq S_{\mathcal W},
 \qquad
 P_{\mathcal W^k(\fg)}(t)
 =\sum_{n\in S_{\mathcal W}}
 \dim H^n(K_{\mathcal W^k(\fg),S_{\mathcal W}})t^n.
\]
\end{corollary}

\begin{proof}
The trivial topological direction leaves the Beilinson--Drinfeld
chiral complex on the curve.

\medskip
\noindent\textbf{Stalk identification by degeneration.}
When $\R = \mathrm{pt}$, the $\SCchtop$ structure degenerates: the
open color (topological direction) contributes only a one-dimensional
$E_1$-algebra, so the Swiss-cheese cochain complex
$\mathrm{C}^\bullet_{\mathrm{ch,top}}(\cA)$ reduces to the purely
holomorphic factorization cochain complex on~$X$ alone. Concretely,
the configuration space factor $\Conf_m(\R)$ collapses to a point for
each $m$, and the brace operations (Definition~\ref{def:brace-ops})
reduce to classical Gerstenhaber braces on the closed-color
(holomorphic) sector. The resulting complex is the BD-chiral
Hochschild complex $\mathrm{C}^\bullet_{\mathrm{BD\text{-}ch}}(\cA)$.
On any small disc $D \subset X$, the local computation is identical
to the genus-$0$ case after choosing a holomorphic coordinate on \(D\);
the formal completion at the chosen point is
\(\operatorname{Spf}\C[[z]]\), and
Theorem~\ref{thm:hochschild-bridge-genus0} identifies the
stalk at any $z_0 \in X$ with the local complex
$\mathrm{C}^\bullet_{\mathrm{ch,top}}(\cA)$. This gives
\[
 \ChirHoch^\bullet(\cA)|_{z_0}
 \;\simeq\;
 \mathrm{C}^\bullet_{\mathrm{BD\text{-}ch}}(\cA)|_{z_0}
\]
as a local quasi-isomorphism of factorization algebras on~$X$.

\medskip
\noindent\textbf{Genus correction from Volume~I.}
When $X = \Sigma_g$ has genus $g \geq 1$, global sections differ
from the local formal-disc complex by the Gauss--Manin and Hodge-line
transition functions of the genus tower. Volume~I computes these
period-matrix corrections explicitly
(Theorem~D\textsubscript{scal}). After a marking and flat
trivialisation the associated graded contains the exterior factor
\(\Lambda^\bullet H^1(\Sigma_g)\); before such a choice the global
object is the corresponding chiral descent local system. In the
trivial-topological-direction specialization, the curvature
reduces to $d^2 = \kappaChHodge(\cA) \cdot \omega_g$ acting purely
on the closed color, which is Volume~I's genus-$g$ obstruction.

\medskip
\noindent\textbf{Global gluing.}
The factorization algebra \(\ChirHoch^\bullet(\cA)\) on \(X\) is
holomorphic and logarithmic, not locally constant. Its global sections
are computed by chiral \v{C}ech--Thom--Sullivan descent on formal discs
and punctured formal neighbourhoods of diagonals. Applying this descent
with the local complex above and the Volume~I transition functions
recovers \(\ChirHoch^\bullet(\cA)\) on \(X\).

\medskip
\noindent\textbf{Family support.}
For $\cA=\mathcal W^k(\fg)$, the datum
$H_H(\cA;S_{\mathcal W})$ supplies a complete chart model, a
deformation retract onto $K_{\cA,S_{\mathcal W}}$, and the
ordered-to-symmetric comparison.  Volume~I Theorem~H therefore gives
the displayed support and Hilbert polynomial.  Continuous
Gel'fand--Fuchs cohomology remains a separate functor of the formal
mode Lie algebra.
\end{proof}

\begin{remark}[HCA descent and the Hochschild--\v{C}ech comparison]
\label{rem:hca-hochschild-descent}
\index{homotopy chiral algebra!Hochschild descent}
The Hochschild bridge theorem above compares two complexes that
compute the same cohomology. At the chain level, the comparison
has additional algebraic structure: the \v{C}ech complex of the
presheaf $U \mapsto C^\bullet_{\mathrm{ch,top}}(\cA|_U)$ carries
a \emph{homotopy chiral algebra} structure $\{F_n\}$ by the
Malikov--Schechtman theorem~\cite{MS24} (Volume~I,
strict \v{C}ech HCA package). The bridge quasi-isomorphism
$\ChirHoch_{\mathrm{ch,top}} \xrightarrow{\sim}
\ChirHoch_{\mathrm{BD\text{-}ch}}$
of Theorem~\ref{thm:hochschild-bridge-genus0} intertwines
this HCA structure with the $A_\infty$ structure on the bar-side
Hochschild complex, as in the cosimplicial Hochschild discussion
of Volume~I.

For the Swiss-cheese cochain complex
$\mathrm{C}^\bullet_{\mathrm{ch,top}}(\cA)$, the \v{C}ech
direction indexes restrictions to open subsets of the boundary
curve, while the Hochschild direction indexes marked points
on configuration spaces. The HCA operations $F_n$ then control
the compatibility between these two directions: $F_2$ is the
brace bracket, $F_3$ is the Jacobiator homotopy constructed in
Volume~I,
and higher $F_n$ vanish on cohomology by
genus-$0$ contractibility.

In the language of Maurer--Cartan elements: the \v{C}ech
convolution algebra
$\check{C}^\bullet(\mathcal{U};\,
\mathrm{C}^\bullet_{\mathrm{ch,top}}(\cA))$ carries a
total differential $D = d_M + d_{\mathrm{Hoch}}$, where
$d_M$ is the \v{C}ech differential and
$d_{\mathrm{Hoch}}$ is the Hochschild differential
on fibers. The MC element $\Phi_{\mathrm{ChirHoch}}$
encoding the Swiss-cheese structure satisfies
$D\,\Phi_{\mathrm{ChirHoch}}
+ \tfrac{1}{2}[\Phi_{\mathrm{ChirHoch}},
\Phi_{\mathrm{ChirHoch}}] = 0$, and
the HCA operations $F_n$ control precisely
the compatibility between $d_M$ and $d_{\mathrm{Hoch}}$:
$F_2$ gives the leading intertwining term, while
$F_n$ ($n \geq 3$) supply the higher homotopy corrections
that ensure $D^2 = 0$ at the total complex level.

Given $H_H(\mathcal W^k(\fg);S_{\mathcal W})$, Volume~I Theorem~H
identifies the chart cohomology with
$H^\bullet(K_{\mathcal W^k(\fg),S_{\mathcal W}})$ and places its
support in~$S_{\mathcal W}$.  The bridge transfers the HCA operations
to this family model.  Their possible cohomological targets are the
degrees in~$S_{\mathcal W}$; an explicit transferred brace calculation
determines which higher operations vanish.
\end{remark}

\subsection{Closed sector from boundary: chiral Hochschild computations}
\label{subsec:bulk-from-boundary}

The algebraic object computed below is the chiral Hochschild cochain
complex of the boundary algebra with its $E_2$-structure
\textup{(}Proposition~\ref{prop:swiss-cheese-from-tamarkin}\textup{)}.
When a standard family is realized by a chosen 3d HT
prefactorization model satisfying Theorem~\ref{thm:physics-bridge}
and the boundary-linear exact-sector hypotheses, Theorem~\ref{thm:bulk_hochschild}
identifies that complex with the corresponding physical bulk
observables. Thus each example computes $C^*(\cA,\cA)$ first:
algebraic closed sector $\simeq$ chiral Hochschild; physical bulk only
after the scoped comparison.

\begin{computation}[Chiral Hochschild closed sector for Heisenberg;
\ClaimStatusProvedElsewhere for the bounded profile; chart reading
conditional on the bounded-to-chart comparison]
\label{comp:bulk-heisenberg}
\index{Heisenberg algebra!Hochschild closed sector}
$\cA=\cH_k$ at nonzero level, identified with the rank-one even
superboson $B_{\mathfrak h}$ of
Proposition~\ref{prop:vol2-superboson-bounded-benchmark}. The
bounded cohomology has the Bakalov--De Sole--Kac profile
\[
 \dim H^n_{\mathrm{ch},b}(\cH_k,\cH_k)=(2,1,0,0,\ldots),
 \qquad
 \text{support }\{0,1\},\qquad
 \text{Hilbert polynomial } 2+t
\]
\textup{(}Bakalov--De Sole--Kac \cite{BDSK21}, Theorem~7.4:
$H^n\cong(S^n(\Pi\mathfrak h))^*\oplus(S^{n+1}(\Pi\mathfrak h))^*$
with $\Pi\mathfrak h$ odd one-dimensional\textup{)}. In particular
$H^1\neq0$: the weight-lowering derivation
\[
 D\colon \alpha\longmapsto|0\rangle,
 \qquad
 D(a_{(n)}b)=(Da)_{(n)}b+a_{(n)}(Db),
\]
is \emph{outer}. An inner derivation is a zero mode $a_{(0)}$ for
some $a\in\cH_k$, and no zero mode realizes $D$:
$\alpha_{(0)}\alpha=0$ by abelianness of the $0$-th product, and
$({:}\alpha^m{:})_{(0)}\alpha=mk\,\partial\,{:}\alpha^{m-1}{:}$
(with the convention ${:}\alpha^0{:}=|0\rangle$,
$\partial|0\rangle=0$) is never the vacuum; derivative generators
only add total derivatives. Hence $[D]\neq0$ spans $H^1$.
$H^2=0$: the level variation $k\to k+\varepsilon$ is exact for
$k\neq0$, absorbed by rescaling the Heisenberg generator, in
agreement with the vanishing bounded $H^2$; the level is a
normalization coordinate of the family, not a fixed-fibre
deformation class. The Gerstenhaber bracket on the surviving
classes is trivial: $[D,D]_G=2D^2=0$ since
$D^2\alpha=D|0\rangle=0$. When the bounded-to-chart comparison
$\chi^{\mathrm{bd}}_{B_{\mathfrak h}}$ of
Proposition~\ref{prop:vol2-superboson-bounded-benchmark} is a
quasi-isomorphism, the chart closed sector carries the same profile
$(2,1,0,\ldots)$. The fixed-fibre \emph{reduced} presentation of
Computation~\ref{comp:tamarkin-e2-heisenberg} is a different
complex (reduced at the vacuum, fixed fibre); the comparison
between the two normalizations is the grading obligation of
Computation~\ref{comp:drinfeld-center-heisenberg}.
\end{computation}

\begin{computation}[Chiral Hochschild closed sector for affine $\widehat{\mathfrak{sl}}_2$;
\ClaimStatusProvedHere]
\label{comp:bulk-sl2}
\index{Kac--Moody!Hochschild closed sector}
$\cA=\widehat{\mathfrak{sl}}_2{}_k$. The three currents
$e,f,h$ generate $\cH^0=\Bbbk$.
$H^1(C^*(\widehat{\mathfrak{sl}}_2,
\widehat{\mathfrak{sl}}_2))=0$: the adjoint maps
$\operatorname{ad}_{e_0}$,
$\operatorname{ad}_{f_0}$, $\operatorname{ad}_{h_0}$
are inner derivations (hence coboundaries),
and $\widehat{\mathfrak{sl}}_2{}_k$ has no outer
derivations at generic level.
$H^2=\Bbbk\langle\Theta_k\rangle$
(one-dimensional, generated by the $2$-cocycle recording
the Sugawara variation $\partial_k T_{\mathrm{Sug}}$
under $k\to k+\varepsilon$; unlike the Heisenberg level coordinate,
the affine simple-pole term prevents rescaling from absorbing it).

The Hochschild closed-sector algebra is
$V_T\cong U(\mathfrak{sl}_2)^{\mathfrak{sl}_2}
\;\cong\;\Bbbk[[\Omega]]$
(the completed ring of Casimir invariants).
The Gerstenhaber bracket:
$[\Omega,\Omega]_G=0$ (Casimir is central).
The closed-sector-to-boundary comparison map $\beta(\Omega)
=\sum_a J^a_{(-1)}J^a$ is the Sugawara
construction.

At the critical level $k=-2$: the Sugawara construction
degenerates, and $V_T$ acquires the Feigin--Frenkel center
$\mathfrak z(\widehat{\mathfrak{sl}}_2{}_{-2})
\cong\Bbbk[[\partial^n\mathfrak p]]$
(differential polynomials in the oper coordinate~$\mathfrak p$).
The closed-sector algebra jumps from $\Bbbk[[\Omega]]$ (generic~$k$)
to $\Bbbk[[\mathfrak p,\partial\mathfrak p,\dots]]$ ($k=-2$).
\end{computation}

\begin{computation}[Chiral Hochschild closed sector for Virasoro;
\ClaimStatusConditional]
\label{comp:bulk-virasoro}
\index{Virasoro algebra!Hochschild closed sector}
For $\cA=\mathrm{Vir}_c$, Bakalov--De Sole--Kac
\cite{BDSK21} \textup{(}Theorem~7.2\textup{)} compute the bounded cohomology
polynomial
\[
 P^{\mathrm{bd}}_{\mathrm{Vir}_c}(t)=1+t^2+t^3.
\]
The degree-two class contains the central-charge cocycle
\[
 \Theta_c(T,T)=\partial_c\{T_\lambda T\}=\lambda^3/12.
\]
If $\chi_c^{\mathrm{bd}}$ is a quasi-isomorphism, the chart closed
sector has the same polynomial and support~$\{0,2,3\}$.  A physical
bulk realization additionally uses the open--closed comparison map,
and the $E_2$ bracket is transported through its brace-compatible
form.
\end{computation}

\begin{computation}[Chiral Hochschild closed sector for $\mathcal W_3$;
\ClaimStatusProvedHere]
\label{comp:bulk-w3}
\index{W3@$\mathcal W_3$!Hochschild closed sector}
$\cA=\mathcal W_3{}_c$. $H^0=\Bbbk$.
$H^1=0$ (all derivations of $\mathcal W_3$ are inner
for generic~$c$).
$H^2=\Bbbk\langle\Theta_c^T,\,\Theta_c^W\rangle$
(two-dimensional: the $2$-cocycles recording
central-charge variation for each generator).
The Hochschild closed-sector algebra is
$V_T\cong\Bbbk[[c]]\otimes
\operatorname{Sym}^*\!(L_n^*\oplus W_m^*)$
with Poisson bracket from the $\mathcal W_3$ algebra.
The quartic self-interaction ${:}TT{:}$ in the $WW$-OPE
produces a nonlinear Poisson bracket:
$\{W_m,W_n\}_G$ involves terms quadratic in $L_p$
(the $\Lambda$-field contribution).

The two-dimensional $H^2$ produces the $2\times 2$
genus-$1$ Hessian
$H^{(1)}=\mathrm{diag}(\kappa_c,\,\kappa_c+48/(5c+22))$,
computed in Volume~I:
the eigenvalues are the curvatures of the two independent
closed-sector deformation directions.
\end{computation}

\begin{computation}[Chiral Hochschild closed sector for $\mathcal W_4$;
\ClaimStatusProvedHere]
\label{comp:bulk-w4}
\index{W4@$\mathcal W_4$!Hochschild closed sector}
$\cA=\mathcal W_{4,c}$. $H^0=\Bbbk$.
$H^1=0$ (all derivations of $\mathcal W_4$ are inner
for generic~$c$; Zamolodchikov--Fateev \cite{FL88}).
$H^2=\Bbbk\langle\Theta_c^T,\,\Theta_c^{W_3},\,
\Theta_c^{W_4}\rangle$ (three-dimensional: the $2$-cocycles
recording central-charge variation for each generator
of spin $s=2,3,4$).
The Hochschild closed-sector algebra is
$V_T\cong\Bbbk[[c]]\otimes
\operatorname{Sym}^*\!(L_n^*\oplus W^{(3)}_m{}^*
\oplus W^{(4)}_p{}^*)$
with Poisson bracket from the $\mathcal W_4$ algebra.

The three-dimensional $H^2$ produces the $3\times 3$
genus-$1$ Hessian $H^{(1)}_4$
(Computation~\ref{comp:bulk-wn}): its eigenvalues are
the curvatures of the three independent closed-sector deformation
directions. The complementarity constant is
$\alpha_4 = 246$ (Proposition~\ref{prop:wn-complementarity}),
with self-dual point $c = 123$.

\emph{Verification of~$H^1=0$ and~$H^2=\Bbbk^3$.}
$\mathcal W_4$ is simple for generic~$c$
(Zamolodchikov--Fateev \cite{FL88}). Every derivation
$D\colon \mathcal W_4 \to \mathcal W_4$ preserving
conformal weight is inner, so $H^1=0$. The maps
$\Theta_c^T$, $\Theta_c^{W_3}$, $\Theta_c^{W_4}$
are $2$-cochains recording
the infinitesimal shifts of the OPE structure constants under
$c \to c + \varepsilon$. These are non-trivial $2$-cocycles
(they shift structure constants without being coboundaries),
giving $H^2 = \Bbbk^3$.
\end{computation}

\subsection{The bulk--boundary factorization theorem}
\label{subsec:bulk-boundary-factorization-theorem}

The following theorem assembles the preceding results. The affine and
boundary-linear inputs are proved; the nonlinear KZ boundary vertex
algebra input is conditional on the analytic SDR package:
\begin{itemize}
\item physical bulk $\simeq$ chiral Hochschild cochains only in the
 scoped HT prefactorization and boundary-linear exact-sector comparison
 (Theorem~\ref{thm:bulk_hochschild});
\item boundary factorization algebra is a vertex algebra in the affine
 lane, and in the nonlinear finite-type KZ lane under
 Assumption~\ref{assum:kz-analytic-sdr-package}
 (Theorem~\ref{thm:general-half-space-bv});
\item the boundary-linear exact-sector comparison identifies physical
 bulk with the Morita-invariant derived centre of the boundary
 category
 (Theorem~\ref{thm:CDG_compatibility}\,(iii));
\item line operators $\simeq$ $\cA^!_{\mathrm{line}}$-modules
 on the chirally Koszul locus
 (Theorem~\ref{thm:lines_as_modules}).
\end{itemize}
The package below collects these into a single bulk--boundary--line datum.

\begin{theorem}[Bulk--boundary--line factorization correspondence]%
\label{thm:bulk-boundary-line-factorization}
\ClaimStatusConditional.
Let $\mathsf{Obs}$ be an HT prefactorization realization in the
scope of Theorem~\ref{thm:physics-bridge}, and $\cA_\partial$ its
boundary chiral algebra. Assume its
product-formal local shadow is a logarithmic $\SCchtop$-algebra
\textup{(}Definition~\textup{\ref{def:log-SC-algebra})}.
Then:
\begin{enumerate}[label=\textup{(\roman*)}]
\item \textbf{Bulk factorization algebra.}
The BV observable prefactorization algebra
$\mathsf{Obs}(\C \times \R)$ is a factorization algebra on
$\C \times \R$ whose product-formal local shadow is governed by
the Swiss-cheese operad $\SCchtop$.
Its cohomology $H^\bullet(\mathsf{Obs}, Q)$ is a commutative
Poisson vertex algebra with $(-1)$-shifted $\lambda$-bracket.

\item \textbf{Boundary vertex algebra.}
The boundary restriction $\cA_\partial$ is a vertex algebra in the
verified boundary sectors cited below.

\item \textbf{Bulk--centre comparison.}
There is a quasi-isomorphism of dg algebras
\[
\mathsf{Obs}(\C \times \R)
\;\simeq\;
C^\bullet_{\mathrm{ch}}(\cA_\partial, \cA_\partial),
\]
for HT prefactorization realizations in the scope of
Theorem~\ref{thm:bulk_hochschild}. If, in addition, the
compact-generation and derived-center hypotheses of
Theorem~\ref{thm:boundary-linear-bulk-boundary} hold
\textup{(}in particular, in the boundary-linear exact sector\textup{)},
then on cohomology
$H^\bullet(\mathsf{Obs}) \simeq Z_{\mathrm{der}}(\cA_\partial)$,
the derived center. The first equivalence is the
\(\Psi_{\mathrm{Ran}}\) comparison, or equivalently the inverse
\(\chi_{\mathrm{HT},A_\partial}\); it is not a physical bulk
construction from \(\cA_\partial\) alone.

\item \textbf{Line operators.}
If\/ $\cA_\partial$ is chirally Koszul, then the category of perturbative line operators
is equivalent to modules for the open-colour Koszul dual:
the category of line operators
$\cC_{\mathrm{line}} \simeq \cA_{\partial,\mathrm{line}}^!\text{-}\mathrm{mod}$,
where $\cA_{\partial,\mathrm{line}}^! =
H^*(\barB^{\mathrm{ord}}(\cA_\partial))^\vee$ is the open-colour
Koszul dual \textup{(}obtained by Verdier duality on the ordered
bar, not by cobar; the ordinary cobar $\Omega(\barB(\cA_\partial))$
recovers the boundary algebra itself\textup{)}.

\item \textbf{Spectral $R$-matrix.}
The braiding of perturbative line operators is encoded by the spectral
$R$-matrix $r(z)$, computed by the bulk-to-boundary
propagator exchange:
\[
r(z) \;\simeq\; \int_0^\infty d\lambda\, e^{-\lambda z}\,
\{{\cdot}{}_\lambda{\cdot}\},
\]
where $\{{\cdot}{}_\lambda{\cdot}\}$ is the $(-1)$-shifted
PVA $\lambda$-bracket on $H^\bullet(\mathsf{Obs})$.
\end{enumerate}
\end{theorem}

\begin{proof}
Each part is proved by a specific theorem.

\emph{Part (i).}
The observable prefactorization algebra of a perturbative QFT on a
manifold is a factorization algebra by the Costello--Gwilliam
recognition theorem~\cite{CG17, CG21}. That the cohomology is a
commutative PVA with $(-1)$-shifted bracket is
Theorem~\ref{thm:cohomology_PVA}. The local Swiss-cheese shadow
(holomorphic factorization in $\C$ and $E_1$ factorization in $\R$)
follows from the local-shadow theorem
(Theorem~\ref{thm:recognition-SC}): the weight forms factor as
$\omega_k^{\mathrm{hol}} \otimes \omega_k^{\mathrm{top}}$
(Definition~\ref{def:log-SC-algebra}), and this product decomposition
is exactly the input for the product-formal Swiss-cheese shadow.

\emph{Part (ii).}
For the affine family, the boundary vertex algebra is verified by
Proposition~\ref{prop:affine-is-log-SC} and
Theorem~\ref{thm:affine-half-space-bv}. For finite-type nonlinear KZ
PVAs with inner holomorphic translation,
Theorem~\ref{thm:general-half-space-bv} gives the boundary
factorization algebra as a vertex algebra only under
Assumption~\ref{assum:kz-analytic-sdr-package}. For Virasoro and
$W_N$, Corollary~\ref{cor:virasoro-wn-are-SC} proves the inherited
logarithmic $\SCchtop$ shadow; the all-loop half-space quantization is
not a formal consequence of that shadow alone.

\emph{Part (iii).}
This is Theorem~\ref{thm:bulk_hochschild} (quasi-isomorphism of the
bulk observable complex with chiral Hochschild cochains) combined with
Theorem~\ref{thm:CDG_compatibility}\,(iii) (the map to the center) in
the boundary-linear exact sector.

\emph{Part (iv).}
This is Theorem~\ref{thm:lines_as_modules}, whose hypothesis is exactly
the chirally Koszul condition stated in the theorem.

\emph{Part (v).}
The spectral $R$-matrix formula is
Theorem~\ref{thm:two-color-master}\,(iii): the braiding is computed by
a single bulk-to-boundary propagator exchange on the half-space
$\C \times \R_{\ge 0}$, and the Laplace transform formula follows from
the exponential decay of the propagator in the topological direction.
\end{proof}

\begin{remark}[Package content and scope]
Each ingredient was proved independently. The theorem is a single
reference point for a logarithmic $\SCchtop$-algebra after a chosen
HT prefactorization realization supplies the product-formal local
shadow and the boundary-linear exact-sector hypotheses of
Part~\textup{(iii)}. Under that scope, the bulk observable complex,
boundary algebra, and spectral braiding fit into one datum
\[
(\mathsf{Obs},\; \cA_\partial,\;
Z_{\mathrm{der}}(\cA_\partial),\;
C_{\mathrm{line}},\; r(z))
\]
while the derived-center and line-module identifications retain the
scope of \textup{(iii)}--\textup{(iv)}. An abstract logarithmic
$\SCchtop$-algebra without such a realization has only the algebraic
local-shadow package; it does not determine a physical bulk
observable complex. This is the boundary-side holographic datum.
The open question is bulk reconstruction: given the boundary datum,
recover the bulk theory and prove uniqueness.
\end{remark}

\begin{remark}[Beyond the package]
Two claims remain open.
\begin{enumerate}[leftmargin=2.4em]
\item \textbf{Bulk reconstruction.} Given the boundary vertex
 algebra $\cA_\partial$, reconstructing the unique
 bulk factorization algebra $\mathsf{Obs}$ with boundary restriction
 $\cA_\partial$ is open.
 Theorem~\ref{thm:bulk_hochschild} goes in the opposite direction:
 for a chosen physical prefactorization model, its bulk observables are
 compared with boundary chiral Hochschild cochains by
 \(\Psi_{\mathrm{Ran}}\) and \(\chi_{\mathrm{HT},A_\partial}\), with
 no uniqueness statement.
\item \textbf{Large-$N$ identification.} The genus expansion
 of the modular bar complex (Theorem~\ref{thm:modular-bar})
 produces a stable-graph sum with $1/|\Aut(\Gamma)|$ symmetry
 factors and genus parameter $\hbar$. No theorem identifies $\hbar$ with a large-$N$ parameter.
 The formal resemblance to the 't Hooft expansion is structural,
 not proved.
\end{enumerate}
\end{remark}

\begin{remark}[Bi-based Ran-space factorisation on the
bulk--boundary datum at the $K3\times E$ specialisation;
\ClaimStatusConditional]
\label{rem:hoch-bibased-Ran}
\index{Ran space!bi-based factorisation}%
\index{K3!bulk--boundary datum}%
The bulk--boundary--line datum
$(\mathsf{Obs},\,\cA_\partial,\,Z_{\mathrm{der}}(\cA_\partial),\,
\cC_{\mathrm{line}},\,r(z))$ of
Theorem~\ref{thm:bulk-boundary-line-factorization} is assembled
over a single Ran-space base $\Ran(C)$. At the Vol~III
$\Phi$-specialisation to the product Calabi--Yau input $K3\times E$
(with $E$ an elliptic curve and the K3 contributing the Mukai datum
of Remark~\ref{rem:hoch-CY2-K3-shift}), the datum upgrades to a
\emph{bi-based} Ran-space factorisation:
\[
\Ran^{\mathrm{bi}}(K3\times E)
\;=\;
\underbrace{\Ran\!\bigl(E^{\mathrm{nod,sm}}_{24}\bigr)}_{\text{chiral
base}}
\;\times\;
\underbrace{\overline{A_2}}_{\text{parameter base}},
\]
where $E^{\mathrm{nod,sm}}_{24}$ is the universal elliptic curve with
$24$ marked nodal-smooth degenerations (the Mathieu-equivariant
$24$-puncture configuration; nearby-cycles at each node
contribute one monodromy stratum, total $24$) and $\overline{A_2}$ is
the compactified polarisation-lattice parameter base. On this
bi-based site, the SC$^{\mathrm{ch,top}}$-structure on the chiral
derived-centre pair
$(C^\bullet_{\mathrm{ch}}(\cA_\partial,\cA_\partial),\,\cA_\partial)$
(Proposition~\ref{prop:derived-center-E2} and
Theorem~\ref{thm:chiral-higher-deligne-hoch-echo}) specialises at
$\cA_\partial=\cA_{K3}\otimes\cA_E$ to the pair
\[
\bigl(\Ran(E^{\mathrm{nod,sm}}_{24}),\,\overline{A_2}\bigr),
\]
with the $24$ nearby-cycle components indexing the chiral-base
factorisation and the $\overline{A_2}$-factor indexing a compatible
parameter coproduct. Compatibility of the two coproducts is the
the \emph{bi-base factorisation identity}
$(\Delta_{\mathrm{ch}}\otimes\id)\Delta_{A_2}
 =(\id\otimes\Delta_{A_2})\Delta_{\mathrm{ch}}$,
the chiral-side expression of the Vol~III $\Phi$-bridge splitting
$K3\times E \mapsto (K3, E)$
(compare Remark~\ref{rem:bar-cobar-4-bibased-ran} of the bar-cobar
review for the upstream form of the same identity on the bar
coalgebra). This specialisation refines the ``bulk-reconstruction''
open question of the preceding remark: the bi-based factorisation is
the shape the bulk acquires once the K3 input is switched on.
\end{remark}

\begin{theorem}[Global triangle for boundary-linear families; \ClaimStatusProvedHere]\label{thm:global-triangle-boundary-linear}%
\index{global triangle!boundary-linear families|textbf}%
\index{boundary-linear Landau--Ginzburg!global triangle closure}%
Let\/ $\cA$ be a chirally Koszul boundary-linear logarithmic\/
$\SCchtop$-algebra on the Gaussian or Lie comparison surface, or on a
contact surface equipped with an explicit nullhomotopy of its quartic
contact class. Thus the theorem applies to Heisenberg, lattice, and
generic affine\/ $V_k(\mathfrak{g})$ models, and to contact models only
after the quartic nullhomotopy has been chosen.
Then the global triangle closes: there are canonical
quasi-isomorphisms
\[
\Abulk
\;\simeq\;
\Zder(\Bbound)
\;\simeq\;
\Zder(\cC_{\mathrm{line}})
\;\simeq\;
\HH^\bullet_{\mathrm{ch}}(\cA^!_{\mathrm{line}}).
\]
\end{theorem}

\begin{proof}
The triangle is the composition of three equivalences.

\emph{Bulk as derived center of boundary.}
By hypothesis, $\cA$ is boundary-linear, so
Theorem~\ref{thm:boundary-linear-bulk-boundary} applies:
the derived center of the boundary algebra recovers the bulk,
$\Abulk \simeq \Zder(\Bbound)$.
The key input is that the boundary superpotential is linear in
the normal directions, so the derived HKR theorem identifies
Hochschild cochains of~$\Bbound$ with functions on the shifted
cotangent bundle of the derived zero locus
(Theorem~\ref{thm:dcrit-shiftedcot}).

\emph{Vanishing of the quartic obstruction.}
The Gaussian surface is strict, the Lie surface has only the finite
Lie--Jacobi cubic layer, and the contact surface is admitted here only
after a nullhomotopy of $\mathfrak{q}_4(\cA)$ has been specified.
Proposition~\ref{prop:quartic-contact-vanishes-BL} supplies this
nullhomotopy in the strict boundary-linear Landau--Ginzburg sector.
Hence no weight-$4$ planted-forest class remains in the present
triangle.

\emph{Formal reduction.}
The chirally Koszul hypothesis gives
$\cC_{\mathrm{line}} \simeq
\cA^!_{\mathrm{line}}\text{-}\mathrm{mod}$
(Theorem~\ref{thm:lines_as_modules}).
Proposition~\ref{prop:formal-global-triangle} then completes
the chain: derived centers are Morita invariant, so
$\Zder(\Bbound) \simeq \Zder(\cC_{\mathrm{line}})
\simeq \HH^\bullet_{\mathrm{ch}}(\cA^!_{\mathrm{line}})$.
\end{proof}

\begin{remark}[Class~$M$ gap: Virasoro and $\cW_N$]%
\label{rem:class-M-global-triangle-gap}%
\index{global triangle!class M obstruction|textbf}%
\index{Virasoro algebra!global triangle obstruction}%
For class~$M$ algebras (Virasoro, $\cW_N$ for $N \geq 3$), the
global triangle is proved in the boundary-linear sector but
\emph{not} globally. Two obstructions are active:
\begin{enumerate}[label=(\roman*)]
\item \textbf{Non-formal $A_\infty$ structure.}
 The higher operations $m_k \neq 0$ for all $k \geq 3$
 (infinite algebraic depth), so
 $\mathfrak{q}_4 \neq 0$
 (Remark~\ref{rem:quartic-contact-beyond-BL}).
 The quartic resonance class
 $\mathfrak{Q}^{\mathrm{contact}}_{\mathrm{Vir}} =
 10/[c(5c+22)]$ is a concrete nonzero instance.
 The pointwise reduction of the ordered bar complex to
 the derived center is obstructed at degree three and above.

\item \textbf{Nonlinear DS superpotential.}
 The Virasoro and $\cW_N$ families arise via Drinfeld--Sokolov
 reduction from affine algebras. The DS boundary conditions
 produce a superpotential that is nonlinear in the normal
 directions, so the derived HKR theorem
 (Theorem~\ref{thm:boundary-linear-bulk-boundary}) does not
 apply. A direct computation on the ordered bar complex,
 avoiding the boundary-linear shortcut, would resolve the
 question but is open.
\end{enumerate}

For obstruction~(ii), partial progress comes from the
nodal degeneration gluing programme of
Nafcha~\cite{Nafcha26}: the factorisation-algebraic
gluing theorem for nodal curves provides chain-level
control of the bar complex across degenerations, which
is the geometric input needed to extend the boundary-linear
argument past the linear-in-normals restriction.

\emph{Status update.} The extension to full DS boundary conditions is a
conditional implication on the stated locus of
Theorem~\ref{thm:holographic-recon-class-m}: the Costello--Gaiotto
$3$d HT theory, the completed DS bar-coalgebra/coderivation transfer
package, and the DS--Hochschild bridge
(Theorem~\ref{thm:ds-hochschild-bridge}) give
the identification
$\mathrm{Bulk}(\cA)\simeq
Z^{\mathrm{der}}_{\mathrm{ch}}(\cA)\simeq
\mathcal{O}(M_{\mathrm{vac}}^{\mathrm{der}})$
cohomologically, and at chain level in the weight-completed/pro
ambient. On the raw original complex this identification requires the
bounded conformal-weight-shift HPL criterion of
Remark~\ref{rem:class-M-hpl-convergence}. The obstructions listed
above are the reasons the boundary-linear HKR shortcut is not a
substitute for the DS bridge.
\end{remark}

\begin{remark}[Class~$M$ via DS transport]%
\label{rem:class-M-DS-transport-strategy}%
Strategy~C reduces the Virasoro global triangle to
Hochschild--DS compatibility.
The triangle holds for $V_k(\mathfrak{sl}_2)$ (class~$L$).
The Drinfeld--Sokolov reduction
$V_k(\mathfrak{sl}_2)\to \Vir_c$ is a strong deformation retract,
and the ordered bar construction commutes with DS.
Commutation of chiral Hochschild cochains with DS would transport the triangle
for $V_k(\mathfrak{sl}_2)$ to the Virasoro algebra.
The missing input is an HPL transfer for the internal Hom functor
$\RHom_A(M,N)$ through the DS strong deformation retract.
The ghost-number bottleneck that trivializes the DS coproduct acts on the same filtration and forces strictness of the transferred Hom.
At fixed conformal weight this is finite-dimensional linear
algebra: SDR data are explicit; the higher homotopies of the transferred $\RHom$ are computed weight by weight.

\emph{Cohomological compatibility.}
The DS-induced map on chiral Hochschild cohomology is verified at all degrees:
$\mathrm{ChirHoch}^0$: $\C \xrightarrow{\sim} \C$ (centre to centre);
$\mathrm{ChirHoch}^1$: $\mathfrak{sl}_2 \to 0$ (three outer derivations killed by BRST-exactness, inner absorption, screening);
$\mathrm{ChirHoch}^2$: $\C \xrightarrow{\sim} \C$
($\Theta_k \mapsto (dc/dk)\,\Theta_c$ with
$dc/dk = -6(k{+}1)(k{+}3)/(k{+}2)^2 \neq 0$ at generic $k$);
$\mathrm{ChirHoch}^{\ge 3}$: $0 \to 0$.
The chain-level obstruction is the infinite $\Ainf$ tower of class~$M$ producing chain-level corrections invisible in cohomology.
\end{remark}

\subsection{\texorpdfstring{The circle bar chain model and factorisation homology of $S^1$}{The circle bar chain model and factorisation homology of S1}}
\label{subsec:circle-bar-chain-model}

The interval bar complex $\Barchord(\cA)$ computes factorisation
homology on the interval: $\int_I \cA \simeq \Barchord(\cA)$.
Its chain-level structure — ordered configurations on $\R$
with OPE residues in the holomorphic direction — is the
engine of the genus-$0$ Koszul dual.
Replacing the interval by the circle produces the
chain-level model for Hochschild homology via factorisation homology; the new ingredient is wrap-around monodromy.

\begin{definition}[Circle bar chain model]%
\ClaimStatusProvedHere
\label{def:circle-bar-chain-model}%
\index{circle bar chain model|textbf}%
\index{factorisation homology!circle|textbf}%
\index{bar complex!circle}%
Let $\cA$ be a strongly admissible augmented
$E_1$-chiral algebra with ordered bar coalgebra
$C=\Barchord(\cA)$ and diagonal bicomodule
$C_\Delta$. The \emph{circle bar chain model} is the
chain complex
\[
B^{S^1}_\bullet(\cA)
\;:=\;
C_\bullet\!\bigl(
\overline{\FM}{}^{\mathrm{ann}}_{r,s};\,
\omega_{\mathrm{bar}}
\bigr)
\;\cong\;
\bigoplus_{n\ge 0}
\mathrm{cyc}^n(C,C_\Delta),
\]
where $\overline{\FM}{}^{\mathrm{ann}}_{r,s}$ is the
annular configuration space
($r$~holomorphic points on~$\C$,
$s$~cyclically ordered points on~$S^1$, with
Fulton--MacPherson compactification in the
holomorphic direction) and
$\mathrm{cyc}^n(C,C_\Delta)=
C_\Delta\widehat\otimes_C C^{\widehat\otimes\,n}$ is
the cyclic cotensor.
This is the annular bar complex of
Definition~\textup{\ref{def:annular-bar}}, presented
here as a factorisation (co)homology computation.
\end{definition}

\begin{theorem}[Factorisation homology of the circle;
\ClaimStatusProvedHere]%
\label{thm:S1-factorisation-homology}%
\index{Hochschild homology!as factorisation homology of circle|textbf}%
Let $\cA$ be a strongly admissible $E_1$-chiral
algebra. There is a canonical quasi-isomorphism
\begin{equation}\label{eq:S1-fact-hom}
\int_{S^1}\!\cA
\;\simeq\;
B^{S^1}_\bullet(\cA)
\;\simeq\;
\HH^{\mathrm{ch}}_\bullet(\cA),
\end{equation}
where $\int_{S^1}\cA$ is the factorisation homology
of the $E_1$-algebra underlying~$\cA$ on the circle,
and $\HH^{\mathrm{ch}}_\bullet(\cA)$ is chiral
Hochschild homology.
The content of this identification is that the
chain complex
$B^{S^1}_\bullet(\cA)$, assembled from
annular configurations with the four-term
differential
$d^{\mathrm{ann}}=d_{\mathrm{int}}+d_{\mathrm{res}}
+d_\Delta+d_{\mathrm{wrap}}$
\textup{(}Theorem~\textup{\ref{thm:annular-bar-differential}}\textup{)},
is a concrete chain model for the abstract
factorisation homology $\int_{S^1}\cA$.
\end{theorem}

\begin{proof}
The factorisation homology of an $E_1$-algebra $A$ on
$S^1$ is, by definition, the colimit
$\int_{S^1}A=A\otimes_{A\otimes A^{\mathrm{op}}}^{\mathbf{L}}A$
(the derived self-trace). This is computed by the
cyclic bar complex $B^{\mathrm{cyc}}_\bullet(A)$.
For the $E_1$-chiral algebra~$\cA$, the ordered bar
coalgebra $C=\Barchord(\cA)$ dualises the algebra
structure, and the cyclic cotensor
$\bigoplus_n\mathrm{cyc}^n(C,C_\Delta)$ is the
coalgebraic model of the same derived self-trace:
it computes
$\operatorname{coHH}_\bullet^{\mathrm{ch}}(C,C_\Delta)
\simeq\HH^{\mathrm{ch}}_\bullet(\cA)$
by the bimodule/bicomodule equivalence
(Theorem~\ref{thm:annular-HH-vol2}; the underlying
bimodule/bicomodule argument is inherited from Volume~I).
The four-term differential reflects the four types of
codimension-$1$ boundary strata of the annular
configuration space: internal bubbles, OPE collisions,
non-wrapping coproduct faces, and the wrap-around
at the seam.
\end{proof}

\begin{remark}[Wraparound monodromy for $E_1$-chiral algebras]%
\label{rem:wraparound-monodromy}%
\index{R-matrix!wraparound monodromy|textbf}%
\index{circle bar chain model!wraparound monodromy}%
The wrap-around differential $d_{\mathrm{wrap}}$
encodes the periodicity of $S^1$ at the chain level.
For an $E_\infty$-chiral algebra (with $R$-matrix
$R(z)$ derived from local OPE), the wraparound
acts by the full $R$-matrix monodromy: transport of
the last tensor factor $a_n$ past the seam $0\sim 1$
applies
$\mathrm{Mon}(R)=\lim_{z\to e^{2\pi i}z}R(z)$.
For the pole-free sub-class (BD commutative chiral
algebras), $R(z)=\tau$ is the identity flip
and $d_{\mathrm{wrap}}$ reduces to standard
cyclic rotation: the classical cyclic bar complex.
For vertex algebras with OPE poles,
$\mathrm{Mon}(R)\neq\mathrm{id}$: the Heisenberg algebra has $R(z)=\exp(k\hbar/z)$ and
monodromy $\exp(-2\pi i k)$.

For a genuinely $E_1$-chiral algebra ($R(z)$ an
independent input, not derived from local OPE),
the wraparound involves the full
$R$-matrix monodromy as an independent algebraic
datum. The $E_\infty$ vs.\ $E_1$ distinction is visible at the chain level: in the $E_\infty$ case
$d_{\mathrm{wrap}}$ is determined by $d_{\mathrm{res}}$
(the OPE collision differential); in the $E_1$ case
$d_{\mathrm{wrap}}$ carries genuinely new information
from the $R$-matrix.
\end{remark}

\begin{proposition}[$E_2$-structure from the Deligne conjecture;
\ClaimStatusProvedElsewhere]%
\label{prop:derived-center-E2}%
\index{Deligne conjecture!E2 on derived center@$E_2$ on derived center|textbf}%
\index{derived center!E2 structure@$E_2$ structure}%
Let $\cA$ be an $E_1$-chiral algebra. The derived center
\[
Z_{\mathrm{der}}(\cA)
\;:=\;
\mathrm{RHom}_{\mathrm{Fun}(\cC_\partial,\cC_\partial)}(\mathrm{id},\mathrm{id})
\;\simeq\;
C^\bullet_{\mathrm{ch}}(\cA,\cA)
\]
carries a natural $E_2$-algebra structure, as a
consequence of the Deligne conjecture
\textup{(}proved by
Kontsevich~\textup{\cite{Kon03}},
Tamarkin~\textup{\cite{Tamarkin00}}\textup{)}
applied to $E_1$-algebras
\textup{(}Theorem~\textup{\ref{thm:tamarkin-higher-structure}}\textup{)}.
Dually, the factorisation homology
$\int_{S^1}\cA\simeq\HH^{\mathrm{ch}}_\bullet(\cA)$
carries an $S^1$-action encoding the cyclic rotation of
the circle, whose homotopy fixed points recover the
negative cyclic homology $\mathrm{HC}^-_\bullet(\cA)$.
\end{proposition}

\begin{theorem}[Chiral Higher Deligne, Hochschild-chapter echo;
\ClaimStatusConditional{} for the two-coloured universal and strict
$E_3$ refinements; \ClaimStatusProvedElsewhere{} for the brace
structure]%
\label{thm:chiral-higher-deligne-hoch-echo}%
\phantomsection\label{thm:chiral-higher-deligne-hoch}%
\index{Deligne conjecture!chiral higher|textbf}%
\index{higher Deligne!chiral|textbf}%
\index{E3@$E_3$!on chiral derived center}%
Let $\cA$ be a logarithmic $\SCchtop$-algebra (for example any
logarithmic chiral algebra; Volume~I Definition of logarithmic
SC-algebra). Write
$Z^{\mathrm{der}}_{\mathrm{ch}}(\cA):=
C^\bullet_{\mathrm{ch}}(\cA,\cA)$ for the chiral derived center.
Then:
\begin{enumerate}
 \item[\textup{(i)}] $Z^{\mathrm{der}}_{\mathrm{ch}}(\cA)$ carries a
 canonical chiral brace, equivalently $E_2$-chiral, structure
 refining Proposition~\ref{prop:derived-center-E2}.
 \item[\textup{(ii)}] The pair
 $(Z^{\mathrm{der}}_{\mathrm{ch}}(\cA),\cA)$ is the two-coloured
 $\SCchtop$ bulk-boundary pair when the heptagon hypotheses and the
 required two-coloured cobar contracting homotopy are supplied; the
 bulk-to-boundary action factors through the universal brace
 (\S\ref{sec:chiral-Hochschild-brace}).
 \item[\textup{(iii)}] A strict $E_3$-chiral refinement is available
 only under the explicit topologisation hypotheses of
 Theorem~\ref{thm:chiral-higher-deligne} and
 Convention~\ref{conv:chd-e3-chiral-meaning}. Volume~I Theorem~H
 remains the independent concentration input; it is not derived here
 from a generic strict $E_3$ action.
\end{enumerate}
\end{theorem}

\begin{proof}
This is a chapter-local echo of
Theorem~\ref{thm:chiral-higher-deligne}. Part~(i) is the
Kontsevich--Tamarkin/chiral-brace action already recorded in
Proposition~\ref{prop:derived-center-E2}. Part~(ii) is the
two-coloured heptagon statement, with the same conditional
two-coloured cobar homotopy as in
Proposition~\ref{prop:chd-two-coloured-cobar-obstruction}. Part~(iii) is exactly
the topologised lane of Convention~\ref{conv:chd-e3-chiral-meaning}.
\end{proof}

\begin{remark}[Climax placement: universal holography master and the $E_\infty$ ladder]%
\label{rem:higher-deligne-climax-placement}%
\index{universal holography!chiral Higher Deligne input}%
\index{$E_\infty$-topological horizon!chiral Higher Deligne input}%
\index{Virasoro!chiral Higher Deligne instance}%
The chiral Higher-Deligne theorem above is the load-bearing
\emph{open}/$E_3$-on-bulk side of the universal holography master
theorem (Theorem~\ref{thm:universal-holography-master}). The
specialisation ladder is:
\begin{itemize}[nosep]
\item \textbf{Rung 1 (Virasoro, $N=2$).} For $\cA = \Vir_c$
 with the universal Sugawara antighost $G^{(2)}$ this realises the Virasoro
 climax. The $E_3$-action on $Z^{\mathrm{der}}_{\mathrm{ch}}(\Vir_c)$ is
 the chain-level avatar of 3D quantum gravity, with central curvature
 $\kappaChHodge(\Vir_c) = c/2$. This
 holds unconditionally on the non-logarithmic $C_2$-cofinite standard
 landscape via tempered-stratum membership
 (Chapter~\ref{ch:tempered-stratum-characterization},
 Theorem~\ref{thm:tempered-stratum-contains-virasoro}).
\item \textbf{Rung $N$ ($\cW_N$ principal).}
 The higher-spin tower
 $\{T = T^{(2)},\,W^{(3)},\ldots,W^{(N)}\}$ furnishes
 $N{-}1$ inner stress tensors, each admitting a chain-level antighost
 primitive $G^{(s)}$ with $T^{(s)} = [Q_{\mathrm{tot}}, G^{(s)}]$ on
 cohomology. Iterated Sugawara
 (Theorem~\ref{thm:iterated-sugawara-construction}) lifts $E_3$ on
 $Z^{\mathrm{der}}_{\mathrm{ch}}(\cW_{N,c})$ to $E_{N+1}$-topological. The
 universal ratio is $\beta_N = 12(H_N - 1)$
 (Chapter~\ref{ch:beta-N-closed-form-all},
 Conjecture~\ref{conj:beta-N-harmonic-closed-form}).
\item \textbf{Rung $\infty$ ($\cW_\infty$).}
 The full higher-spin Casimir tower drives the
 specialisation $\bigsqcup_{N\ge 2} \mathrm{Rung}\,N \to E_\infty$-topological
 in the inverse limit; $E_3$ on $Z^{\mathrm{der}}_{\mathrm{ch}}$ at
 each finite $N$ extends to $E_{k+2}$-top for any $k \le N$, and the
 limit is the $E_\infty$-topological horizon of the programme.
\end{itemize}
The chiral Higher Deligne action of the present theorem is the
$N=2$ entry point; the climax chapter reads the
ladder as a single statement, the higher-spin antighost tower being
recorded chapter-by-chapter through the conditional harmonic
$\beta_N$ law.
\end{remark}

\begin{remark}[Three Hochschild invariants]%
\label{rem:three-hochschild-invariants}%
\index{Hochschild!three variants|textbf}%
The literature carries three invariants sharing the
name ``Hochschild'' or ``derived centre'':
\begin{enumerate}
 \item[\textup{(i)}] $\mathrm{ChirHoch}^*(\cA)$: chiral Hochschild
 cochains. A family datum $H_H(\cA;S_\cA)$ computes their support
 inside~$S_\cA$.
 \item[\textup{(ii)}] $\HH^*(\cA_{\mathrm{mode}})$: classical
 Hochschild cochains of the mode algebra.
 For the Weyl mode algebra in the free-field example,
 $\HH^0\cong\Bbbk$ and $\HH^n=0$ for $n>0$.
 \item[\textup{(iii)}] $H^*_{\mathrm{GF}}(\mathrm{Lie}(\cA))$:
 Gel'fand--Fuchs continuous Lie cohomology of the mode Lie algebra.
 Its critical affine instances contain the familiar unbounded
 polynomial families.
\end{enumerate}
These are three typed invariants of three typed source algebras.
For a Weyl mode algebra, item~\textup{(ii)} is concentrated in
degree~$0$; item~\textup{(iii)} supplies the unbounded
Gel'fand--Fuchs polynomial. At critical level, Frenkel--Teleman compute
unbounded self-Ext of the vacuum in the oper lane. Transport to
item~\textup{(i)} requires the completed bar--Ext-to-chart comparison
map. Theorem~\ref{thm:three-hochschild-unification} names the available
comparisons.
\end{remark}

\begin{remark}[Chiral Hochschild comparison at the K3
specialisation via the Universal $\Phi$-Trace Identity;
\ClaimStatusProvedElsewhere]
\label{rem:hoch-trinity-phi-trace}
\index{Hochschild!comparison at K3}%
\index{Phi-trace identity@$\Phi$-trace identity!chiral Hochschild comparison}%
\index{Borcherds pushforward!kappa BKM}%
The Universal $\Phi$-Trace Identity
(in the upstream form
Remark~\ref{rem:bar-cobar-5-phi-trace-conductor}) organises the
three Hochschild invariants of
Remark~\ref{rem:three-hochschild-invariants} into a single
cross-volume trace. At the K3 specialisation
$\cA=\cA_{K3}$, write
$K(\cA_{K3})=\tr_{Z^{\mathrm{der}}_{\mathrm{ch}}(\cA_{K3})}(K_{\cA_{K3}})$
for the Koszul conductor evaluated on the chiral derived centre
(= variant (i) of Remark~\ref{rem:three-hochschild-invariants}).
On the \emph{free-field branch} of the $K3$ landscape (the
$\cA_{K3}=V_{L_{K3}}$ lattice-VA presentation via the Mukai rank-$24$
lattice), the conductor vanishes,
\[
K(\cA_{K3})
\;=\;
\tr_{Z^{\mathrm{der}}_{\mathrm{ch}}(\cA_{K3})}(K_{\cA_{K3}})
\;=\;0,
\]
by the class-$\mathsf{G}$ assignment of the free-field branch
(Theorem~C complementarity constant $K=0$ on the affine lineage,
shared with Vol~I). Under the Vol~III $\Phi$-functor, the
Borcherds singular-theta pushforward on the BKM side reads the
non-vanishing $\kappa_{\mathrm{BKM}}$-invariant off the
\emph{Monstrous $E_1$-lift} input rather than the free-field input:
\[
\Phi_\ast\bigl(K(\cA_{\mathfrak{g}_{\Delta_5}})\bigr)
\;=\;
\kappa_{\mathrm{BKM}}(\mathfrak{g}_{\Delta_5})
\;=\;\frac{c_{\phi_{0,1}^{K3}}(0)}{2}\;=\;5,
\]
with $c_{\phi_{0,1}^{K3}}(0)$ the K3 elliptic-genus coefficient
entering the Borcherds lift and $\Delta_5$ the primitive
Gritsenko--Nikulin denominator; its scalar Igusa square is
$\Phi_{10}^{\mathrm{un}}=\Delta_5^2$.
The three-Hochschild comparison at the K3 specialisation is
branch-sensitive. The free-field branch gives $K=0$ (trivial conductor, matching
variants~(ii)/(iii) of Remark~\ref{rem:three-hochschild-invariants}
being finite on the smooth Weyl/GF content), while the Monstrous
$E_1$-lift branch gives
$\kappa_{\mathrm{BKM}}(\mathfrak g_{\Delta_5})=c_{\phi_{0,1}^{K3}}(0)/2=5$
(the weight of
$\Delta_5$, loaded into variant~(i) via the
$\eta_{\mathrm{mode}}$-insensitive chiral Hochschild content of
Theorem~\ref{thm:three-hochschild-unification}\,(i) in degree
$\geq 2$). Theorem~\ref{thm:three-hochschild-unification} is
accordingly the chiral-side form of the Universal
$\Phi$-Trace Identity, with the branch split realised as the
degree-$0$ (conductor) versus
degree-$\geq 2$ (KZ-data) refinement of~$\eta_{\mathrm{mode}}$.
\end{remark}

\begin{theorem}[Three-Hochschild unification; \ClaimStatusProvedHere]%
\label{thm:three-hochschild-unification}%
\index{Hochschild!unification theorem|textbf}%
\index{Hochschild!three variants, unification}%
Let $\cA$ be a logarithmic chiral algebra with mode algebra
$\cA_{\mathrm{mode}}$ and mode Lie algebra $\mathrm{Lie}(\cA)$.
There are canonical natural transformations
\[
\begin{tikzcd}[row sep=small, column sep=large]
\mathrm{ChirHoch}^*(\cA)
\arrow[r,"\eta_{\mathrm{mode}}"]
\arrow[d,"\eta_{\mathrm{GF}}"']
& \HH^*(\cA_{\mathrm{mode}},\cA_{\mathrm{mode}}) \\
H^*_{\mathrm{GF}}(\mathrm{Lie}(\cA))
\arrow[ur,"\zeta"'] &
\end{tikzcd}
\]
such that:
\begin{enumerate}
 \item[\textup{(i)}] $\eta_{\mathrm{mode}}$ is an isomorphism in
 degree $0$ on the non-critical scalar-centre locus and an injection
 in degree $1$ on the finite PBW/Koszul locus. It is \emph{not} an
 isomorphism in degree $\ge 2$: $\mathrm{ChirHoch}^2$ carries
 genus-$0$ KZ data that is invisible to $\HH^*$. At the critical
 affine level $k=-h^\vee$, degree $0$ is the
 Feigin--Frenkel zero-mode/one-point specialization
 \[
 \Fun(\mathrm{Op}_{\fg^\vee}(D))
 \longrightarrow
 \HH^0(\cA_{\mathrm{mode}},\cA_{\mathrm{mode}}),
 \]
 with infinite-dimensional kernel; it is not an isomorphism in this
 uncompleted mode-algebra ambient.
 \item[\textup{(ii)}] $\eta_{\mathrm{GF}}$ factors through the
 continuous cohomology of the positive half
 $\mathrm{Lie}^+(\cA)$ and is an isomorphism on
 the Cartan-stable part. At critical level, $\eta_{\mathrm{GF}}$
 is surjective onto $H^*_{\mathrm{GF}}$; off critical, it lands in
 degrees $\le 2$ only.
 \item[\textup{(iii)}] $\zeta$ is the CE/HKR character sending a
 continuous cochain to the scalar Hochschild class obtained from its
 deformation-quantisation trace.  When \(\cA_{\mathrm{mode}}\) is a
 Weyl algebra, \(\HH^\bullet(\cA_{\mathrm{mode}},
 \cA_{\mathrm{mode}})=k\) in degree \(0\); consequently \(\zeta\)
 identifies only the degree-$0$ quotient of
 \(H^\bullet_{\mathrm{GF}}(\mathrm{Lie}(\cA))\) with the unit summand
 of \(\HH^0\), and all positive-degree Gel'fand--Fuchs classes lie in
 its kernel.
\end{enumerate}
The classical comparison therefore has the opposite shape from a
chiral-only uniqueness claim: in the Weyl and mode-algebra cases
\(\HH^*(\cA_{\mathrm{mode}})\) is \emph{more} concentrated than
\(\mathrm{ChirHoch}^*(\cA)\), while the chiral complex retains the
spectral and higher-configuration data.
\end{theorem}

\begin{proof}
$\eta_{\mathrm{mode}}$: Restrict the configuration data defining
$\mathrm{ChirHoch}$ to a formal disk at a single point and
mode-expand. This kills the spectral parameters, producing a cochain
of $\cA_{\mathrm{mode}}$; the resulting map is the coderivation
comparison induced by mode expansion (Volume~I modes-as-BD
restriction).  On the non-critical scalar-centre locus, the only
degree-$0$ central chiral class is the scalar vacuum class, and it maps
to the scalar centre of the mode algebra.  The degree-$1$ injection on
the finite PBW/Koszul locus is the same direct mode-expansion
calculation applied to infinitesimal automorphisms.  At critical level
this proof no longer gives an isomorphism: the degree-$0$ source is
\(\Fun(\mathrm{Op}_{\fg^\vee}(D))\), and the mode comparison is the
Feigin--Frenkel--Reshetikhin zero-mode/one-point specialization with
infinite-dimensional kernel.
$\eta_{\mathrm{GF}}$: Apply the Chevalley--Eilenberg functor to
the boundary $E_1$-chiral structure, obtaining a continuous Lie
cohomology class. Cartan-stable-part compatibility: at non-critical
level, the CE differential is deformed by the Sugawara vector, and
the comparison is by explicit spectral sequence (Feigin--Frenkel;
Arakawa).
$\zeta$: HKR for Weyl algebras is classical in the following precise
form:
\(\HH^\bullet(\mathrm{Weyl}_n,\mathrm{Weyl}_n)=k\) in degree \(0\).
The continuous Gel'fand--Fuchs classes therefore cannot all survive in
Weyl Hochschild cohomology.  The map \(\zeta\) is the
deformation-quantisation trace character; it identifies the
degree-$0$ quotient of the GF algebra with the unit class and kills the
positive-degree GF generators.  Critical level singularity:
at $k=-h^\vee$, the Feigin--Frenkel centre makes
$\mathrm{ChirHoch}^*$ infinite-dimensional in degree $0$; the ordinary
mode Hochschild complex remains finite in the uncompleted ambient, so
the comparison is the finite quotient described in clause~(i), not a
degree-$0$ isomorphism.
\end{proof}

\begin{proposition}[Algebraic--geometric chiral Hochschild models are equivalent;
\ClaimStatusConditional]%
\label{thm:chiral-hochschild-models-equivalent}%
\index{chiral Hochschild!algebraic vs geometric models|textbf}%
\index{FM-calculus!chiral Hochschild model}%
\index{chiral Hochschild!algebraic-geometric equivalence}%
Let $\cA$ be a logarithmic $\SCchtop$-algebra. Define:
\begin{enumerate}
 \item[\textup{(a)}] The \emph{algebraic} model
 $C^\bullet_{\mathrm{ch,alg}}(\cA)
 :=\prod_{n\ge 0}\mathrm{End}^{\mathrm{ch}}_\cA(n+1)[-n]$,
 where $\mathrm{End}^{\mathrm{ch}}_\cA(n)=
 \mathrm{Hom}(\cA^{\otimes n},\cA(\!(\lambda_1)\!)\cdots(\!(\lambda_{n-1})\!))$
 with differential $\delta f=[m,f]_{\mathrm{G}}$ (Gerstenhaber bracket).
 \item[\textup{(b)}] The \emph{geometric} model
 $C^\bullet_{\mathrm{ch,geom}}(\cA)$: sections over $\FM_{n+1}(X)$
 with logarithmic forms, differential $d=d_{\mathrm{int}}+d_{\mathrm{fact}}+d_{\mathrm{config}}$.
\end{enumerate}
After local coordinate trivialisation on a formal disk (four-step
chain: choose point, choose coordinate, trivialise D-module,
identify $\lambda_i$ with relative positions), the two models are
naturally quasi-isomorphic as brace dg algebras:
\[
\Psi\colon C^\bullet_{\mathrm{ch,alg}}(\cA)
\xrightarrow{\;\simeq\;} C^\bullet_{\mathrm{ch,geom}}(\cA).
\]
The quasi-isomorphism intertwines the Gerstenhaber brackets
and, on cohomology, the $E_3$-structures of
Theorem~\ref{thm:chiral-higher-deligne}.
Chain-level filtrations (excess of closed inputs; genus filtration)
are preserved up to explicit spectral-sequence corrections
on the geometric side.
\end{proposition}

\begin{proof}
$\Psi$ is the BD local-to-global comparison (Volume~I
FM-calculus), assembled with the four-step trivialisation of the
BD chiral operad against the algebraic $\mathrm{End}^{\mathrm{ch}}$
(choose a point, choose a formal coordinate, trivialise the
$D$-module, identify spectral variables with relative positions).
Quasi-isomorphism is by the standard argument: both sides compute
$\mathrm{RHom}_{\SCchtop}(\cA,\cA)$ in the derived category of
$\SCchtop$-algebras over the formal disk.
Brace compatibility: the algebraic brace operations
(\S\ref{sec:chiral-Hochschild-brace}) agree with the geometric brace
operations (iterated fibre integration over FM strata) after
$\Psi$, by Volume~I Proposition
\ref*{prop:geometric-braces-well-defined}.
$E_3$-structure preserved: the pentagon construction of
Theorem~\ref{thm:chiral-higher-deligne} uses only data common to both
models (two-colour SC operad).
\end{proof}

\begin{definition}[Explicit half-braiding]%
\label{def:half-braiding-explicit}%
\index{half-braiding!explicit construction|textbf}%
\index{R-matrix!half-braiding origin}%
Let $\cA$ be an $E_1$-chiral bialgebra with spectral coproduct
$\Delta_z\colon\cA\to\cA\otimes\cA$
(Vol~III notation; Sweedler: $\Delta_z(a)=\sum a_{(1),z}\otimes a_{(2),z}$).
For a left $\cA$-module $N$ and $a\in\cA$, $n\in N$, set
\[
\sigma_\cA(z)(a\otimes n)
\;:=\;
\sum\bigl(\Delta_z(a)\bigr)_{(2)}\cdot n\;\otimes\;\bigl(\Delta_z(a)\bigr)_{(1)}
\;=\;\sum a_{(2),z}\cdot n\otimes a_{(1),z}.
\]
This is the \emph{half-braiding} of $\cA$ acting on $N$. The family
$\{\sigma_\cA(z)\}$ is the $R$-matrix of $\cA$, viewed as a natural
transformation of bifunctors on $\mathrm{Mod}_\cA\times\mathrm{Mod}_\cA$.
\end{definition}

\begin{proposition}[Half-braiding is $R(z)$ and makes $\mathrm{Mod}_\cA$ braided;
\ClaimStatusProvedHere]%
\label{prop:half-braiding-R-matrix}%
\index{R-matrix!from half-braiding}%
\index{quantum group!braiding construction}%
The half-braiding $\sigma_\cA(z)$ of
Definition~\ref{def:half-braiding-explicit}:
\begin{enumerate}
 \item[\textup{(i)}] satisfies the hexagon axioms (quasi-triangularity
 translated to the half-braiding);
 \item[\textup{(ii)}] equals the spectral $R$-matrix
 $R(z)\in\mathrm{End}(V\otimes V)(\!(z)\!)$ on evaluation modules
 $V_u,V_v$ with $z=u-v$;
 \item[\textup{(iii)}] makes $\mathrm{Mod}_\cA$ a braided monoidal
 category (=$E_2$ in $\mathrm{Cat}$), \emph{not} $\cA$ itself.
\end{enumerate}
The formula $R(z)=(\mathrm{id}\otimes S)\circ\Delta_z(1_\cA)$
collapses to $1\otimes 1$ by the counit and therefore cannot encode
braiding. Definition~\ref{def:half-braiding-explicit} is the
nontrivial half-braiding formulation.
\end{proposition}

\begin{proposition}[DS--Hochschild class dichotomy: chain-level vs
cohomological scope; \ClaimStatusProvedHere{} as a class-stratified
criterion]%
\label{prop:ds-hochschild-class-dichotomy}%
\index{DS reduction!Hochschild compatibility!class dichotomy|textbf}%
\index{class M@class~$\mathsf{M}$!chain-level gap}%
Let $\cA$ be a chiral algebra carrying a conformal vector at
non-critical level, and fix the DS--Hochschild comparison map
\[
\mathrm{HKR}^{\mathrm{DS}}\colon
C^\bullet_{\mathrm{ch}}(\cA,\cA)\longrightarrow
\cO\bigl(M_{\mathrm{vac}}^{\mathrm{der}}(\cA)\bigr).
\]
Then the comparison splits by shadow class as follows.
\begin{enumerate}[label=\textup{(\roman*)}]
\item\label{item:ds-dich-GLC} \emph{Classes $\mathsf{G},\mathsf{L},\mathsf{C}$
 (chain level).} The transferred $A_\infty$-tower on
 $H^\bullet(\barB(\cA))$ has $m_k^{\mathrm{tr}}=0$ for $k$ beyond the
 class-dependent depth $r_{\max}(\cA)\le 4$ (depth $2$ for $\mathsf{G}$,
 $3$ for $\mathsf{L}$, $4$ for $\mathsf{C}$). The HPL transfer through
 the DS strong deformation retract terminates at arity $r_{\max}(\cA)$,
 and $\mathrm{HKR}^{\mathrm{DS}}$ is a quasi-isomorphism of chain
 complexes.
\item\label{item:ds-dich-M-coh} \emph{Class $\mathsf{M}$ (cohomological).}
 For $\cA=\cW^k(\fg)$ (principal DS at non-critical level, so class
 $\mathsf{M}$), the cohomological comparison is conditional on the
 smooth chiral-HKR/derived-vacuum identification and the completed
 DS coderivation transfer package.  Generic non-critical level is not
 a $C_2$-cofinite input.  When the smooth-HKR package is supplied,
 $\mathrm{HKR}^{\mathrm{DS}}$ is a quasi-isomorphism on cohomology.
\item\label{item:ds-dich-M-chain} \emph{Class $\mathsf{M}$ (chain level, open).}
 The class-$\mathsf{M}$ transferred tower $m_k^{\mathrm{tr}}$ does not
 truncate ($r_{\max}=\infty$; Proposition~on $\cW_N$ class $\mathsf M$
 above). Consequently chain-level Stasheff cancellation of cross-arity
 contractions in the HPL transfer requires one extra ingredient:
 \emph{convergence of the HPL transfer in the pro-nilpotent completion
 associated to the conformal-weight filtration}
 $F^w\cA=\bigoplus_{\Delta\ge w}\cA_\Delta$.
\end{enumerate}
The gap in~\ref{item:ds-dich-M-chain} is precisely the statement that,
for $\cA$ of class $\mathsf{M}$, the formal sum
$\sum_{k\ge 3} h_k\circ m_k^{\mathrm{tr}}\circ h_k$ of the HPL
correction recursion converges weight-by-weight in $F^\bullet\cA$ with
bounded conformal-weight shift at each arity. Equivalently, the
cross-arity Stasheff cancellation diagrams
$\mathrm{St}_n\colon \bigoplus_{i+j=n}m_i\circ_\ell m_j=0$ stabilise
inside each $F^w$-graded piece after finitely many recursion steps.
\end{proposition}

\begin{proof}
(i) For classes $\mathsf{G},\mathsf{L},\mathsf{C}$ the $A_\infty$-depth
bound $r_{\max}\le 4$ is the operadic complexity theorem (Vol~I,
Theorem~\ref*{V1-thm:operadic-complexity-detailed}; see the shadow-depth
table in Section~\ref{subsec:punctured-disk-bar}). HPL transfer through
the DS contraction~\cite{Arakawa15} is a finite sum and thus converges
on the nose; the resulting chain map is a quasi-isomorphism because the
DS cohomology is concentrated in degree $0$ (Arakawa).
(ii) For $\cA=\cW^k(\fg)$, the finite-dimensionality used by the
comparison is per conformal weight and comes from the Kazhdan/PBW
filtration and finite strong-generation profile, not from generic
$C_2$-cofiniteness.  The smooth chiral-HKR hypothesis identifies the
associated graded of the chiral Hochschild complex with the
polyvector algebra on the derived vacuum space, while the completed
DS coderivation transfer identifies the DS-reduced side.  Under those
two supplied data, the bar/conformal-weight spectral sequence compares
the same \(E_2\)-page and gives a cohomology isomorphism.
(iii) Chain-level promotion fails generically: with $r_{\max}=\infty$,
the HPL recursion produces a formal infinite tower of corrections, each
shifting conformal weight by a bounded amount; convergence is weight-
filtration-adic and is not implied by $E_2$-degeneration alone (the
latter controls only the graded pieces, not the completion of the
tower). The named convergence criterion is what would close the gap.
\end{proof}

\begin{remark}[Ingredient that would close class~$\mathsf{M}$ chain level]%
\label{rem:class-M-hpl-convergence}%
\index{class M@class~$\mathsf{M}$!HPL convergence criterion}%
Fix $\cA=\cW^k(\fg)$. Let $(\pi,\iota,h)$ be the DS contraction onto
$H^\bullet_{\mathrm{DS}}(\cA)$. A sufficient condition for
Proposition~\ref{prop:ds-hochschild-class-dichotomy}\ref{item:ds-dich-M-chain}
to upgrade from cohomological to chain-level equivalence is:
\begin{quote}
\emph{Bounded conformal-weight shift HPL}: there exists a constant
$C=C(\cA)$ such that for every transferred operation
$m_k^{\mathrm{tr}}$ the homotopy correction $h\circ m_k^{\mathrm{tr}}$
shifts conformal weight by at most $C$, independent of $k$.
\end{quote}
Under this criterion, the formal infinite tower evaluates weight-by-
weight inside each conformal-weight stratum.  The relevant strata are
finite-dimensional by the Kazhdan/PBW finite-weight profile of the
BRST complex; this is different from \(C_2\)-cofiniteness and holds
at generic non-critical level without asserting lisse-ness.  The HPL
recursion then converges in the $F^\bullet$-adic completion. The
criterion is known for Virasoro at
$c\notin\{c_{p,q}\}_{\mathrm{minimal}}$ by direct Feigin--Fuchs weight
counting~\cite{FeiginFuchs90}; it is open for
$\cW_N$ with $N\ge 3$ and for generic $\cW^k(\fg)$. Establishing the
bounded-shift criterion would close the DS--Hochschild compatibility
at chain level for class~$\mathsf{M}$; the present manuscript uses the
cohomological form~\ref{item:ds-dich-M-coh} wherever chain-level
class-$\mathsf{M}$ data is not independently available.
\end{remark}

\begin{theorem}[DS--Hochschild bridge; \ClaimStatusProvedHere
\textup{(}scoped by Proposition~\textup{\ref{prop:ds-hochschild-class-dichotomy}}\textup{)}]%
\label{thm:ds-hochschild-bridge}%
\index{DS reduction!Hochschild compatibility|textbf}%
\index{bulk!universal at chain level (class M)}%
\index{class M!HKR identification}%
Let $\cA=\cW^k(\fg)$ be a $W$-algebra obtained by Drinfeld--Sokolov
reduction of an affine Kac--Moody algebra $V^k(\fg)$ at non-critical
level. Write $M_{\mathrm{vac}}^{\mathrm{der}}$ for the derived moduli
stack of vacuum solutions of the DS-modified BV complex.
Then there is a natural HKR-type comparison
\[
\mathrm{HKR}^{\mathrm{DS}}\colon
C^\bullet_{\mathrm{ch}}(\cW^k(\fg),\cW^k(\fg))
\;\longrightarrow\;
\mathcal{O}(M_{\mathrm{vac}}^{\mathrm{der}}).
\]
It is a cohomological quasi-isomorphism for class~$\mathsf M$ under
Proposition~\ref{prop:ds-hochschild-class-dichotomy}\ref{item:ds-dich-M-coh}.
It is a chain-level quasi-isomorphism only in the pro/weight-completed
ambient, or on the original complex after the bounded conformal-weight
shift HPL criterion of
Remark~\ref{rem:class-M-hpl-convergence} has been verified.
Consequently, the universal algebraic bulk identification
$\mathrm{Bulk}(\cW^k(\fg))\simeq Z^{\mathrm{der}}_{\mathrm{ch}}(\cW^k(\fg))$
holds cohomologically for class~$\mathsf M$, and at chain level only
in the completed/pro sense just stated.
\end{theorem}

\begin{proof}
(1) Costello--Gaiotto exhibit a 3d holomorphic-topological theory
whose boundary is $\cW^k(\fg)$ via holomorphic Chern--Simons with
DS boundary conditions; the improvement term
$T^{\mathrm{imp}}=[Q_{\mathrm{tot}},G']$ is Cartan-valued
(Vol II Theorem \ref{thm:E3-topological-DS}, \ref{thm:E3-topological-DS-general}).
(2) The smooth chiral-HKR/derived-vacuum package supplies the
cohomological comparison target.  The finite-weight Kazhdan/PBW
profile supplies the associated graded finiteness per conformal
weight; no generic \(C_2\)-cofiniteness assertion is used.
(3) The chain-level HKR map is constructed from the DS-modified
BV complex: the $Q_{\mathrm{tot}}$-cohomology is by construction
$\mathcal{O}(M_{\mathrm{vac}}^{\mathrm{der}})$; the chiral Hochschild
cochains compute the same object via edge 4--5 of the pentagon
(convolution model, Theorem~\ref{thm:chiral-higher-deligne}).
(4) Compatibility with DS: the bar-cobar/DS intertwining
(cross-volume bridge DS-bar, Volume~I ds-koszul-intertwine)
commutes with HKR on the cohomological comparison. Chain-level
compatibility for class~$\mathsf M$ is the completed/pro statement
unless the bounded-shift criterion of
Remark~\ref{rem:class-M-hpl-convergence} is supplied for the original
complex. This furnishes the class~$\mathsf M$ universal algebraic bulk
cohomologically, and records the exact extra input required for an
uncompleted chain-level identification.
\end{proof}

\begin{remark}[HKR scope and chiral analogue]%
\label{rem:hkr-scope-and-chiral-analogue}%
\index{HKR theorem!scope}%
\index{HKR theorem!chiral analogue}%
The HKR isomorphism enters this chapter in two guises whose scopes
differ, and the Universal Holography bridge depends on the distinction.
\emph{Classical HKR.} For a smooth commutative algebra~$A$ over a
characteristic-zero base,
$\HH^{\bullet}(A)\simeq\mathcal{O}(T^*[-1]\Spec A)$
as $\mathbb{P}_2$-algebras; smoothness is essential, as it identifies
the cotangent complex with the module of K\"ahler differentials
concentrated in degree zero.
\emph{Chiral HKR (Khan--Zeng).} For a freely-generated PVA~$V$ with
polynomial $\gr_{\mathrm{Li}}V$, the chiral Hochschild complex admits
the analogous identification
$\ChirHoch^{\bullet}(V)\simeq
 \mathcal{O}(T^*[-1]\mathrm{DerM}_{\mathrm{vac}}(V))$
as chiral $\mathbb{P}_2$-algebras, the isomorphism coming from the
Li filtration together with the smoothness of $\mathrm{DerM}_{\mathrm{vac}}(V)$
on the freely-generated locus.
\emph{Scope (generic central charge).} Virasoro at generic~$c$, and
more generally any class~$M$ vertex algebra with smooth
$\mathrm{DerM}_{\mathrm{vac}}$ on its universal locus, satisfies the
chiral HKR hypothesis; the Universal Holography bridge
(Theorem~\ref{thm:ds-hochschild-bridge}) gives
$\mathrm{Bulk}\simeq Z^{\mathrm{der}}_{\mathrm{ch}}$ cohomologically
and in the weight-completed/pro chain-level ambient.  The original
complex requires the bounded-shift HPL or Banach estimate stated in
Remark~\ref{rem:class-M-hpl-convergence}.
\emph{Scope break (minimal models).} At minimal central charges
$c_{p,q}$, null-vector ideals cut out singular strata of
$\mathrm{DerM}_{\mathrm{vac}}$; chiral HKR remains a quasi-isomorphism
generically on the smooth locus but fails along the singular strata,
where the cotangent complex is no longer concentrated in a single
degree. Consequently, for minimal-model central charges the Universal
Holography identification $\mathrm{Bulk}\simeq Z^{\mathrm{der}}_{\mathrm{ch}}$
is conditional on a derived resolution of
$\mathrm{DerM}_{\mathrm{vac}}(\mathrm{Vir}_{c_{p,q}})$ along the
null-vector strata (conjectural; compare
Remark~\ref{rem:ds-hochschild-downstream} on class~$M$ generic case).
The forthcoming downstream consequences are stated for the generic-$c$
stratum; minimal-model extensions are flagged where they arise.
\end{remark}

\begin{remark}[Downstream consequences for class~M bulk--boundary identifications]%
\label{rem:ds-hochschild-downstream}%
Theorem~\ref{thm:ds-hochschild-bridge} combined with
Theorem~\ref{thm:chiral-higher-deligne} gives the global triangle
\[
\mathrm{Bulk}\simeq Z^{\mathrm{der}}_{\mathrm{ch}}\simeq
\HH^{\mathrm{ch}}
\]
on the admissible boundary-linear/DS locus. For class~$\mathsf M$ the
triangle is cohomological or weight-completed/pro unless the
bounded-shift HPL criterion of
Remark~\ref{rem:class-M-hpl-convergence} supplies raw original-complex
descent. Thus class~$\mathsf M$ is reached by the DS route, not by the
boundary-linear HKR shortcut.
\end{remark}

\begin{theorem}[Feigin--Frenkel chiral Hochschild at critical level;
\ClaimStatusProvedHere]%
\label{thm:feigin-frenkel-chirhoch}%
\index{Feigin-Frenkel@Feigin--Frenkel!chiral Hochschild|textbf}%
\index{critical level!ChirHoch@$\ChirHoch$}%
\index{class FF}%
\index{opers!Langlands dual disc}%
Let $\cA = V_{-h^\vee}(\fg)$ be the affine Kac--Moody vertex algebra at
the critical level $k=-h^\vee$. Then:
\begin{enumerate}[label=\textup{(\roman*)}]
\item\label{item:ff-chh0} The degree-$0$ chiral Hochschild cohomology
identifies with the Feigin--Frenkel centre:
\[
\ChirHoch^0\bigl(V_{-h^\vee}(\fg),V_{-h^\vee}(\fg)\bigr)
\;\simeq\;
\mathfrak z\bigl(\widehat{\fg}\bigr)
\;\simeq\;
\Fun\bigl(\mathrm{Op}_{\fg^\vee}(D)\bigr),
\]
the ring of functions on the scheme of classical $\fg^\vee$-opers on
the formal disc $D$.
\item\label{item:ff-higher} Under the completed bar--oper comparison
of Volume~I Theorem~\ref*{V1-thm:fifth-class-FF},
\[
 \ChirHoch^n\bigl(V_{-h^\vee}(\fg)\bigr)
 \cong \Omega^n\!\bigl(\mathrm{Op}_{\fg^\vee}(D)\bigr)
 \qquad(n\geq0).
\]
The completed critical chart therefore has unbounded support.
\item\label{item:ff-threeway} In the three-Hochschild diagram of
Theorem~\ref{thm:three-hochschild-unification}, the critical level
$k=-h^\vee$ carries the degree-zero comparison
\[
 \eta_{\mathrm{mode}}^0\colon
 \Fun(\mathrm{Op}_{\fg^\vee}(D))
 \longrightarrow \HH^0(\cA_{\mathrm{mode}}).
\]
Its source is the infinite-dimensional Feigin--Frenkel centre; the
map itself determines the mode-algebra quotient.
\item\label{item:ff-ds} Principal Drinfeld--Sokolov reduction at the
critical level recovers the classical $\cW$-algebra:
$H^0_{\mathrm{DS,prin}}(V_{-h^\vee}(\fg))\simeq
\Fun(\mathrm{Op}_{\fg^\vee}(D))^{\mathrm{cl}}$ (Arakawa, Frenkel).
Each non-principal nilpotent orbit requires its corresponding critical
DS calculation.
\end{enumerate}
\end{theorem}

\begin{proof}
Part~\ref{item:ff-chh0}: By the Feigin--Frenkel
theorem~\cite{FeiginFrenkel1992,Frenkel2007}, the centre of the vertex
algebra $V_{-h^\vee}(\fg)$ coincides with the polynomial algebra
$\mathfrak z(\widehat\fg)$ generated by $\mathrm{rank}(\fg)$ families of
Sugawara-type fields of conformal weights $d_i+1$ (exponents of $\fg$),
and is isomorphic to $\Fun(\mathrm{Op}_{\fg^\vee}(D))$ via the
Miura map.  Degree-zero chiral Hochschild cohomology is the chiral
centre by definition of the derived endomorphism complex.
Part~\ref{item:ff-higher} is the completed bar--oper comparison of
Volume~I Theorem~\ref*{V1-thm:fifth-class-FF}.  Part
\ref{item:ff-threeway} is the degree-zero component of the native
comparison map $\eta_{\mathrm{mode}}$.
Part~\ref{item:ff-ds}: Arakawa~\cite{Arakawa15} computes
$H^0_{\mathrm{DS,prin}}(V_{-h^\vee}(\fg))$ as the classical
$\cW$-algebra $\Fun(\mathrm{Op}_{\fg^\vee}(D))^{\mathrm{cl}}$,
completing the square (Feigin--Frenkel duality + classical limit at
$k=-h^\vee$).
\end{proof}

\begin{remark}[Compatibility with twisted $D$-modules on $\Bun_G$ at the critical level]%
\label{rem:twisted-d-mod-critical-level}%
\index{Bun G@$\Bun_G$!critical twist}%
\index{Beilinson-Drinfeld@Beilinson--Drinfeld!Hecke eigensheaves}%
\index{Frenkel-Gaitsgory@Frenkel--Gaitsgory!localisation}%
\index{opers!global}%
Theorem~\ref{thm:feigin-frenkel-chirhoch} has a global geometric
counterpart on the moduli stack $\Bun_G$ of principal $G$-bundles on
$X$. The category $D_{\mathrm{crit}}(\Bun_G)$ is the category of
$D$-modules on $\Bun_G$ twisted by the critical line bundle: the
square root $\Omega^{1/2}_{\Bun_G}$ of the canonical line bundle, whose
curvature matches the critical-level shift $k=-h^\vee$ on $V_k(\fg)$.
Beilinson and Drinfeld~\cite[Chapter~7]{BD04} prove that the
Hecke eigensheaves on $D_{\mathrm{crit}}(\Bun_G)$ are parametrised by
global $\fg^\vee$-opers on $X$; the local (formal-disc) version of this
correspondence is precisely the Feigin--Frenkel identification
$\mathfrak z(\widehat\fg)\simeq\Fun(\mathrm{Op}_{\fg^\vee}(D))$ used in
part~\ref{item:ff-chh0} above. In the programme's language,
Theorem~\ref{thm:feigin-frenkel-chirhoch} is the chiral-algebraic
(local, formal-disc) incarnation of the Beilinson--Drinfeld
Hecke-at-critical correspondence: the ring $\ChirHoch^0$ on the local
side matches the spectral parameter space of Hecke eigensheaves on the
global side.

Compatibility between the two frameworks is realised by the
Frenkel--Gaitsgory localisation
functor~\cite{FG06,FrenkelLanglands}
\[
\Loc_{\mathrm{crit}}\colon D_{\mathrm{crit}}(\Bun_G)\longrightarrow
\widehat{\fg}_{\mathrm{crit}}\text{-}\Mod
\;\simeq\;\Rep\bigl(V_{-h^\vee}(\fg)\bigr),
\]
which sends a critically twisted $D$-module on $\Bun_G$ to its
completed sections as a module over the affine Kac--Moody algebra at
the critical level; global opers parametrising Hecke eigensheaves land
in the Feigin--Frenkel centre under $\Loc_{\mathrm{crit}}$, aligning
the two spectral decompositions.

The construction is established for simply-connected classical $G$ by
Beilinson--Drinfeld~\cite{BD04} combined with
Frenkel--Gaitsgory~\cite{FG06}; extensions of the local and global
pictures to non-simply-laced and adjoint groups are available in
Frenkel~\cite{FrenkelLanglands} and subsequent work, at a level
sufficient to cover all simple types relevant to the programme.
\end{remark}

\begin{remark}[Class $\mathrm{FF}$ structural invariants; companion to
$\mathsf{G}/\mathsf{L}/\mathsf{C}/\mathsf{M}$]%
\label{rem:class-FF-structural}%
\index{class FF@class $\mathrm{FF}$!structural invariants|textbf}%
\index{fingerprint!coset FF@$\mathrm{coset}_{\mathrm{FF}}$}%
The fifth class $\mathrm{FF}$ (Feigin--Frenkel) collects chiral algebras
with $\kappaChHodge=0$; it is a \emph{companion} stratum of the
$\mathsf{G}/\mathsf{L}/\mathsf{C}/\mathsf{M}$ quadrichotomy, not a region
of exclusion. Structural fingerprint slots
$\varphi(\cA)=(p_{\max},r_{\max},\chi_{\mathrm{VOA}},n_{\mathrm{strong}},\mathrm{coset})$
specialise as follows for $\mathrm{FF}$.
\begin{enumerate}[label=\textup{(\roman*)}]
\item $\kappaChHodge=0$ (defining condition); the Sugawara tensor
$T^{\mathrm{Sug}}=\frac{1}{2(k+h^\vee)}\sum_a{:}J^a J^a{:}$
(Vol~I convention, OPE-mode normalisation; the Vol~II $\lambda$-bracket
presentation carries an additional $1/n!$ divided-power factor)
degenerates because its normalisation diverges at $k=-h^\vee$, so the
usual conformal vector is
absent and the two colours of $\SCchtop$ do not topologise into
$E_3^{\mathrm{top}}$. The open colour remains $E_1^{\mathrm{top}}$ and
the closed colour remains $E_2^{\mathrm{chiral}}$; the pair is
$\SCchtop$, not upgraded.
\item $r_{\max}=\infty$ (as for class $\mathsf{M}$), but distinguished
from $\mathsf{M}$ by the \emph{shape} of the shadow tower: class
$\mathsf{M}$ has quartic OPE pole and generic central charge
$c\notin\{c_{p,q}\}$; class $\mathrm{FF}$ has degenerate Sugawara,
polynomial $\mathrm{ChirHoch}^0$ on opers, and a distinct set of
shadow-tower coefficients (Gelfand--Dikii generators contribute in
every degree).
\item The fifth fingerprint slot is replaced by
$\mathrm{coset}_{\mathrm{FF}}=(\mathfrak z(\widehat\fg),\mathrm{Op}_{\fg^\vee}(D))$:
the commutant pair encodes the Feigin--Frenkel centre together with its
Langlands-dual oper geometry.
\item Koszul structure: the bar complex
$\bar B^{\mathrm{ch}}(V_{-h^\vee}(\fg))$ is a well-defined
$\mathrm{Ass}^{\mathrm{ch}}$-coalgebra; its Koszul dual is
$\Fun(\mathrm{Op}_{\fg^\vee}(D))$, which is itself polynomial
(Koszul-self-dual up to the Gelfand--Dikii shift). The bar-cobar
inversion $\Omega^{\mathrm{ch}}\bar B^{\mathrm{ch}}\simeq\cA$ does
\emph{not} hold on the nose at $k=-h^\vee$: both sides become
commutative vertex algebras, but the identification is via the
Feigin--Frenkel isomorphism, not via the usual pro-nilpotent
completion.
\end{enumerate}
\textbf{Scope.} The theorem above is proved for affine KM at the
critical level (Feigin--Frenkel 1992). Principal DS reduction at the
critical level reduces class $\mathrm{FF}$ to classical
$\cW$-algebras, an observation due to Arakawa and Frenkel. Non-principal
DS at the critical level, and the $\kappaChHodge=0$ locus for
other families (critical-level extended algebras, logarithmic
$\cW(p)$ at $p=1$) remain open and are targets for a companion
theorem.
\textbf{Nomenclature.} The fifth class is named $\mathrm{FF}$
(Feigin--Frenkel) after its degree-$0$ oper content; the informal
synonym ``wild'' refers to the same stratum, the structural content
being the $\kappaChHodge=0$ critical-level fingerprint of this remark.
\end{remark}

\begin{remark}[Igusa $\Delta_5$ as the
Hochschild-cohomological gauge comparator of the $K3$ chiral
Hall--Drinfeld double; \ClaimStatusProvedElsewhere]
\label{rem:hoch-igusa-gauge-class}
\index{Igusa form!Hochschild gauge comparator}%
\index{Hall--Drinfeld double!K3!Hochschild invariant}%
On the Monstrous $E_1$-lift
$\mathfrak{g}_{\Delta_5}$ (Vol~III Part~\ref{part:frontier},
the K3 Hall--Drinfeld branch; compare
Remark~\ref{rem:bar-cobar-1-igusa-gauge}), the Vol~II Hochschild
complex supplies only the deformation chart. The comparison with
Igusa's line is conditional on the finite Hall/CoHA source, the
Mukai/Serre pairing, PBW/no-extra-relations, radical quotient,
$\Z/2$ parity, filtered completion, inverse-limit, and EK/Heegner
comparison gates. After those gates, the target map is
\[
\widehat H^2_{\mathrm{red}}\!
  \bigl(\mathfrak{g}_{\Delta_5}\bigr)^{\Z/2,\,K(1)}
\;\longrightarrow\;
\C\cdot\Delta_5,
\]
to the line generated by Igusa's Siegel weight-$5$ cusp form
$\Delta_5$ on $\mathsf{Sp}_4(\Z)$. Under the Hochschild presentation
of the chiral derived centre
$Z^{\mathrm{der}}_{\mathrm{ch}}(\cA_{\mathfrak{g}_{\Delta_5}})
=C^\bullet_{\mathrm{ch}}(\cA,\cA)$
(Definition~\ref{def:chiral_hochschild}), this map is a comparator for
the $K3$ chiral Hall--Drinfeld double: a candidate class in
degree-$2$ chiral Hochschild must survive the $\Z/2$ parity involution,
the $K(1)$ isotropic-lattice reduction, and the radical quotient before
it can be compared with $\Delta_5$. It is \emph{not} a member of the
class-FF stratum above (whose degree-$0$ content is the
Feigin--Frenkel oper algebra). Instead, the target line
$\C\cdot\Delta_5$ is paired in degree~$2$ against the
Koszul-conductor class
$K(\cA_{\mathfrak{g}_{\Delta_5}})$ of
Remark~\ref{rem:hoch-trinity-phi-trace} via the Universal
$\Phi$-Trace Identity, with expected weight-$5$ output
$\kappa_{\mathrm{BKM}}(\mathfrak{g}_{\Delta_5})=c_{\phi_{0,1}^{K3}}(0)/2=5$.
This is the Hochschild-cohomological reading of the Vol~III
$\Phi$-bridge as a gated comparison: the Igusa cusp form is the
automorphic target line for the Hall--Drinfeld double, with the pairing
structure sitting alongside the FF fingerprint slots of
Remark~\ref{rem:class-FF-structural}, not replacing them).
\end{remark}

\begin{remark}[Triplet $\cW(p)$: class~M shadow is sound; $\Ethree$-topologization criterion via $\Z_p$ orbifold]%
\label{rem:wp-class-M-and-E3}%
\index{triplet algebra!class M!E_3-topological via orbifold}%
The logarithmic triplet algebra $\cW(p)$ ($p\ge 2$) is
$C_2$-cofinite, non-rational, and is realized as the kernel of
the long screening operator acting on the rank-$1$ lattice
vertex algebra $V_L$ with $L=\sqrt{2p}\,\Z$
\cite{AdamovicMilas08,FeiginTipunin10}.
Three points, in first-principles (a)/(b)/(c) format.
\textbf{(a) What class~M assignment gets right.} The shadow
tower of $\cW(p)$ is computed Feigin--Tipunin-equivariantly:
the Virasoro subalgebra $\mathrm{Vir}_{c_{p,1}}\hookrightarrow\cW(p)$
with $c_{p,1}=13-6(p+p^{-1})$ is honestly class~M because the
stress tensor retains its quartic OPE pole
$T(z)T(w)\sim(c_{p,1}/2)(z-w)^{-4}$ after passing to the kernel of
the screening (screenings are weight-$1$ primaries, they do not
touch the quartic central term). The $T$-channel therefore forces
$r_{\max}^T=\infty$ regardless of $c$ exactly as recorded in
Vol.~I \cite[\S\ref{V1-sec:wp-shadow-depth}]{VolI}. Class~M for
$\cW(p)$ is unconditional at the level of the shadow tower.
\textbf{(b) What a naive reading gets wrong.} Logarithmic
vertex algebras carry nilpotent $L_0$ action on projective
modules (Jordan blocks of size $2$ for $p\ge 2$), and one might
worry that this Jordan structure truncates or nilpotently shifts
the shadow tower. It does not: the nilpotent part of $L_0$ lives
on \emph{modules}, whereas the shadow depth $r_{\max}$ is an
invariant of the VA OPE, i.e.\ of the adjoint action of $\cW(p)$
on itself, where $L_0$ is diagonalizable with integer weights.
The Jordan blocks enter the module category (hence affect $\Cat_{\Mod}$
and the logarithmic modular $S$-action), not the self-OPE. The
class~M assignment of Vol.~I
Table~\ref{tab:landscape-shadow-classes} is accordingly sound.
\textbf{(c) Correct $\Ethree$-topological statement.} The
topologization climax
(Theorem~\ref{thm:E3-topological-km}, Theorem~\ref{thm:E3-topological-DS-general})
does not apply \emph{verbatim} to $\cW(p)$ because $\cW(p)$ is
not freely generated (strong generation is not free: four strong
generators are known to satisfy at least one further relation),
and is not a DS reduction of an
affine at non-critical level. The correct route is the
orbifold route of case~\ref{item:scope-orbifold} of
Remark~\ref{rem:E3-top-free-PVA-scope}: $\cW(p)$ is the
$\Z_p$-invariant sub-VA of $V_L\otimes V_{L^*}$ under the
long-screening $\Z_p$ symmetry, and $V_L$ is class~G and is
$\Ethree$-topological via abelian holomorphic Chern--Simons.
The Monster orbifold route of
Theorem~\ref{thm:monster-orbifold-e3}/Remark~\ref{rem:monster-orbifold-route}
generalizes only as a descent criterion: $\cW(p)$ inherits
$\Ethree$-topological structure from $V_L$ provided the
Dijkgraaf--Witten anomaly class
$\alpha_p\in H^3(B\Z_p,U(1))$ is trivialized by an explicit normalized
cocycle/coboundary computation, or by a type-specific cited theorem,
and provided finite-orbifold BV descent identifies the descended
boundary with the triplet algebra. \emph{Scope.} For $p=2$ the
orbifold is the symplectic-fermion orbifold, but the normalized
$\Z/2$ test is still the same one: triple sign \(+1\) is the zero
class and triple sign \(-1\) is the generator. Even lattice data and
modularity of the fermion partition function are compatibility checks,
not a BV anomaly computation. Thus the $\Ethree$-topological lift of
$\cW(p)$ is conditional for every $p\ge 2$ until the corresponding
local class \(\alpha_p\) and BV descent are supplied; the
modular-tensor-category structure of
Flandoli--Lentner--Runkel~\cite{FlandoliLentner18} is evidence for the
expected vanishing, not a replacement for the cocycle calculation.
\textbf{Summary.} The class~M shadow statement is unconditional, while
the $\Ethree$-topological lift reduces to the finite-orbifold problem
\(\alpha_p=0\) plus BV descent. The argument separates the
strong-vs-free-generation distinction, the Heisenberg-as-abelian-KM
topologization route, the Monster-orbifold criterion, and the orbifold
case of Remark~\ref{rem:E3-top-free-PVA-scope}.
\end{remark}

\begin{remark}[Genus-$1$ curvature]%
\label{rem:circle-bar-curvature}%
\index{curvature!genus-1 on circle bar complex}%
The genus-$1$ curvature of the bar complex is visible
in the circle bar chain model.
The curved differential on
$B^{S^1}_\bullet(\cA)$ satisfies
\begin{equation}\label{eq:circle-bar-curvature}
(d^{\mathrm{ann}})^2 \;=\; \kappaChHodge(\cA)\cdot\omega_1,
\end{equation}
where $\omega_1$ is the Hodge class on
$\overline{\cM}_{1,1}$ and $\kappaChHodge(\cA)$ is the central
curvature of~$\cA$
(Theorem~\ref{thm:hochschild-bridge-higher-genus}\,(iii)
and Volume~I, Theorem~D).
When $\kappaChHodge(\cA)\neq 0$ (all non-trivial standard
families), the circle bar complex is \emph{curved}:
it is not a chain complex but a curved chain complex,
and the factorisation homology
$\int_{S^1}\cA$ is correspondingly a curved object
over~$\overline{\cM}_{1,1}$.

The physical interpretation is that
$\kappaChHodge(\cA)\cdot\omega_1$ is the one-loop anomaly of
the chiral algebra on the torus: the failure of the
genus-$0$ differential to square to zero when the
topological direction is compactified reflects the
Hodge-theoretic obstruction to extending $\cA$
from the plane to genus~$1$.
\end{remark}

\begin{remark}[Modular group action at chain level]%
\label{rem:SL2Z-chain-level}%
\index{modular group!chain-level action|textbf}%
\index{SL(2,Z)@$\mathrm{SL}(2,\Z)$!chain-level action on circle bar}%
The torus $T^2=S^1_A\times S^1_B$ has
$\mathrm{SL}(2,\Z)$ as its mapping class group.
The circle bar chain model
$B^{S^1}_\bullet(\cA)$ is constructed by
compactifying the topological direction ($A$-cycle);
the holomorphic direction ($B$-cycle) enters through
the OPE data in $d_{\mathrm{res}}$.
The $S$-transformation exchanges the two cycles
and therefore interchanges the roles of the
topological wraparound
($d_{\mathrm{wrap}}$, $A$-cycle monodromy)
and the holomorphic collision structure
($d_{\mathrm{res}}$, $B$-cycle OPE data).
This exchange acts at the chain level on
$B^{S^1}_\bullet(\cA)$, permuting the
components of $d^{\mathrm{ann}}$.

The $T$-transformation $\tau\mapsto\tau+1$
acts by the Dehn twist along the $A$-cycle,
which inserts a monodromy factor
$\mathrm{Mon}(R)$ into $d_{\mathrm{wrap}}$.
For $E_\infty$-chiral algebras, this is
determined by the OPE poles; for $E_1$-chiral
algebras, it is the independent $R$-matrix
monodromy.
The full $\mathrm{SL}(2,\Z)$-action at chain
level is the modular content of the
genus-$1$ bar complex, connecting to the
modular PVA quantisation programme
\textup{(}Section~\textup{\ref{subsec:cyclic-pairing-bar}}
of the modular PVA chapter\textup{)}
and to the annular bar complex of
Chapter~\textup{\ref{chap:ordered-associative}}
\textup{(}Definition~\textup{\ref{def:annular-bar}}\textup{)}.
\end{remark}

\subsection{\texorpdfstring{The punctured disk bar chain model and the $E_2$--$E_1$ gap}{The punctured disk bar chain model and the E2--E1 gap}}
\label{subsec:punctured-disk-bar}

The circle $S^1$ and the punctured disk
$D^*:=D\setminus\{0\}$ are homotopy equivalent: the
retraction $D^*\to S^1$ collapsing the radial direction
is an equivalence. At the level of homology,
$H^{\mathrm{ch}}_*(D^*,\cA)\cong\HH^{\mathrm{ch}}_*(\cA)$.
At the level of chains the two models differ:
the punctured disk chain complex
$C^{\mathrm{ch}}_*(D^*,\cA)$ retains the full
holomorphic collision data (OPE poles, spectral
parameters, $E_2$ structure), while the circle chain
complex $C^{\mathrm{top}}_*(S^1,\cA)$ retains only the
topological/cyclic data ($E_1$ structure).
This chain-level discrepancy is the bar-complex
manifestation of the $E_2$--$E_1$ gap, and is the load-bearing transition from Rung~1 to Rung~2 in the geometric escalation (Preface, \S\,0).

\subsubsection{\texorpdfstring{Configuration spaces on $D^*$ versus $S^1$}{Configuration spaces on punctured D versus S1}}

\begin{definition}[Punctured disk configuration space]%
\label{def:punctured-disk-config}%
\index{configuration space!punctured disk|textbf}%
\index{punctured disk!configuration space}%
Let $D^*=D\setminus\{0\}\subset\C$.
The \emph{punctured disk configuration space}
$\FM_k(D^*)$ is the Fulton--MacPherson
compactification of the space of $k$~distinct
labelled points in~$D^*$: it is the pullback
\[
\FM_k(D^*)
\;:=\;
\FM_k(\C)\times_{\C^k}(D^*)^k,
\]
retaining all codimension-$1$ boundary strata from
collisions of subsets $S\subseteq\{1,\ldots,k\}$
($|S|\ge 2$), together with the boundary stratum
at $|z_i|\to 0$ (approach to the puncture) and
the constraint $|z_i|<1$ (containment in the disk).

Compare this with
$\Conf_k(S^1)=\{(\theta_1,\ldots,\theta_k)\in
(S^1)^k:\theta_i\neq\theta_j\}$,
the (uncompactified) configuration space of $k$~distinct
points on the circle, which carries only angular/cyclic
data with no holomorphic collision structure.
\end{definition}

The retraction $D^*\to S^1$, $z\mapsto z/|z|$, collapses the
radial coordinate. For \(k>1\) it does \emph{not} identify the full
configuration-space operadic modules:
\[
H_\bullet(\Conf_k(D^*))\not\simeq H_\bullet(\Conf_k(S^1))
\quad\text{as operadic modules}.
\]
It gives a quasi-isomorphism only after passing to the angular
\(E_1\)-quotient. At the level of differential forms the pullback
$\Omega^\bullet(\Conf_k(S^1))\hookrightarrow
\Omega^\bullet_{\log}(\FM_k(D^*))$ is an inclusion whose
cokernel contains the logarithmic forms with holomorphic poles
at collision divisors -- the forms encoding OPE residues. The
relative \(E_2/E_1\) gap is generated by the angular winding classes
\[
\theta_i=d\arg z_i
\]
and the pairwise logarithmic collision classes
\[
\omega_{ij}=d\log(z_i-z_j).
\]
The angular quotient remembers only the cyclic order and wraparound
monodromy; the full punctured-disk model still remembers the
holomorphic collision forms.

\subsubsection{The punctured disk bar chain model}

\begin{definition}[Punctured disk bar chain model
$B^{D^*}_\bullet(\cA)$; \ClaimStatusProvedHere]%
\label{def:punctured-disk-bar}%
\index{bar complex!punctured disk|textbf}%
\index{punctured disk!bar chain model|textbf}%
Let $\cA$ be a strongly admissible $E_1$-chiral algebra
with ordered bar coalgebra $C=\Barchord(\cA)$.
The \emph{punctured disk bar chain model} is the
chain complex
\begin{equation}\label{eq:punctured-disk-bar}
B^{D^*}_\bullet(\cA)
\;:=\;
\bigoplus_{k\ge 0}
\Omega^\bullet_{\log}\bigl(\FM_k(D^*)\bigr)
\,\widehat\otimes\,
\cA^{\widehat\otimes\,k},
\end{equation}
equipped with the differential
\begin{equation}\label{eq:punctured-disk-diff}
d^{D^*}
\;=\;
d_{\mathrm{int}} + d_{\mathrm{res}} + d_{\mathrm{cyc}},
\end{equation}
where:
\begin{enumerate}[label=\textup{(\roman*)}]
\item $d_{\mathrm{int}}$ is the internal differential
 (BV-BRST operator $Q=m_1$ on each tensor factor);
\item $d_{\mathrm{res}}$ is the collision residue
 differential. For each subset $S\subseteq\{1,\ldots,k\}$
 with $|S|\ge 2$, the codimension-$1$ boundary stratum
 $D_S\subset\FM_k(D^*)$ parametrises configurations where
 the points $\{z_i\}_{i\in S}$ collide. On a generator
 $\omega\otimes(a_1\otimes\cdots\otimes a_k)$, the
 contribution from $D_S$ is
 \[
 d_{\mathrm{res}}^S
 \;=\;
 \operatorname{Res}_{D_S}\!\bigl(\omega\bigr)
 \,\otimes\,
 \bigl(a_1\otimes\cdots\otimes
 m_{|S|}(a_{i_1},\ldots,a_{i_{|S|}})
 \otimes\cdots\otimes a_k\bigr),
 \]
 where $S=\{i_1<\cdots<i_{|S|}\}$,
 $\operatorname{Res}_{D_S}$ extracts the logarithmic
 residue along $D_S$ (the coefficient of
 $d\log(z_{i_1}-z_{i_2})\wedge\cdots$ in the
 expansion of $\omega$ near $D_S$), and
 $m_{|S|}$ is the $|S|$-ary $\Ainf$-operation.
 The total collision differential is
 $d_{\mathrm{res}}=\sum_{|S|\ge 2}d_{\mathrm{res}}^S$,
 identical to the collision differential on $\FM_k(\C)$;
\item $d_{\mathrm{cyc}}$ is the cyclic/wraparound
 differential encoding the periodicity of the
 angular coordinate on $D^*\simeq(0,1)\times S^1$,
 with $R$-matrix monodromy insertion as in the
 circle bar chain model
 \textup{(}Remark~\textup{\ref{rem:wraparound-monodromy}}\textup{)}.
\end{enumerate}
The identity $(d^{D^*})^2=0$ follows from the same
codimension-$2$ cancellation argument as
Theorem~\textup{\ref{thm:bar_differential_squared}},
applied to the boundary complex of $\FM_k(D^*)$.
\end{definition}

The circle bar chain model
$B^{S^1}_\bullet(\cA)$
of Definition~\ref{def:circle-bar-chain-model}
is the topological shadow: $\Conf_k(S^1)$ in place of $\FM_k(D^*)$, differential $d^{\mathrm{ann}}$ retaining only cyclic/topological collision data. The chain complex
$B^{D^*}_\bullet(\cA)$ is strictly larger: it contains
the full logarithmic forms encoding OPE poles,
spectral parameters, and the radial collisions
invisible to $S^1$.

\begin{proposition}[Punctured disk versus circle: rotation gap;
\ClaimStatusProvedHere]
\label{prop:punctured-disk-circle-rotation-gap}
The symmetric punctured-disk chain model
\begin{equation}\label{eq:punctured-disk-symmetric-bar}
B^{D^*,\Sigma}_\bullet(\cA)
=
\bigoplus_{n\ge0}
C_\bullet(\Conf_n(D^*))\otimes_{\Sigma_n}\cA^{\otimes n}
\end{equation}
maps by radial projection to the circle bar model
\[
B^{S^1}_\bullet(\cA)
=
\bigoplus_{n\ge0}
C_\bullet(\Conf_n(S^1))\otimes_{\Sigma_n}\cA^{\otimes n}.
\]
The induced map is a quasi-isomorphism on the angular \(E_1\)
subcomplex, but it is not an equivalence of the full
configuration-space operadic modules. The \(E_2\)-to-\(E_1\) loss is
the loss of the rotation BV operator
\begin{equation}\label{eq:punctured-disk-rotation-bv}
\Delta_{\mathrm{rot}}
=
\sum_i\iota_{\partial_{\arg z_i}}.
\end{equation}
This operator survives on \(D^*\), where the classes
\(\theta_i=d\arg z_i\) and \(\omega_{ij}=d\log(z_i-z_j)\) live in the
logarithmic configuration complex. On \(S^1\) it becomes ordinary
cyclic rotation of the bar word.

Consequently the circle bar computes the annular trace
\[
B^{S^1}_\bullet(\cA)\rightsquigarrow
\operatorname{Tr}^{\mathrm{ann}}(\cA),
\]
whereas the punctured-disk model computes the chiral centre with
rotation,
\[
B^{D^*}_\bullet(\cA)\rightsquigarrow
Z_{\mathrm{der}}^{\mathrm{ch}}(\cA)^{S^1\text{-}\mathrm{rot}}.
\]
They become comparable after angular \(E_1\)-truncation; they are not
the same chain model.
\end{proposition}

\subsubsection{The forgetful quasi-isomorphism}

\begin{proposition}[Punctured disk to circle:
angular \(E_1\) quasi-isomorphism;
\ClaimStatusProvedHere]%
\label{prop:punctured-disk-S1-qiso}%
\index{punctured disk!forgetful quasi-isomorphism|textbf}%
\index{quasi-isomorphism!punctured disk to circle}%
\index{state--field correspondence!chain-level}%
Let $\cA$ be a strongly admissible $E_1$-chiral algebra
with ordered bar coalgebra $C=\Barchord(\cA)$.
Let \(B^{D^*,\mathrm{ang}}_\bullet(\cA)\) denote the angular
\(E_1\)-quotient of \(B^{D^*}_\bullet(\cA)\), obtained by killing the
relative \(E_2/E_1\) subcomplex generated by the rotation and
logarithmic collision forms in
Proposition~\ref{prop:punctured-disk-circle-rotation-gap}. The
retraction \(D^*\to S^1\) induces a quasi-isomorphism
of \(E_1\)-chain complexes
\begin{equation}\label{eq:punctured-disk-S1-qiso}
\varphi\colon
B^{D^*,\mathrm{ang}}_\bullet(\cA)
\;\xrightarrow{\;\sim\;}
B^{S^1}_\bullet(\cA)
\end{equation}
that is \emph{surjective} on angular chains. Concretely, $\varphi$ is
the composite of:
\begin{enumerate}[label=\textup{(\alph*)}]
\item integration over the radial fibre
 $(0,1)\subset D^*$, collapsing $\FM_k(D^*)$ to
 $\Conf_k(S^1)$;
\item restriction of the logarithmic forms
 $\Omega^\bullet_{\log}(\FM_k(D^*))$ to the
 cyclic forms $\Omega^\bullet(\Conf_k(S^1))$,
 discarding the holomorphic pole data in the
 collision divisors.
\end{enumerate}
At the level of the angular quotient, $\varphi$ induces an
isomorphism
\begin{equation}\label{eq:punctured-disk-homology}
H_*\bigl(B^{D^*,\mathrm{ang}}_\bullet(\cA)\bigr)
\;\cong\;
H_*\bigl(B^{S^1}_\bullet(\cA)\bigr)
\;\cong\;
\HH^{\mathrm{ch}}_*(\cA).
\end{equation}
This is not a claim that the full punctured-disk configuration module is
equivalent to the circle configuration module; the full \(D^*\) model
retains the rotation BV operator \(\Delta_{\mathrm{rot}}\).
\end{proposition}

\begin{proof}
Two stages: $\varphi$ is a chain map on the angular quotient, and the
radial fibre is contractible after the \(E_2/E_1\) relative forms have
been quotiented out. Both rest on flatness of the connection defined by
$d_{\mathrm{res}}$ on the ordered bar complex, and homotopy invariance
of the associated local system on $\C^\times$.

\emph{The flat connection and its local system.}
At tensor degree $k$, the bar complex
${\Barch}_k^{\mathrm{ord}}(\cA)$ defines a local
system $\mathcal{F}_k$ on
$\mathrm{Conf}_k^{\mathrm{ord}}(\C^\times)$ with
stalk $(s^{-1}\bar\cA)^{\otimes k}$ and flat
connection $\nabla$ induced by $d_{\mathrm{res}}$
(Construction~\ref{constr:r-matrix-monodromy};
flatness is $d_{\mathrm{res}}^2=0$, equivalent to $(d^{D^*})^2=0$). For $k=2$,
$\nabla=d-\sum_{n\ge 0}r_n\,z^{-n-1}\,dz$ on
$\C^\times$ with fibre $\bar\cA\otimes\bar\cA$,
monodromy around the generator $\gamma$ of
$\pi_1(\C^\times)=\Z$ the $R$-matrix
$R(z)=\operatorname{Mon}_\gamma(\nabla)$
(equation~\eqref{eq:r-matrix-monodromy}).

The connection $\nabla$ is defined on the
\emph{algebraic} de~Rham complex of
$\C^\times=\operatorname{Spec}k[z,z^{-1}]$: the
coefficient $r(z)=\sum_{n\ge 0}r_n\,z^{-n-1}$
is a formal Laurent series in
$\operatorname{End}(\bar\cA\otimes\bar\cA)(\!(z)\!)$,
and the flatness identity $d_{\mathrm{res}}^2=0$ is an
algebraic relation among formal power series.
The monodromy $\operatorname{Mon}(R)$ is therefore
defined as the holonomy of a flat algebraic
connection on the smooth variety $\C^\times$, not as a
contour integral of a convergent function.

\emph{Chain map property.}
The identity
\(\varphi\circ d^{D^*}=d^{\mathrm{ann}}\circ\varphi\)
holds component by component.

For $d_{\mathrm{int}}$: this acts on $\cA$-factors and commutes
with $\varphi$ by functoriality.

For $d_{\mathrm{res}}$: the collision residue at a divisor
$D_S\subset\FM_k(D^*)$ depends only on the relative positions
of the colliding points (the local coordinate
$z_i - z_j$), not on the global radial coordinate $|z_i|$.
The retraction $\varphi$ acts on the global radial coordinate
and is the identity on relative local coordinates near $D_S$,
so $\varphi$ commutes with $d_{\mathrm{res}}$.

For $d_{\mathrm{cyc}}$: the wraparound differential inserts
the monodromy $\operatorname{Mon}(R)$ when transporting a tensor
factor around the angular coordinate. By flatness of $\nabla$,
the holonomy of a flat connection depends only on the homotopy
class of the path. The inclusion $S^1\hookrightarrow D^*$ induces
an isomorphism $\pi_1(S^1)\xrightarrow{\sim}\pi_1(D^*)=\Z$: every
circle $\{|z|=r\}\subset D^*$ represents the same generator of
$\pi_1(D^*)$. Therefore $\operatorname{Mon}(R)$ computed on any
such circle agrees with $\operatorname{Mon}(R)$ on $S^1=\{|z|=1\}$,
by homotopy invariance of holonomy for the flat connection $\nabla$.
This is an algebraic statement: the holonomy of a flat connection
on the algebraic de~Rham complex of $\C^\times$ depends only on
the class in $\pi_1^{\mathrm{alg}}(\C^\times)=\widehat{\Z}$
(or $\pi_1^{\mathrm{top}}(\C^\times)=\Z$ in the analytic
topology). No convergence hypothesis on $R(z)$ is needed;
flatness alone suffices.

The retraction $D^*\to S^1$ preserves winding numbers, so
$\varphi$ commutes with $d_{\mathrm{cyc}}$. This holds
uniformly for $E_\infty$ and $E_1$ algebras.

\emph{Angular quotient.}
The punctured disk $D^*$ and the circle $S^1$ are related by the
inclusion $\iota\colon S^1\hookrightarrow D^*$, which is a homotopy
equivalence for the \emph{one-point} space. Configuration spaces with
several points carry extra braid and logarithmic collision classes, so
the same sentence is false as a statement about the full operadic
modules. After quotienting by the relative \(E_2/E_1\) subcomplex of
Proposition~\ref{prop:punctured-disk-circle-rotation-gap}, the tensor
degree~\(k\) source is the angular de~Rham complex
\[
\Omega^\bullet_{\mathrm{ang}}(\Conf_k(D^*);\mathcal F_k)
:=
\Omega^\bullet_{\mathrm{alg}}(\FM_k(D^*);\mathcal F_k)/
\langle\theta_i,\omega_{ij}\rangle_{\mathrm{rel}},
\]
where \(\mathcal F_k\) is the local system defined by \(\nabla\). This
angular quotient is pulled back from
\(\Conf_k(S^1)\), with coefficients in the pulled-back local system
\(\iota^*\mathcal F_k\).

By homotopy invariance of sheaf cohomology for locally constant
coefficients on this angular quotient, the restriction
\[
\iota^*\colon
H^\bullet_{\mathrm{ang}}(\Conf_k(D^*);\mathcal{F}_k)
\;\xrightarrow{\;\sim\;}
H^\bullet(\Conf_k(S^1);\iota^*\mathcal{F}_k)
\]
is an isomorphism for each~$k$. At the chain level, this means
\(\varphi\) is a quasi-isomorphism on
\(B^{D^*,\mathrm{ang}}_\bullet(\cA)\). The relative classes
\(\theta_i\), \(\omega_{ij}\), and the contraction operator
\(\Delta_{\mathrm{rot}}\) are not declared exact; they are precisely
the full punctured-disk data forgotten by the angular \(E_1\)
projection.

The angular-quotient argument is uniform in the $E_\infty$/$E_1$ distinction:
the only input is flatness of $\nabla$
($d_{\mathrm{res}}^2=0$), the identification
$\pi_1(D^*)=\pi_1(S^1)=\Z$, and homotopy invariance
of cohomology with locally constant coefficients. None of
these require convergence of $R(z)$, holomorphicity of a
radial retraction, or any property specific to
$E_\infty$-chiral algebras.
\end{proof}

\begin{remark}[The $E_1$ case: what the quasi-isomorphism
forgets]%
\label{rem:punctured-disk-S1-E1-gap}%
\index{punctured disk!E1 quasi-isomorphism@$E_1$ quasi-isomorphism|textbf}%
\index{fiber integration!E1 case@$E_1$ case}%
The angular quasi-isomorphism $\varphi$ of
Proposition~\ref{prop:punctured-disk-S1-qiso} is stated
and proved for all strongly admissible $E_1$-chiral
algebras (a class that includes all $E_\infty$-chiral
algebras). The chain map property on the angular quotient uses three inputs:
(i)~collision residues depend on local relative
coordinates,
(ii)~the $R$-matrix $R(z)$ is meromorphic on~$\C^\times$
with poles only at~$z=0$
(Definition~\ref{def:e1-chiral-algebra},
Proposition~\ref{prop:field-theory-r}), and
(iii)~by flatness of the shifted KZ connection
and homotopy invariance of holonomy, the monodromy
$\mathrm{Mon}(R)$ is independent of the radius.
These hold for both $E_1$ and $E_\infty$ algebras.

The distinction between $E_1$ and $E_\infty$ is
\emph{not} about the existence of the angular projection but about
what the projection discards. For $E_\infty$-chiral
algebras, the $R$-matrix $R(z)$ is derived from the
local OPE, and the full holomorphic collision data
encoded in $B^{D^*}_\bullet(\cA)$ is recoverable from
$B^{S^1}_\bullet(\cA)$ via the Deligne conjecture
(which reconstructs the $E_2$-structure on homology).
For genuinely $E_1$-chiral algebras, the independent
$R$-matrix $R(z)$ carries information beyond what the
OPE determines, and the passage from \(D^*\) to the angular
\(S^1\) quotient genuinely loses this data: the circle model
remembers only the monodromy $\mathrm{Mon}(R)$ (a
single operator), not the full spectral family
$R(z)$. The spectral parameter dependence of
$R(z)$, visible in $B^{D^*}_\bullet(\cA)$ through
the holomorphic forms on $\FM_k(D^*)$, is the
chain-level datum killed by the angular \(E_1\)-truncation.
\end{remark}

\begin{remark}[State--field correspondence at chain level]%
\label{rem:state-field-chain-level}%
\index{state--field correspondence!chain vs homology|textbf}%
The angular quasi-isomorphism $\varphi$ of
Proposition~\ref{prop:punctured-disk-S1-qiso}
is the chain-level avatar of the state--field
correspondence. At the homology level, the
state--field correspondence is an isomorphism:
states (elements of the Fock space, i.e., the
fibre at $0\in D$) are identified with fields
(operator-valued distributions on~$D^*$). At
the chain level, this identification factors through the angular
quotient and \emph{forgets} holomorphic
structure: the full punctured disk chain complex
$B^{D^*}_\bullet(\cA)$ carries the full field
data (OPE poles, spectral parameters, radial
ordering), while the circle chain complex
$B^{S^1}_\bullet(\cA)$ carries only the state
data (cyclic ordering, topological wraparound).
The angular projection $\varphi$ passes from fields
to states by integrating out the holomorphic
direction.
\end{remark}

\subsubsection{\texorpdfstring{The $E_2$--$E_1$ gap}{The E2--E1 gap}}

\begin{remark}[The chain-level $E_2$--$E_1$ gap;
punctured disk versus circle]%
\label{rem:E2-E1-chain-gap}%
\index{E2-E1 gap@$E_2$--$E_1$ gap|textbf}%
\index{punctured disk!E2-E1 gap@$E_2$--$E_1$ gap}%
The two chain complexes $B^{D^*}_\bullet(\cA)$
and $B^{S^1}_\bullet(\cA)$ compute the same
homology $\HH^{\mathrm{ch}}_*(\cA)$, which carries
an $E_2$-structure by the Deligne conjecture
\textup{(}Proposition~\textup{\ref{prop:derived-center-E2}}\textup{)}.
The $E_2$-structure is \emph{invisible} in
$B^{S^1}_\bullet(\cA)$: the circle model sees
only the $E_1$-cyclic structure (the wraparound
differential $d_{\mathrm{wrap}}$, the cyclic
ordering on $\Conf_k(S^1)$). The $E_2$-structure
\emph{is} visible in $B^{D^*}_\bullet(\cA)$: the
holomorphic collision data in
$\Omega^\bullet_{\log}(\FM_k(D^*))$ (specifically,
the OPE pole residues and the radial collision
strata) encode the braiding coherences that promote
$E_1$ to $E_2$.

The angular quasi-isomorphism $\varphi$ is therefore
\emph{not} an equivalence of $E_2$-algebras: it is
an equivalence after \(E_1\)-truncation that forgets the
$E_2$-refinement. The $E_2$-structure on
$\HH^{\mathrm{ch}}_*(\cA)$ can be
\emph{recovered} from the circle model $B^{S^1}_\bullet(\cA)$
only via the Deligne conjecture (an existence theorem),
while it is \emph{exhibited} by the punctured disk
model $B^{D^*}_\bullet(\cA)$ (a constructive chain-level
witness).

In the language of the geometric escalation:
$S^1$ is a Rung~1 space (topological, $E_1$)
and $D^*$ is a Rung~2 space (holomorphic, $E_2$).
The projection \(\varphi\) is the \(E_1\)-angular truncation
from Rung~2 to Rung~1. Its angular quotient is homologically
equivalent to the circle model, while the full relative gap complex
contains the rotation BV operator and logarithmic collision classes.
The $E_2$ operations live in this chain-level data.
\end{remark}

\subsubsection{Relation to other bar chain models}

The punctured disk bar chain model sits in the
following network of bar chain models, interpolating
between the formal disk and the circle.

\begin{enumerate}[label=\textup{(\roman*)}]
\item \emph{Formal disk $D$}.
 The vertex algebra structure on $\cA$ is defined on~$D$,
 with configuration spaces $\FM_k(D)$. The punctured
 disk $D^*=D\setminus\{0\}$ is obtained by removing the
 insertion point. The inclusion $D^*\hookrightarrow D$
 induces a restriction map on bar chain models: the
 formal disk carries the full vertex algebra data
 (including the vacuum and the state--field map at
 $z=0$), while $D^*$ carries only the field data away
 from the origin.

\item \emph{Circle $S^1$
 \textup{(}Rung~1, cyclic bar complex\textup{)}}.
 The retraction $D^*\to S^1$ induces the forgetful
 quasi-isomorphism $\varphi$ of
 Proposition~\ref{prop:punctured-disk-S1-qiso},
 discarding holomorphic data.

\item \emph{Annulus $\mathrm{Ann}_\tau\subset D^*$
 \textup{(}Section~\textup{\ref{subsec:annular-bar-complex}}\textup{)}}.
 The annulus $\mathrm{Ann}_\tau=\{r_1<|z|<r_2\}$ embeds
 in $D^*$ as the complement of two concentric disks.
 The annular bar complex
 $B^{\mathrm{ann}}_\bullet(\cA)$
 is the restriction of the punctured disk bar chain
 model to configurations supported on
 $\mathrm{Ann}_\tau$, with the additional structure of
 two boundary circles providing the bimodule datum.
 The conformal modulus $\tau$ of the annulus is visible
 at the chain level in $B^{D^*}_\bullet(\cA)$ through
 the radial position of the configurations, connecting
 to the chain-level modulus dependence of
 Remark~\ref{rem:annular-modulus-dependence}.
\end{enumerate}

\subsection{The annular bar complex}
\label{subsec:annular-bar-complex}

The annulus $\mathrm{Ann}_\tau=\{r_1<|z|<r_2\}$ (conformally equivalent
to $\C^*/q^{\Z}$ with $q=e^{2\pi i\tau}$) has two boundary
circles: an inner circle $S^1_{\mathrm{in}}$ at $|z|=r_1$ and
an outer circle $S^1_{\mathrm{out}}$ at $|z|=r_2$. The boundary
algebra $\cA$ restricts to each boundary component, giving
$\mathrm{Ann}_\tau$ the structure of an $(\cA,\cA)$-bimodule in the
$E_1$-chiral category: the two bulk-to-boundary maps
$\beta_{\mathrm{in}}\colon V_T\to\operatorname{End}(\cA)$
and $\beta_{\mathrm{out}}\colon V_T\to\operatorname{End}(\cA)$
equip the annular observables with a left action from the outer
circle and a right action from the inner circle.

\begin{construction}[Annular bar complex $B^{\mathrm{ann}}(\cA)$;
\ClaimStatusProvedHere]%
\label{constr:annular-bar-vol2}%
\index{annular bar complex!stratified configuration model|textbf}%
\index{bar complex!annular stratified}%
The annular bar complex is the chain complex on the
\emph{annular stratified configuration space}: configurations of
$r$~holomorphic points in the bulk of $\mathrm{Ann}_\tau$ and
$s$~topologically ordered points on each boundary circle,
with Fulton--MacPherson compactification in the holomorphic
direction and cyclic ordering on each $S^1$-factor.
Concretely,
\[
B^{\mathrm{ann}}_\bullet(\cA)
\;:=\;
C_\bullet\!\bigl(
\overline{\FM}{}^{\,\mathrm{ann}}_{r,s};\,
\omega_{\mathrm{bar}}
\bigr)
\;\cong\;
\bigoplus_{n\ge 0}
\mathrm{cyc}^n(C,C_\Delta),
\]
where $C=\Barchord(\cA)$ is the ordered bar coalgebra and
$C_\Delta$ is the diagonal bicomodule
\textup{(}Volume~I, Definition~\textup{\ref{def:annular-bar}}\textup{)}.
The differential decomposes into four components
\textup{(}Volume~I,
Theorem~\textup{\ref{thm:annular-bar-differential}}\textup{)}:
\[
d^{\mathrm{ann}}
\;=\;
d_{\mathrm{int}}+d_{\mathrm{res}}+d_\Delta+d_{\mathrm{wrap}},
\]
where $d_{\mathrm{int}}$ is the internal differential on each
tensor factor, $d_{\mathrm{res}}$ encodes OPE collision residues
in the holomorphic direction, $d_\Delta$ is the coproduct
differential at non-wrapping boundary collisions, and
$d_{\mathrm{wrap}}$ is the wrap-around differential at the seam
of each boundary circle. The two boundary circles contribute
independently to $d_{\mathrm{wrap}}$: the inner-circle seam wraps
the last right-coaction factor, while the outer-circle seam wraps
the first left-coaction factor. The identity
$(d^{\mathrm{ann}})^2=0$ is the algebraic manifestation of
$\partial^2=0$ on the manifold with corners
$\overline{\FM}{}^{\,\mathrm{ann}}_{r,s}$.
\end{construction}

\begin{theorem}[Annular bar computes chiral Hochschild homology;
\ClaimStatusProvedHere]%
\label{thm:annular-HH-vol2}%
\index{annular bar complex!Hochschild identification|textbf}%
\index{Hochschild homology!annular bar model}%
Let $\cA$ be a strongly admissible $E_1$-chiral algebra. There
is a canonical quasi-isomorphism
\begin{equation}\label{eq:annular-HH-vol2}
B^{\mathrm{ann}}_\bullet(\cA)
\;\simeq\;
\operatorname{coHH}^{\mathrm{ch}}_\bullet(C,C_\Delta)
\;\simeq\;
\HH^{\mathrm{ch}}_\bullet(\cA),
\end{equation}
where the first equivalence is tautological
\textup{(}the cyclic cotensor is the coHochschild complex\textup{)}
and the second is the bimodule/bicomodule equivalence
\textup{(}proved in Volume~I\textup{)}.
\end{theorem}

\begin{proof}
The cyclic cotensor
$\bigoplus_n\mathrm{cyc}^n(C,C_\Delta)$ is, by
definition, the coHochschild complex of the coalgebra~$C$
with coefficients in the bicomodule~$C_\Delta$: the coface
maps $\delta_1,\ldots,\delta_{n-1}$ produce $d_\Delta$,
while $\delta_0$ and $\delta_n$ produce $d_{\mathrm{wrap}}$.
The identification
$\operatorname{coHH}^{\mathrm{ch}}_\bullet(C,C_\Delta)
\simeq\HH^{\mathrm{ch}}_\bullet(\cA)$ then follows from
the bimodule/bicomodule equivalence of
Volume~I, which identifies the self-cotrace of the bar
coalgebra with the self-trace of the algebra.
\end{proof}

\begin{remark}[Modulus dependence: chain level vs.\ homology]%
\label{rem:annular-modulus-dependence}%
\index{annular bar complex!modulus dependence|textbf}%
\index{modulus!chain-level vs homology}%
The annular bar complex $B^{\mathrm{ann}}_\bullet(\cA)$
depends on the conformal modulus $\tau$ of the annulus at
the chain level: the collision residue differential
$d_{\mathrm{res}}$ involves the propagator on $\mathrm{Ann}_\tau$,
which is a holomorphic function of~$\tau$. This dependence
is genuinely holomorphic: the annular bar complex forms a
holomorphic family of chain complexes over the moduli space
$\mathfrak{h}/\Z$ of annuli.

At the level of homology, this dependence disappears:
$\HH^{\mathrm{ch}}_\bullet(\cA)$ is a topological invariant
of the circle (the deformation retract of the annulus),
independent of~$\tau$. The chain-level $\tau$-dependence that
washes out in homology is the bar-complex manifestation of the
$E_2$ vs.\ $E_1$ distinction for the circle model. The
$E_2$-structure on
$\HH^{\mathrm{ch}}_\bullet(\cA)$
\textup{(}Proposition~\textup{\ref{prop:derived-center-E2}}\textup{)}
sees only the topological circle; the $E_1$-chain model
$B^{\mathrm{ann}}_\bullet(\cA)$ sees the holomorphic annulus
and its modulus. The spectral sequence from chain-level
$\tau$-dependence to $\tau$-independent homology is the
Hodge-to-de~Rham degeneration for the annular bar complex:
the $E_1$-page carries the holomorphic data, and the
$E_2$-page is the topological invariant.
\end{remark}

\subsection{The pair-of-pants bar chain model}
\label{subsec:pants-bar-chain-model}

The pair of pants $P = \mathbb{P}^1\setminus\{0,1,\infty\}$ is a
genus-$0$ surface with three boundary circles
$\partial_{\mathrm{in}}^1$, $\partial_{\mathrm{in}}^2$,
$\partial_{\mathrm{out}}$. It is the multiplication cobordism:
its bar chain model encodes the product on chiral Hochschild
homology. The preceding annular bar complex
(Construction~\ref{constr:annular-bar-vol2}) lives on a single
boundary circle; the pair-of-pants model composes two such
circles into a third.

\begin{construction}[Pair-of-pants bar chain model;
\ClaimStatusProvedHere]%
\label{constr:pants-bar-model}%
\index{pair of pants!bar chain model|textbf}%
\index{bar complex!pair-of-pants|textbf}%
\index{configuration space!three-punctured sphere}%
Let $\cA$ be a strongly admissible augmented $E_1$-chiral algebra
with ordered bar coalgebra $C = \Barchord(\cA)$ and diagonal
bicomodule~$C_\Delta$. The \emph{pair-of-pants bar chain model}
$B^P_\bullet(\cA)$ is the chain complex on the stratified
configuration space
$\overline{\FM}_n^{\mathrm{ord}}(P)$
of ordered configurations on $P$ with Fulton--MacPherson
compactification. Concretely:
\begin{enumerate}[label=\textup{(\roman*)}]
\item \emph{Generators.}
 A degree-$n$ chain in $B^P_n(\cA)$ is a configuration of
 $n$ labelled points in $P$ equipped with an element of
 $\cA^{\otimes n}$, modulo the bar suspension shift.

\item \emph{Configuration spaces.}
 The compactification
 $\overline{\FM}_n^{\mathrm{ord}}(P)$ has boundary strata
 recording collisions near each of the three punctures
 $\{0\}$, $\{1\}$, $\{\infty\}$. Denote by
 $\partial_0$, $\partial_1$, $\partial_\infty$
 the boundary divisors corresponding to clusters of points
 colliding at $0$, $1$, $\infty$ respectively.

\item \emph{Differential.}
 The total differential decomposes as
 \begin{equation}\label{eq:pants-differential}
 d^P \;=\; d_{\mathrm{int}} + d_{\mathrm{res}}
 + d_{\partial_0} + d_{\partial_1} + d_{\partial_\infty},
 \end{equation}
 where $d_{\mathrm{int}}$ is the internal differential
 on~$\cA$, $d_{\mathrm{res}}$ encodes OPE collision residues
 from the holomorphic direction
 (via $\FM_k(\mathbb{P}^1\setminus\{0,1,\infty\})$),
 and $d_{\partial_0}$, $d_{\partial_1}$,
 $d_{\partial_\infty}$ are the boundary strata contributions
 from collisions at the three punctures. Each boundary
 stratum $\partial_i$ carries an $E_1$ structure: the
 local coordinate at puncture~$i$ identifies collisions
 there with the ordered bar differential on a half-line.

\item \emph{Three $E_1$ structures.}
 The three punctures induce three independent $E_1$ structures:
 \begin{itemize}
 \item $\partial_0$, $\partial_1$ are the two
 \emph{input} circles, each carrying an annular bar complex
 $B^{\mathrm{ann}}_\bullet(\cA)$
 (Construction~\ref{constr:annular-bar-vol2});
 \item $\partial_\infty$ is the \emph{output} circle,
 carrying $B^{\mathrm{ann}}_\bullet(\cA)$.
 \end{itemize}
 The pair-of-pants multiplication composes the two input
 annular bar complexes through $P$ into the output:
 \[
 \mu_P \colon
 B^{\mathrm{ann}}_\bullet(\cA)
 \otimes
 B^{\mathrm{ann}}_\bullet(\cA)
 \longrightarrow
 B^{\mathrm{ann}}_\bullet(\cA).
 \]
 This agrees with the ordered fusion product
 $\mu_P\colon C_\Delta\star C_\Delta\to C_\Delta$
 of Theorem~\textup{\ref{thm:pair-of-pants}}.
\end{enumerate}
\end{construction}

\begin{proposition}[Pair-of-pants product on $\HH^{\mathrm{ch}}_*$;
\ClaimStatusProvedHere]%
\label{prop:pants-product-HH}%
\index{Hochschild homology!pair-of-pants product|textbf}%
\index{pair of pants!product on Hochschild homology}%
The pair-of-pants multiplication $\mu_P$ of
Construction~\textup{\ref{constr:pants-bar-model}} is a chain map
with respect to $d^P$. On homology it induces the associative
product
\begin{equation}\label{eq:pants-HH-product}
 \mu_P^H \colon
 \HH^{\mathrm{ch}}_*(\cA)
 \otimes
 \HH^{\mathrm{ch}}_*(\cA)
 \longrightarrow
 \HH^{\mathrm{ch}}_*(\cA)
\end{equation}
that endows $\HH^{\mathrm{ch}}_*(\cA)$ with the structure of a
graded associative algebra.
\end{proposition}

\begin{proof}
\emph{Chain-map property.}
The map $\mu_P$ is defined by restriction to the
codimension-$1$ boundary stratum of
$\overline{\FM}_{n_1+n_2}^{\mathrm{ord}}(P)$
corresponding to separation of the configuration
into two clusters, one near $\{0\}$ and one near
$\{1\}$. To verify that $\mu_P$ commutes with the
differentials, decompose $d^P$ as
in~\eqref{eq:pants-differential}.
The internal differential $d_{\mathrm{int}}$ acts
tensorially and commutes with $\mu_P$ by functoriality.
For $d_{\mathrm{res}}$: a bulk collision among
points in the first (resp.\ second) input cluster
is a boundary stratum of
$\overline{\FM}_{n_1}^{\mathrm{ord}}$
(resp.\ $\overline{\FM}_{n_2}^{\mathrm{ord}}$), and
$\mu_P$ restricts to the identity on each factor.
A bulk collision involving points from both clusters
lies in codimension~$2$ of
$\overline{\FM}_{n_1+n_2}^{\mathrm{ord}}(P)$
(the cross-ratio must degenerate simultaneously with
the collision), so contributes zero by the
codimension-$2$ vanishing principle. The boundary
terms $d_{\partial_0}$, $d_{\partial_1}$,
$d_{\partial_\infty}$ commute with $\mu_P$ because
$\mu_P$ respects the local coordinate
identifications at each puncture. Since $P$ has
genus~$0$, the bar differential is flat
($(d^P)^2 = 0$, no curvature correction), so $\mu_P$
is a chain map.

\emph{Associativity.}
The two compositions
$\mu_P\circ(\mu_P\otimes\mathrm{id})$ and
$\mu_P\circ(\mathrm{id}\otimes\mu_P)$ correspond to
the two boundary strata of the moduli space
$\overline{\cM}_{0,4}\cong\mathbb{P}^1$:
the stratum at $\lambda=0$ (pinching the circle
separating $\{0,1\}$ from $\{\lambda,\infty\}$) and
the stratum at $\lambda=\infty$ (pinching the
circle separating $\{0,\lambda\}$ from
$\{1,\infty\}$). The interval
$\overline{\cM}_{0,4}\setminus\{0,\infty\}$
connecting these two boundary strata provides a
chain homotopy $h$ satisfying
$\mu_P\circ(\mu_P\otimes\mathrm{id})
-\mu_P\circ(\mathrm{id}\otimes\mu_P)
= d^P\circ h + h\circ d^P$,
constructed by integrating the bar chain model
over the cross-ratio parameter~$\lambda$.
Passage to homology then gives strict associativity
of $\mu_P^H$. The argument at the coalgebraic level is
Theorem~\textup{\ref{thm:pair-of-pants}}.
\end{proof}

\begin{remark}[Degeneration: annulus pinch to pair of pants]%
\label{rem:annulus-pinch-to-pants}%
\index{pair of pants!annulus degeneration}%
\index{degeneration!annulus to pair of pants}%
The product $\mu_P^H$ on $\HH^{\mathrm{ch}}_*(\cA)$ arises
geometrically from the degeneration of an annulus: pinch the waist
circle of an annulus $\mathrm{Ann}_\tau$ to produce a nodal surface
with two components, each a disk with one puncture. Smoothing the
node produces the pair of pants $P$. At the chain level, the
annular bar complex $B^{\mathrm{ann}}_\bullet(\cA)$ degenerates:
the thin-annulus limit $\mathrm{Im}(\tau)\to\infty$ collapses the
annulus into the sewing of two half-planes, and the pair-of-pants
multiplication $\mu_P$ is the chain-level operation that composes
through this sewing. The identity
$(d^P)^2 = 0$ at genus~$0$ reflects the
absence of curvature in the degeneration: the smooth locus of
$\overline{\cM}_{0,3}$ is a single point, and no non-trivial
Arakelov correction appears.
\end{remark}

\subsection{The nodal bar chain model}
\label{subsec:nodal-bar-chain-model}

At a node, a smooth curve degenerates into irreducible components
joined at ordinary double points. The bar chain model on a nodal
curve $\Sigma_{\Gamma}$, indexed by a stable graph
$\Gamma\in\StGraph(g,n)$ labelling the boundary stratum of
$\overline{\cM}_{g,n}$, is the tensor product of vertex bar
complexes sewn at edges by the annular bar complex. The
curvature $\dfib^{\,2} = \kappaChHodge(\cA)\cdot\omega_g$ is
concentrated on the boundary strata: the smooth locus carries a
flat differential, and the nodal degenerations produce the global
$(1,1)$-form $\omega_g$.

\begin{construction}[Nodal bar chain model; \ClaimStatusProvedHere]%
\label{constr:nodal-bar-model}%
\index{nodal curve!bar chain model|textbf}%
\index{bar complex!nodal curve|textbf}%
\index{stable graph!bar complex indexing|textbf}%
Let $\cA$ be a logarithmic $\SCchtop$-algebra
(Definition~\ref{def:log-SC-algebra}), and let
$\Gamma\in\StGraph(g,n)$ be a stable graph with vertex set
$V(\Gamma)$, edge set $E(\Gamma)$, and leg set
$L(\Gamma)$ of cardinality~$n$. Each vertex~$v\in V(\Gamma)$
carries a genus label $g(v)\ge 0$ and a valence
$n(v) = |E_v| + |L_v|$, where $E_v$ is the set of half-edges at~$v$
and $L_v$ is the set of legs at~$v$. The total genus is
$g = \sum_{v\in V(\Gamma)} g(v) + b_1(\Gamma)$,
where $b_1(\Gamma) = |E(\Gamma)| - |V(\Gamma)| + 1$ is the first
Betti number.

The \emph{nodal bar complex on the stratum indexed by $\Gamma$}
is:
\begin{equation}\label{eq:nodal-bar-complex}
 B_\Gamma(\cA)
 \;:=\;
 \biggl(\,
 \bigotimes_{v\in V(\Gamma)}
 B^{(g(v))}_{n(v)}(\cA)
 \biggr)
 \bigotimes_{e\in E(\Gamma)}
 B^{\mathrm{ann}}_e(\cA),
\end{equation}
where the data are as follows.
\begin{enumerate}[label=\textup{(\roman*)}]
\item \emph{Vertex bar complexes.}
 For each vertex $v\in V(\Gamma)$, the complex
 $B^{(g(v))}_{n(v)}(\cA)$ is the genus-$g(v)$ bar complex
 with $n(v)$ marked points: the bar complex of~$\cA$ on the
 smooth component $\Sigma_v$ of genus~$g(v)$ with $n(v)$
 punctures. When $g(v) = 0$, this is the flat bar complex
 (the ordinary bar differential with $(d^{(0)})^2 = 0$).
 When $g(v)\ge 1$, the vertex bar complex carries the
 curvature $\dfib^{\,2} = \kappaChHodge(\cA)\cdot\omega_{g(v)}$ of
 the smooth component.

\item \emph{Edge sewing.}
 For each edge $e\in E(\Gamma)$, connecting vertices~$v_+$
 and~$v_-$ (or a self-edge at a single vertex in the
 non-separating case), the complex $B^{\mathrm{ann}}_e(\cA)$
 is the annular bar complex
 (Construction~\ref{constr:annular-bar-vol2}) placed at the node
 corresponding to~$e$. Concretely, the node is the thin-annulus
 limit of a plumbing collar: a neighbourhood of the node in the
 smoothing family is an annulus $\{|z\,w| < |t|\}$ with plumbing
 parameter~$t$, and $B^{\mathrm{ann}}_e(\cA)$ is the annular bar
 complex on this collar.

 The sewing map at edge~$e$ is
 \begin{equation}\label{eq:sewing-map}
 \sigma_e \colon
 B^{(g(v_+))}_{n(v_+)}(\cA)
 \otimes
 B^{\mathrm{ann}}_e(\cA)
 \otimes
 B^{(g(v_-))}_{n(v_-)}(\cA)
 \longrightarrow
 B_\Gamma(\cA),
 \end{equation}
 given by identifying the output boundary circle of the
 annulus with a puncture of $\Sigma_{v_+}$ and the input
 boundary circle with a puncture of~$\Sigma_{v_-}$.

\item \emph{The tensor product structure.}
 The tensor product in~\eqref{eq:nodal-bar-complex} is taken
 over all vertices and edges simultaneously: the edge annuli
 are the sewing connectors between vertex components. For a
 non-separating edge (self-edge at vertex~$v$), the two
 half-edges both belong to the same vertex and the annular bar
 complex sews two punctures of a single $\Sigma_v$, raising
 the genus by~$1$.

\item \emph{Differential.}
 The total differential on $B_\Gamma(\cA)$ is the sum of the
 vertex differentials and the annular sewing differentials:
 \begin{equation}\label{eq:nodal-bar-differential}
 d_\Gamma \;=\;
 \sum_{v\in V(\Gamma)} d_v
 \;+\;
 \sum_{e\in E(\Gamma)} d^{\mathrm{ann}}_e,
 \end{equation}
 where $d_v$ is the bar differential on the vertex component
 $B^{(g(v))}_{n(v)}(\cA)$ and $d^{\mathrm{ann}}_e$ is the
 annular differential
 (Theorem~\ref{thm:annular-bar-differential})
 on the sewing collar at edge~$e$.
\end{enumerate}
\end{construction}

\begin{proposition}[Curvature localisation on nodal strata;
\ClaimStatusProvedHere]%
\label{prop:nodal-curvature-source}%
\index{curvature!localisation on nodal strata|textbf}%
\index{nodal curve!curvature localisation}%
\index{Arakelov form!nodal degeneration}%
On the smooth locus of $\overline{\cM}_{g,n}$
(the open stratum $\cM_{g,n}$), the bar differential
$d^{(g)}$ on the genus-$g$ bar complex
$\barB^{(g)}(\cA)$ satisfies
$(d^{(g)})^2 = \kappaChHodge(\cA)\cdot\omega_g$ whenever the scalar
continuation of Volume~I, Theorem~D applies; this is
unconditional at genus~$1$ and on the proved scalar
higher-genus lane. Here the curvature arises from the
$\bar\partial$-defect of the Arakelov propagator. The restriction to
the boundary stratum indexed by $\Gamma$ satisfies:
\begin{equation}\label{eq:nodal-curvature-decomposition}
 \dfib^{\,2}\big|_\Gamma
 \;=\;
 \sum_{v\in V(\Gamma)}
 \kappaChHodge(\cA)\cdot\omega_{g(v)}
 \;+\;
 \sum_{e\in E(\Gamma)}
 \kappaChHodge(\cA)\cdot\delta_e,
\end{equation}
where $\omega_{g(v)}$ is the Arakelov $(1,1)$-form on the smooth
component $\Sigma_v$ (normalised so that
$\int_{\Sigma_v}\omega_{g(v)}=1$; see the explicit
formula in the proof below), and $\delta_e$ is the
distributional $(1,1)$-current supported on the node
corresponding to edge~$e$. This current arises from
the $\bar\partial$-defect of the Arakelov propagator
$S^{(g)}(z,w):=\partial_z G_{\mathrm{Ar}}^{(g)}(z,w)$
(where $G_{\mathrm{Ar}}^{(g)}$ is the Arakelov Green
function on $\Sigma_g$) in the thin-annulus limit of
the sewing collar. On a smooth genus-$g$ surface,
$\bar\partial_z S^{(g)}(z,w) = \pi\,\delta^{(2)}(z-w)
- \pi\,\omega_g(z)$, so the propagator is holomorphic
up to the harmonic correction $\omega_g$. In the
thin-annulus limit $t\to 0$ of the plumbing collar
$\{|z\,w|<|t|\}$ at edge~$e$, this defect localises:
\begin{equation}\label{eq:node-dbar-defect}
 \bar\partial_z S^{(g)}(z,w)\big|_{e}
 \;=\; \pi\,\delta_e(z),
\end{equation}
where $\delta_e$ is the $(1,1)$-current of unit mass
at the node (the limit of $\omega_g$ restricted to the
degenerating collar).
In particular:
\begin{enumerate}[label=\textup{(\alph*)}]
\item At genus~$0$
 \textup{(}$\Gamma$ a tree, $g(v)=0$ for all~$v$,
 $b_1(\Gamma)=0$\textup{)}: the curvature vanishes,
 $\dfib^{\,2}|_\Gamma = 0$, and the bar complex is flat.
 This is why the pair-of-pants multiplication
 \textup{(}Construction~\textup{\ref{constr:pants-bar-model}}\textup{)}
 is strictly associative at genus~$0$.

\item At genus~$1$
 \textup{(}$\Gamma$ has $b_1(\Gamma)=1$ or a single vertex
 with $g(v)=1$\textup{)}: the curvature is
 $\kappaChHodge(\cA)\cdot\omega_1$, matching the one-loop anomaly
 of the annular bar complex
 \textup{(}equation~\eqref{eq:circle-bar-curvature}\textup{)}.

\item At genus $g\ge 2$: the curvature
 $\kappaChHodge(\cA)\cdot\omega_g$ decomposes over the stable graph
 into vertex and edge contributions as
 in~\eqref{eq:nodal-curvature-decomposition}, and the genus
 tower is computed stratum by stratum: the genus-$g$ shadow
 obstruction tower is the direct sum of vertex contributions
 $\kappaChHodge(\cA)\cdot\omega_{g(v)}$ and edge contributions
 $\kappaChHodge(\cA)\cdot\delta_e$ over all strata of
 $\overline{\cM}_{g,n}$.
\end{enumerate}
\end{proposition}

\begin{proof}
Statement~(a) is the flatness of the genus-$0$ bar differential:
when $g = 0$ and $\Gamma$ is a tree, the propagator on
$\mathbb{P}^1$ is the globally meromorphic function
$K(z,w) = 1/(z-w)$, which satisfies $\bar\partial_z K = 0$
away from the diagonal. Since $\mathbb{P}^1$ has no period
matrix and no harmonic correction term, the
$\bar\partial$-defect that produces curvature at higher genus
is absent, and $(d^{(0)})^2 = 0$.
(At genus~$1$, the canonical Arakelov form is
$\omega_1 = \frac{i}{2\operatorname{Im}(\tau)}\,
dz\wedge d\bar{z}$ with $\int_{\Sigma_1}\omega_1 = 1$;
at genus $g\ge 2$, the Arakelov form
$\omega_g = \frac{i}{2g}\sum_{j,k=1}^g
(\operatorname{Im}\Omega)^{-1}_{jk}\,
\omega_j\wedge\bar\omega_k$
where $\{\omega_j\}$ is a basis of holomorphic
differentials and $\Omega$ is the period matrix,
normalised so that $\int_{\Sigma_g}\omega_g = 1$.
Neither exists at genus~$0$.)

Statement~(b) follows from
equation~\eqref{eq:circle-bar-curvature}: the genus-$1$ curvature
$\kappaChHodge(\cA)\cdot\omega_1$ was established as the one-loop anomaly
of the circle bar chain model. In the stable graph description,
the genus-$1$ stratum is either a smooth torus ($\Gamma$ a single
vertex with $g(v) = 1$) or a rational curve with a single node
($\Gamma$ a single vertex with $g(v) = 0$ and one self-edge,
$b_1(\Gamma) = 1$). In the latter case the curvature comes
entirely from the edge term $\kappaChHodge(\cA)\cdot\delta_e$. In the
former, it is $\kappaChHodge(\cA)\cdot\omega_1$ on the smooth torus.
Both contribute $\kappaChHodge(\cA)\cdot\omega_1$ to the total
curvature, as computed in Volume~I.

Statement~(c) is proved by induction on the number of edges
$|E(\Gamma)|$.

\emph{Base case}: $|E(\Gamma)|=0$, so $\Gamma$ has a single
vertex with $g(v) = g$ and no edges. Then
$B_\Gamma(\cA) = B^{(g)}_{n}(\cA)$ and the curvature is
$\dfib^{\,2} = \kappaChHodge(\cA)\cdot\omega_g$ from the vertex bar
complex alone on that same scalar lane.

\emph{Inductive step}: suppose the decomposition holds for all
stable graphs with $<|E(\Gamma)|$ edges. Adding the edge~$e$
that connects (or self-connects) vertex components introduces
the sewing via $B^{\mathrm{ann}}_e(\cA)$. By
equation~\eqref{eq:node-dbar-defect}, the
$\bar\partial$-defect of the propagator across the sewing
collar localises to the distributional current
$\kappaChHodge(\cA)\cdot\delta_e$ at the node. The sewing map
$\sigma_e$ of~\eqref{eq:sewing-map} is compatible with the
curvature decomposition: the differential
$d_\Gamma = \sum_v d_v + \sum_e d^{\mathrm{ann}}_e$
of~\eqref{eq:nodal-bar-differential} satisfies
$d_\Gamma^2 = \sum_v d_v^2 + \sum_e (d^{\mathrm{ann}}_e)^2$
because the cross-terms $d_v\circ d^{\mathrm{ann}}_e
+ d^{\mathrm{ann}}_e\circ d_v$ vanish: $d_v$ acts on the
vertex component $B^{(g(v))}_{n(v)}(\cA)$ and
$d^{\mathrm{ann}}_e$ acts on the sewing collar
$B^{\mathrm{ann}}_e(\cA)$, and these are supported on
disjoint regions of the nodal curve (the smooth component
interior and the plumbing collar respectively).

The global consistency is the distributional identity for the
Arakelov form on a nodal curve:
\[
 \omega_g
 \;=\;
 \sum_{v\in V(\Gamma)} \iota_{v*}\,\omega_{g(v)}
 \;+\;
 \sum_{e\in E(\Gamma)} \delta_e,
\]
where $\iota_{v*}$ denotes pushforward from the normalisation.
This identity holds because the Arakelov Green function on the
smoothing family $\Sigma_t$ ($t$ the plumbing parameter)
degenerates as $t\to 0$ into the sum of Green functions on
each irreducible component plus the logarithmic contribution
$\log|t|$ at each node, and $\bar\partial\partial\log|t|$
produces $\delta_e$
(cf.\ Wentworth~\cite{Wentworth91},
Wolpert~\cite{Wolpert90}).
Multiplying by $\kappaChHodge(\cA)$ and combining with the inductive
hypothesis on each vertex component
gives~\eqref{eq:nodal-curvature-decomposition}.
\end{proof}

\begin{remark}[Stable graph stratification and the genus tower]%
\label{rem:stable-graph-genus-tower}%
\index{stable graph!genus tower computation|textbf}%
\index{genus tower!stable graph decomposition}%
The stable graph stratification of $\overline{\cM}_{g,n}$ indexes
all degeneration types of a genus-$g$ curve with $n$ marked
points. The nodal bar complex
$B_\Gamma(\cA)$ of
Construction~\ref{constr:nodal-bar-model} provides a bar chain
model on each stratum, and the genus tower of Volume~I is computed
stratum by stratum as follows.

On the smooth locus $\cM_{g,n}\subset\overline{\cM}_{g,n}$,
the bar differential is curved with
$\dfib^{\,2} = \kappaChHodge(\cA)\cdot\omega_g$, and the curvature
class lies in the cohomology of $\cM_{g,n}$ with coefficients
in the chiral Hochschild complex. On each boundary stratum
$\partial_\Gamma\overline{\cM}_{g,n}$, the nodal bar complex
$B_\Gamma(\cA)$ replaces the smooth bar complex: the curvature
decomposes into vertex and edge contributions as
in~\eqref{eq:nodal-curvature-decomposition}. The total curvature
on $\overline{\cM}_{g,n}$ is then
\[
 \dfib^{\,2}\big|_{\overline{\cM}_{g,n}}
 \;=\;
 \kappaChHodge(\cA)\cdot\omega_g
 \;=\;
 \kappaChHodge(\cA)\cdot
 \biggl(
 \sum_{\Gamma\in\StGraph(g,n)}
 \frac{1}{|\Aut(\Gamma)|}
 \,\iota_{\Gamma*}\,
 \omega_\Gamma
 \biggr),
\]
where $\omega_\Gamma = \sum_v \omega_{g(v)} + \sum_e \delta_e$
and $\iota_{\Gamma*}$ is the pushforward from the normalisation
of the boundary stratum to $\overline{\cM}_{g,n}$.

The key structural point is that
\emph{nodal degenerations are where the curvature lives}. The
smooth locus is where the propagator is holomorphic and the
bar differential is flat at leading order; the curvature
$\kappaChHodge(\cA)\cdot\omega_g$ is the failure of the
holomorphic bar complex to extend flatly across the boundary
strata. The Arakelov propagator, which is smooth on the
interior of $\cM_{g,n}$ but has a $\bar\partial$-defect at the
nodes, provides the connection whose curvature is precisely
$\kappaChHodge(\cA)\cdot\omega_g$. This is why the genus tower is
computable: each genus-$g$ obstruction class decomposes into
contributions from stable graphs, and each contribution is
determined by the curvature data on lower-genus components
(the vertex contributions) and the sewing data at nodes
(the edge contributions).
\end{remark}

\begin{remark}[$K^{\kappaChHodge}=8$ on the
Mukai-$K3$ branch and $\hbar^2 K^{\kappaChHodge}=-1$ via
Bruinier--Heegner reciprocity; \ClaimStatusProvedElsewhere]
\label{rem:hoch-Kkappa-Mukai-Bruinier}
\index{Mukai!K3!K-invariant}%
\index{Bruinier--Heegner reciprocity!K-invariant}%
The genus tower
$\dfib^{\,2}|_{\overline{\cM}_{g,n}}=\kappaChHodge(\cA)\cdot\omega_g$ of
Remark~\ref{rem:stable-graph-genus-tower} admits a cross-volume
sharpening on the $K3$ input via the Vol~III $\Phi$-bridge: the
\emph{chiral-side $K$-invariant} of the Koszul-conductor pairing
\[
K^{\kappaChHodge}
\;:=\;
\tr_{Z^{\mathrm{der}}_{\mathrm{ch}}(\cA)}
\bigl(K_\cA\cdot\kappaChHodge(\cA)\bigr),
\]
evaluated on the B-family branch of the $K3$ landscape
($\cA=\cA_{K3}^{\mathrm{B}}$, the Bridgeland-stable presentation via
$D^b\mathrm{Coh}(K3)$), satisfies
\[
K^{\kappaChHodge}(\cA_{K3}^{\mathrm{B}})
\;=\;
2\,c_+\!\bigl(\mathrm{Mukai}(K3)\bigr)
\;=\;8,
\]
where $\mathrm{Mukai}(K3)=H^{2\ast}(K3,\Z)$ is the Mukai lattice of
signature $(4,20)$ and $c_+$ is the positive-definite rank
($c_+(\mathrm{Mukai}(K3))=4$). The factor $2$ is the CY-$2$
degree shift of Remark~\ref{rem:hoch-CY2-K3-shift}. At the
universal $\hbar^2$-level,
\[
\hbar^2\,K^{\kappaChHodge}(\cA_{K3}^{\mathrm{B}})
\;=\;-1,
\]
computed as a Heegner-divisor pushforward via Bruinier's
reciprocity~\cite{Bru02} for the singular-theta lift: the
$\hbar^2$-coefficient of $K^{\kappaChHodge}$ is the
constant term of the Borcherds vector-valued modular form
evaluated on the Heegner divisor class $H_{-1}\subset\mathrm{II}_{2,2}$,
which reads $-1$ by the Eisenstein-expansion reciprocity
(Bruinier~\cite{Bru02}, Prop.~5.1; compare
Remark~\ref{rem:bar-cobar-5-phi-trace-conductor} of the
bar-cobar review for the upstream trace identity). The sign is the
signature of the Mukai pairing under the parity involution
$\Z/2$ of Remark~\ref{rem:hoch-igusa-gauge-class}. The
chain-level content is the localisation of the genus-$g$
$\kappaChHodge\cdot\omega_g$ current against the
Heegner-divisor class on $\overline{\cM}_{g,n}$: the nodal-stratum
decomposition of
Proposition~\ref{prop:nodal-curvature-source} specialises on the
B-family to a Heegner-divisor concentration, whose Bruinier
reciprocity reads off the $\hbar^2$-coefficient.
\end{remark}

\begin{proposition}[Compatibility with the modular
clutching; \ClaimStatusProvedHere]%
\label{prop:nodal-bar-modular-operad}%
\label{rem:nodal-bar-modular-operad}% backward-compatible label
\index{modular operad!nodal bar complex compatibility|textbf}%
\index{clutching map!nodal bar sewing|textbf}%
Let $\xi_\Gamma$ denote the conjectural clutching maps of the modular
operad $\mathcal{O}^{\Ainf\text{-ch}}$
\textup{(}Definition~\textup{\ref{def:modular-operad-ainf-chiral}}\textup{)},
and let $\sigma_e$ denote the sewing maps of
Construction~\textup{\ref{constr:nodal-bar-model}}. Then:
\begin{enumerate}[label=\textup{(\alph*)}]
\item \emph{Operadic identification.}
 For each edge $e\in E(\Gamma)$, the sewing map $\sigma_e$
 of~\eqref{eq:sewing-map} realises the expected clutching map
 $\xi_e$ at the bar-complex level: the bar differential on
 $B_\Gamma(\cA)$ restricts to the vertex and edge
 differentials on each factor,
 as established in~\eqref{eq:nodal-bar-differential}.

\item \emph{Local order-independence.}
 If two edges $e_1,e_2\in E(\Gamma)$ have disjoint sewing
 collars, then the sewing maps $\sigma_{e_1}$ and
 $\sigma_{e_2}$ commute. At the level of differentials, this gives
 $d_\Gamma^2 = \kappaChHodge(\cA)\cdot\omega_\Gamma$ independently of the
 local sewing order for those disjoint collars. This is the
 bar-complex compatibility needed below; it is not a proof of the
 full modular operad associativity for
 $\mathcal{O}^{\Ainf\text{-ch}}$ in general.

\item \emph{$R$-matrix monodromy at sewing.}
 At each edge $e$, the wrap-around differential
 $d_{\mathrm{wrap}}$ of $B^{\mathrm{ann}}_e(\cA)$ inserts
 the $R$-matrix monodromy $\mathrm{Mon}(R)$
 \textup{(}Construction~\textup{\ref{constr:genus1-ainf-chiral-operations}}\textup{)},
 which is the holonomy of the KZ connection around the
 sewing circle. This is the datum that entangles the
 holomorphic and transverse factors at each node and
 prevents the nodal bar complex from splitting as a
 na\"ive tensor product.
\end{enumerate}
\end{proposition}

\begin{proof}
Part~(a): the sewing map $\sigma_e$ is defined by identifying
boundary circles of vertex bar complexes with boundary circles
of the annular bar complex. This identification is the data of
the clutching map $\xi_e$ in the modular operad, which glues
two $(g_i,n_i)$-labelled components along marked points to produce
a $(g,n)$-labelled component. The compatibility of differentials
is the content of~\eqref{eq:nodal-bar-differential}: the total
differential $d_\Gamma = \sum_v d_v + \sum_e d^{\mathrm{ann}}_e$
restricts to the vertex differential on each vertex factor and
to the annular differential on each edge factor.

Part~(b): if edges $e_1,e_2\in E(\Gamma)$ have disjoint sewing
collars and are sewn in two different orders, the resulting
differentials agree because
$d^{\mathrm{ann}}_{e_1}$ and $d^{\mathrm{ann}}_{e_2}$ act on
disjoint sewing collars and therefore commute as operators
on $B_\Gamma(\cA)$. The curvature contributions
$\kappaChHodge(\cA)\cdot\delta_{e_1}$ and
$\kappaChHodge(\cA)\cdot\delta_{e_2}$ are distributional currents
with disjoint support, so their sum is order-independent.

Part~(c) follows from the construction of
$d_{\mathrm{wrap}}$ in the annular bar complex: the monodromy
$\mathrm{Mon}(R)$ is the algebraic datum encoding the failure
of the bar complex to be periodic under the angular coordinate
of the sewing circle, and it enters multiplicatively in the
wrap-around face map.
\end{proof}

\subsection{The Drinfeld center of the chiral double}
\label{subsec:drinfeld-center-chiral-double}

\subsubsection*{Reframing: correctness of the boundary-bulk reconstruction}

Under the Drinfeld double programme
(Remark~\ref{rem:drinfeld-programme-four-conditions} of
Chapter~\ref{chap:ordered-associative}), the chiral double
$U_\cA = \cA \bowtie \cA^!$ is the boundary-line double whose
Drinfeld centre is conjectured in
Conjecture~\ref{conj:drinfeld-center-equals-bulk} to compare with the
bulk derived centre.
Part~(a) of the programme is
the existence of $U_\cA$; part~(b) is its uniqueness via the
antipode. The present section is part~(c), the correctness
clause: the comparison is correct precisely when the centre of the
double recovers the bulk, that is
$Z(U_\cA) \simeq Z^{\mathrm{der}}_{\mathrm{ch}}(\cA) =
\ChirHoch^\bullet(\cA)$. This is an independent compatibility between
the ordered Drinfeld-centre construction and the chiral Hochschild
coderivation construction.

The definition of the ``ordered Drinfeld centre'' that the
thesis requires is the $E_1$-Hochschild construction of
Francis~\cite{francis2013tangent}: $Z(U_\cA) := \mathrm{HH}^0$
of the $E_1$-module category of $U_\cA$, computed on the
ordered bar resolution. In the boundary-bulk picture this is
the line-side centre. The statement of
Conjecture~\ref{conj:drinfeld-center-equals-bulk} is the equivalence of
two a priori different constructions: the $E_1$-centre of the double and
the derived chiral centre of the boundary algebra.

The Drinfeld double programme of Chapter~\ref{chap:ordered-associative}
(Remark~\ref{rem:drinfeld-double-programme}) assembles the boundary
algebra $\cA$ and its Koszul dual $\cA^!$ into a single bar-complex-level
Hopf object $U_\cA = \cA \bowtie \cA^!$ whose module category is conjectured to be
the line-operator category $\cC_{\mathrm{line}}$. Part~(a) of the
programme is the construction of $U_\cA$; part~(b) is the identification
$U_\cA\text{-mod} \simeq \cC_{\mathrm{line}}$. Part~(c), addressed here,
is the Drinfeld-center identification.

In the classical setting, Drinfeld proved that the Drinfeld centre
$Z(D(H))$ of the double $D(H)$ of a Hopf algebra $H$ is the center of
the representation category: the monoidal center of $H\text{-mod}$ on
the categorical side, and the algebra of natural transformations of the
identity functor on the algebraic side. The analogous chiral statement
below conjecturally identifies the center of the chiral double with the
bulk algebra.

\begin{remark}[What ``ordered Drinfeld centre'' means]
\label{rem:ordered-drinfeld-center-definition}
The classical Drinfeld centre is an $E_2$-invariant of a monoidal
category. The chiral double $U_\cA$ is built from the \emph{ordered}
bar coalgebra $B^{\mathrm{ord}}(\cA) = T^c(s^{-1}\bar\cA)$, so the
natural centre on the bar-coalgebra side is an $E_1$-invariant, not
an $E_2$-invariant. In the absence of a fully worked $E_1$-centre
construction at the chiral level, we adopt the following tentative
identification: the ``ordered Drinfeld centre'' $Z(U_\cA)$ is
the derived endomorphism algebra
\[
Z(U_\cA) \;:=\;
\mathrm{HH}^0\bigl(U_\cA\text{-mod}\bigr)
\;=\;
\mathrm{R}\!\operatorname{End}_{U_\cA\otimes U_\cA^{\mathrm{op}}}(U_\cA)
\]
computed on the ordered bar resolution. The expected $E_2$
refinement, via the two-sided brace structure of the chiral Hochschild
complex, is not used in stating
Conjecture~\ref{conj:drinfeld-center-equals-bulk}; it will appear
only after the ordered-to-symmetric descent of
Chapter~\ref{chap:ordered-associative}.
\end{remark}

\begin{conjecture}[Drinfeld double programme, part~(c);
\ClaimStatusConjectured]
\label{conj:drinfeld-center-equals-bulk}
Let $\cA$ be a vertex algebra on a smooth curve $C$ satisfying the
hypotheses of the bulk--boundary--line factorization correspondence
\textup{(}Theorem~\textup{\ref{thm:bulk-boundary-line-factorization}}\textup{)},
and let $U_\cA = \cA \bowtie \cA^!$ be the chiral Drinfeld double of
Remark~\textup{\ref{rem:drinfeld-double-programme}}, assembled from the
ordered bar coalgebra $B^{\mathrm{ord}}(\cA) = T^c(s^{-1}\bar\cA)$
\textup{(}augmentation ideal\textup{)} and its Verdier
dual on $\Ran(C)$. Then the Drinfeld centre of $U_\cA$ is
quasi-isomorphic to the chiral derived center of $\cA$:
\begin{equation}
\label{eq:drinfeld-center-equals-bulk}
Z(U_\cA) \;\simeq\; Z^{\mathrm{der}}_{\mathrm{ch}}(\cA)
\;=\; \ChirHoch^\bullet(\cA).
\end{equation}
The bar used throughout is the ordered bar $B^{\mathrm{ord}}$: the equivalence is between the
ordered Drinfeld centre and the chiral Hochschild complex computed by
ordered configurations on $C$.
\end{conjecture}

\begin{remark}[Four functors, not three]
\label{rem:drinfeld-center-four-functors}
Equation~\eqref{eq:drinfeld-center-equals-bulk} relates two of the four
functors acting on the boundary algebra. The left-hand side is the
Drinfeld centre of the double $U_\cA = \cA \bowtie \cA^!$, an assembly
built from the Verdier functor $\cA \mapsto \cA^!$ and the bar coalgebra
$B^{\mathrm{ord}}(\cA)$. The right-hand side is the chiral Hochschild
complex $\ChirHoch^\bullet(\cA)$, the derived center obtained by the
Hochschild coderivation construction applied to $\cA$. Bar--cobar inversion
$\Omega(B(\cA)) \simeq \cA$ and Verdier duality do not suffice to identify
$Z(U_\cA)$ with $\ChirHoch^\bullet(\cA)$; the Hochschild construction is a
separate input. The conjecture asserts that these two inputs, applied
to the same boundary algebra, produce quasi-isomorphic complexes. It
does not assert that bar--cobar ``produces the bulk.''
\end{remark}

\subsubsection{Proof strategy: three steps}

A plausible proof proceeds in three steps, each with a specific
programmatic obstruction.

\begin{enumerate}[label=\textup{(\arabic*)}, leftmargin=2.4em]

\item \textbf{Pointwise reduction to the classical Drinfeld centre.}
For each point $x \in C$, the fibre $\cA_x$ is a vector space (the
stalk of $\cA$ at $x$), and the algebra $\cA_x \bowtie (\cA^!)_x$ at the
point is a plain associative algebra in $\Vect$. The classical Drinfeld
centre $Z(\cA_x \bowtie (\cA^!)_x)$ is computed by the usual formula: it
is the ordinary Hochschild complex $C^\bullet(\cA_x, \cA_x)$ in the
strictest case. The statement to verify in this step is that the
pointwise Drinfeld centre $Z(U_\cA)_x$ agrees with the pointwise
Hochschild complex $\ChirHoch^\bullet(\cA)_x$ under the identification
of stalks with their classical counterparts.

\item \textbf{Globalization via factorization.}
The sheaf $x \mapsto Z(U_\cA)_x$ on $C$ is not a priori a factorization
algebra; the pointwise centres must assemble, under collisions of points,
into a factorization algebra on $\Ran(C)$. The claim of step~(2) is that
this assembly exists and agrees, as a factorization algebra on $C$, with
the chiral Hochschild factorization algebra of
Theorem~\textup{\ref{thm:bulk_hochschild}}.

\item \textbf{Bar--cobar identification of the global assembly.}
The global factorization algebra of step~(2) is identified with
$\ChirHoch^\bullet(\cA)$ via the bar--cobar adjunction of Volume~I
(Theorem~A) together with the Verdier intertwining
$\mathbb{D}_{\Ran}(\barB^{\mathrm{ord}}(\cA))
\simeq \barB^{\mathrm{ord}}(\cA^!)$. Theorem~A supplies a
quasi-isomorphism between the bar and cobar sides, and the chiral
Hochschild complex sits in this picture as the derived center of the
cobar side.
\end{enumerate}

\begin{remark}[Programmatic obstructions, one per step;
\ClaimStatusConjectured]
\label{rem:drinfeld-center-obstructions}
Each of the three steps carries a concrete obstruction, tagged here as
programmatic rather than technical.

\emph{Obstruction to step (1): non-formal algebras.}
For class $\mathbf{G}$ and class $\mathbf{L}$ boundary algebras the
pointwise reduction is accurate: the fibre at a generic point captures
the associative data of $\cA_x \bowtie (\cA^!)_x$ up to the level of
strict centres. For class $\mathbf{M}$ boundary algebras (Virasoro,
$\cW_N$), the pointwise fibre loses information: the non-vanishing
higher $A_\infty$ operations $m_k$ (class~$\mathbf{M}$ is defined by
shadow depth $d_{\mathrm{alg}} = \infty$) contribute to the global
derived centre but are invisible in any single stalk. The
pointwise-Drinfeld-centre assembly is therefore only a first-order
approximation: the chiral derived centre contains additional
$A_\infty$-pieces that do not come from any stalk. The obstruction is
that class $\mathbf{M}$ breaks the ``classical stalk-wise reduction''
strategy at degree three and above.

\emph{Verdier factorization obstruction for \(\cA^!\).}
The Verdier dual $\cA^!$ is a factorization algebra only in the derived
sense; its strict factorization structure, required to compare
$Z(U_\cA)$ with a chiral Hochschild factorization algebra on equal
footing, is not automatic. Concretely, the factorization coproduct on
$B^{\mathrm{ord}}(\cA^!)$ is the Verdier transform of the deconcatenation
coproduct on $B^{\mathrm{ord}}(\cA)$, and this transform does not commute with stalk restriction
in general. Global descent asks that the family of pointwise
centres assemble into a single factorization algebra; this is a
non-trivial descent claim at the level of factorization homology.

\emph{Bar--cobar does not directly give the
center.}
Theorem~A of Volume~I produces a quasi-isomorphism
$\Omega(\barB(\cA)) \simeq \cA$ and a Verdier intertwining on
$\Ran(C)$. Neither statement says anything about the Hochschild complex
or the Drinfeld centre of the pair $(\cA, \cA^!)$. The bar--cobar
adjunction transports between the two sides of the Koszul pair but does
not identify the derived centre (an output of the Hochschild
coderivation construction) with any component of the bar--cobar
assembly. The missing datum is
a compatibility between the
Hochschild construction and the bar--cobar adjunction specific to the chiral
setting. No such compatibility is proved at present.
\end{remark}

\subsubsection{Cohomological constraint from Theorem~H}

Volume~I Theorem~H assigns to each family datum $H_H(\cA;S_\cA)$ a
model $K_{\cA,S_\cA}$ and an identification
\[
 \ChirHoch^n(\cA)\cong H^n(K_{\cA,S_\cA}),
 \qquad
 \operatorname{Supp}\ChirHoch^\bullet(\cA)\subseteq S_\cA.
\]

\begin{corollary}[Amplitude constraint on the chiral double;
conditional on Conjecture~\textup{\ref{conj:drinfeld-center-equals-bulk}}]
\label{conj:drinfeld-center-amplitude}
Assume Conjecture~\textup{\ref{conj:drinfeld-center-equals-bulk}} and
a family datum $H_H(\cA;S_\cA)$.  Then the Drinfeld centre of the
chiral double has support contained in~$S_\cA$.  When the groups
$H^n(K_{\cA,S_\cA})$ are finite-dimensional, it has the same finite
total dimension as
$\ChirHoch^\bullet(\cA)$:
\begin{equation}
\dim_{\mathrm{tot}}\, H^\bullet(Z(U_\cA))
\;=\; \dim_{\mathrm{tot}}\, \ChirHoch^\bullet(\cA)
\;<\; \infty.
\end{equation}
Its Hilbert polynomial is
$\sum_{n\in S_\cA}\dim H^n(K_{\cA,S_\cA})t^n$.
\end{corollary}

\begin{proof}
If $Z(U_\cA) \cong \ChirHoch^\bullet(\cA)$, the bound transfers
directly from Theorem~H of Volume~I whenever its hypotheses hold.
\end{proof}

\begin{remark}[Amplitude constraint]
\label{rem:drinfeld-center-amplitude-constraint}
The family datum and the centre comparison are independent inputs.
Together they transport the full support~$S_\cA$ and each graded
dimension to $Z(U_\cA)$.  This gives a degree-by-degree test for a
construction of the chiral double.
\end{remark}

\begin{computation}[Heisenberg comparison and grading obligation for
Conjecture~\textup{\ref{conj:drinfeld-center-equals-bulk}};
\ClaimStatusConditional]
\label{comp:drinfeld-center-heisenberg}
Let $\cA=\cH_k$ at nonzero level and identify its vertex presentation
with the rank-one even superboson $B_{\mathfrak h}$. Assume the
proposed first-formal-neighbourhood model for the chiral double is
\begin{equation}
Z(U_{\cH_k})_{\mathrm{fib},1} \;\simeq\;
\Bbbk\oplus \Bbbk\,\varepsilon_k,\qquad \varepsilon_k^2=0,
\end{equation}
with two displayed classes and, at this stage, an unassigned
cohomological grading. Bakalov--De Sole--Kac
\cite{BDSK21} \textup{(}Theorem~7.4\textup{)} compute the bounded profile
$(2,1,0,\ldots)$. Hence the simultaneous assumptions
that $\chi^{\mathrm{bd}}_{B_{\mathfrak h}}$ is a quasi-isomorphism and
that $Z(U_{\cH_k})\simeq\ChirHoch^\bullet(\cH_k)$ require three total
cohomology classes: two in degree zero and one in degree one. Thus the
completed comparison has the exact target
\[
 \dim H^0 Z(U_{\cH_k})=2,
 \qquad
 \dim H^1 Z(U_{\cH_k})=1,
 \qquad
 H^n Z(U_{\cH_k})=0\quad(n\ge2).
\]
Constructing the completed differential and grading proves these
identities from the displayed first-neighbourhood input.
\end{computation}

\begin{remark}[Fixed fibre and formal level family]
\label{rem:drinfeld-center-heisenberg-tamarkin-inconsistency}
The notation $\Bbbk[[\kappa]]$ denotes the completed base of the
one-parameter family of Heisenberg chiral algebras.  Theorem~H concerns
the chart complex specified by $H_H(\cH_k;S_{\cH})$.  Tamarkin's
parameter algebra concerns deformation of the family.  A comparison
between them consists of a cochain map together with its
quasi-isomorphism and completion hypotheses.
\end{remark}

\begin{remark}[Reframing under the unifying thesis: the Tamarkin
power series is the reconstructor deformation parameter]
\label{rem:drinfeld-center-heisenberg-thesis-resolution}
The power series $\Bbbk[[\kappa]]$ of
Computation~\ref{comp:tamarkin-e2-heisenberg} parametrises deformation
of the reconstructor $U_{\cH_k}$.  The fixed-level bulk is the chart
complex $C^\bullet_{\mathrm{ch}}(\cH_k,\cH_k)$, whose support is
computed from $H_H(\cH_k;S_{\cH})$.  The two objects are related by a
specified comparison map. The grading obligation in
Computation~\ref{comp:drinfeld-center-heisenberg} becomes the grading
test: the completed centre must have degree-zero dimension two and
degree-one dimension one in the BDSK comparison regime.
\end{remark}

\begin{remark}[Sanity check at $k=0$]
\label{rem:drinfeld-center-heisenberg-k-zero}
At $k=0$, the classical $r$-matrix $r(z)=k\Omega_{\mathrm{tr}}/z$
vanishes and the proposed double becomes the ordinary tensor product
$\cH_0\otimes\cH_0$.  The zero-level chart is a distinct algebra from
the nondegenerate superboson.  Its comparison therefore uses its own
family datum $H_H(\cH_0;S_0)$ and a direct centre calculation.
\end{remark}

\begin{remark}[Virasoro centre comparison]
\label{rem:drinfeld-center-virasoro-open}
Assume the bounded-to-chart map $\chi_c^{\mathrm{bd}}$ is a
quasi-isomorphism.  The Virasoro chart then has Hilbert polynomial
$1+t^2+t^3$ by Bakalov--De Sole--Kac
\cite{BDSK21} \textup{(}Theorem~7.2\textup{)}. Together with
Conjecture~\textup{\ref{conj:drinfeld-center-equals-bulk}}, this predicts
the same polynomial for the Drinfeld centre of
$\Vir_c\bowtie\Vir_{26-c}$.  A direct completed-centre calculation
tests this prediction.  Each principal $\cW_N$ family similarly asks
for its own datum $H_H(\cW_N;S_N)$ and centre comparison.
\end{remark}

\begin{remark}[Relation to the bulk--boundary--line datum]
\label{rem:drinfeld-center-bulk-boundary}
Conjecture~\textup{\ref{conj:drinfeld-center-equals-bulk}} fits into the
bulk--boundary--line datum of
Theorem~\textup{\ref{thm:bulk-boundary-line-factorization}} as follows.
Part~(iii) of the theorem identifies the bulk observable complex
$\mathsf{Obs}(\C \times \R)$ with the chiral Hochschild complex
$\ChirHoch^\bullet(\cA_\partial)$, and, in the boundary-linear sector,
with the derived centre $Z^{\mathrm{der}}(\cA_\partial)$. The
Drinfeld-centre identification~\eqref{eq:drinfeld-center-equals-bulk}
adds a second geometric presentation of the same object: the bulk is
computed on one hand by chiral Hochschild cochains in the coderivation
model and on the other hand by the
Drinfeld centre of the chiral double (the presentation via the Koszul
pair and its Hopf-algebraic assembly). The conjecture says these two
presentations agree.

This is the expected chiral analogue of Drinfeld's classical theorem
that the Drinfeld centre of the double is the bulk of the boundary
representation category. At the level of line operators, it says that
the centre of the chiral Hopf algebra controlling line-operator braiding
is the same bulk algebra that controls the $(-1)$-shifted PVA bracket of
local operators. A full proof would close the loop: the Hochschild
presentation of the bulk, the Koszul-dual presentation of the boundary,
and the Drinfeld-centre presentation of the double would all produce
the same factorization algebra on $C \times \R$.
\end{remark}

\begin{remark}[Hall--Drinfeld double versus Drinfeld
Yangian, the fourth taxonomic slot of quasi-Hopf chiral quantum
groups; \ClaimStatusConjectured]
\label{rem:hoch-hall-drinfeld-fourth-slot}
\index{Drinfeld Yangian!versus Hall--Drinfeld double}%
\index{Hall--Drinfeld double!chiral taxonomic slot}%
\index{quasi-Hopf quantum group!four slots}%
The chiral double $U_\cA=\cA\bowtie\cA^!$ of
Conjecture~\ref{conj:drinfeld-center-equals-bulk} is one of
\emph{four} distinguishable quasi-Hopf chiral quantum groups in the
Vol~III K3 taxonomy (Vol~III Part~\ref{part:frontier}). The
four slots are:
\begin{enumerate}[label=\textup{(\alph*)}]
\item \emph{Drinfeld Yangian} $Y_\hbar(\fg)$: classical quantum
 group on a point, $E_0$-chiral, quasi-triangular with rational
 $r$-matrix $r^Y(z)=\hbar\Omega/z$. No Hochschild degree~$\geq 2$
 content (the pointwise centre is the Casimir algebra).
\item \emph{Chiral affine} / \emph{spectral Yangian}: the chiral
 enhancement of~(a) on a curve $C$, $E_1$-chiral, with
 $r$-matrix $r^{\mathrm{KM}}(z)=k\Omega_{\mathrm{tr}}/z$ on a level prefix, and
 chiral Hochschild support computed from a family datum
 $H_H(\cA;S_\cA)$; the ordered and symmetric chart models are related
 by the averaging quasi-isomorphism carried by that datum.
\item \emph{Drinfeld chiral double} $U_\cA=\cA\bowtie\cA^!$: the
Koszul double of Conjecture~\ref{conj:drinfeld-center-equals-bulk},
whose Drinfeld centre is presented, under the BZFN/Morita comparison
hypotheses, by \(E_2\)-modules over the chiral derived centre
$Z^{\mathrm{der}}_{\mathrm{ch}}(\cA)$ of the boundary. This is a
categorical presentation, not an equality with the cochain complex
itself; it lives on a curve and carries a bar-coalgebra quasi-Hopf
structure on $B^{\mathrm{ord}}(\cA)$ (not on $\cA$ itself).
\item \emph{Hall--Drinfeld double}: the CY-$2$ specialisation of~(c)
 via the Vol~III $\Phi$-functor on $D^b\mathrm{Coh}(K3)$. The
 boundary input is $\cA_{K3}$, the Mukai datum supplies the
 $[2]$-shift of Remark~\ref{rem:hoch-CY2-K3-shift}, the bulk
is the K3 Hall algebra in the sense of
Kapranov--Vasserot/Cohomological Hall Algebra theory
(compare Vol~III Part~\ref{part:frontier}), and the
Hochschild comparison target is the Igusa line of
Remark~\ref{rem:hoch-igusa-gauge-class}.
\end{enumerate}
The chiral double (c) and the Hall--Drinfeld double (d) are
\emph{not} the same object: (c) works for any logarithmic
$\SCchtop$-algebra $\cA$ on a curve and carries an
$E_1$-dg-coalgebraic structure (ordered bar); (d) is the K3-input
specialisation and carries an $E_2$-dg-coalgebraic structure
(CY-$2$ bi-grading, Serre projective Schur cocycle). The
four-slot taxonomy thus extends the three-slot (a)/(b)/(c) classical
Drinfeld picture to a fourth CY-$2$ slot where the closed-colour
input is K3 and the automorphic comparator is $\Delta_5$; this is the
Vol~III $\Phi$-transport of the Hochschild bulk-reconstruction
question of the preceding remark. The fourth slot is the precise
chiral analogue of the classical Ringel--Hall algebra of a CY-$2$
category, and its Drinfeld-centre calculation is the chiral
counterpart of the Okounkov--Smirnov computation for the
Nakajima quiver variety CoHA~\cite{OkounkovSmirnov16}.
\end{remark}

\section{K3 and class-\texorpdfstring{$\mathcal S$}{S} specialisations}
\label{sec:hoch-supplement}

\begin{remark}[Categorical Hochschild amplitude at the K3 base]
\label{rem:hoch-1}
Vol~I Theorem~H is a curve-level statement about chiral Hochschild
cochains on the ordered bar complex. The HKR theorem computes a
separate categorical Hochschild object for the Calabi--Yau input. For
a K3 surface,
\[
\HH^n(D^b\mathrm{Coh}(K3))
\;\cong\;
\bigoplus_{p+q=n}H^q(K3,\Lambda^pT_{K3})
\]
has Hilbert polynomial
\[
 1+22t^2+t^4 .
\]
The degree-$2$ summand is
\[
H^2(\cO_{K3})\oplus H^1(T_{K3})\oplus H^0(\Lambda^2T_{K3})
\cong \bC\oplus\bC^{20}\oplus\bC,
\]
and the total even cohomology carries the Mukai lattice of signature
$(4,20)$. Thus the K3 input has categorical Hochschild amplitude
\(\{0,2,4\}\).

For \(K3\times E\), K\"unneth gives
\[
P_{\HH(K3\times E)}(t)
=
(1+22t^2+t^4)(1+2t+t^2)
=
1+2t+23t^2+44t^3+23t^4+2t^5+t^6 .
\]
The holomorphic volume class is one generator inside the
$44$-dimensional degree-three group. Any projection of this
categorical amplitude to the curve-level chiral Theorem~H chart passes through
the chiral-specialisation comparison map
\[
 q_{\Phi,C}\colon
 \HH^\bullet(D^b\mathrm{Coh}(X))
 \longrightarrow
 \ChirHoch^\bullet(\Phi^{\mathrm{ch}}_C(X)).
\]
Let $A_X=\Phi_C^{\mathrm{ch}}(X)$ and assume
$H_H(A_X;S_X)$. For every closed categorical cochain~$x$ of degree
$n\notin S_X$, compatibility with Theorem~H requires a cochain
$h_x$ in the chart satisfying
\[
 q_{\Phi,C}(x)=d h_x.
\]
Thus the exact comparison problem is to construct $q_{\Phi,C}$ and
these off-support primitives, together with the homotopies expressing
their functoriality. The automorphic pairings with
$(\Phi_{10}^{\mathrm{un}})^{-1}$ supply scalar evidence for the BKM
branch; the displayed cochain equations supply the chart-level
obligation.
\end{remark}

\begin{remark}[Class-$\mathcal S$ Hochschild trace at $c_{2d} = -214$]
\label{rem:hoch-2}
The class-$\mathcal S$ Hochschild trace of $\mathcal T[A_1,\Sigma_{0,24}]$
at the K3 base point gives central charge
$c_{2d} = -12 c_{4d} = -12\cdot(107/6) = -2\cdot 107 = -214$ via the
Shapere--Tachikawa $2008$ \S 3 normalisation
$c_{4d} = (2n_v + n_h)/12$ applied to the Chacaltana--Distler $2010$
eq.~(5.14) count $(n_v, n_h) = (63, 88)$, yielding $c_{4d} = 107/6$.
The prime factorisation is \( -214=-2\cdot107\). No Mukai-signature
formula is asserted for this central charge: the Mukai lattice controls
the K3 categorical Hochschild pairing of
Remark~\ref{rem:hoch-1}, while the class-\(\mathcal S\) value above is
the Shapere--Tachikawa trace invariant of the four-dimensional theory.
The curve-level chiral Hochschild window remains the Theorem~H window
only after the hypotheses of Remark~\ref{rem:thm-H-three-way-scope} are
met.
Compare Chapter~\ref{V3-ch:k3-chiral-bialgebra-platonic}.
\end{remark}
